% ===========================================================================
%  Appendix B --- Argumentation
%
%  INTEGRATION NOTE.  This file compiles stand-alone AND drops into the
%  manuscript.  To integrate it as a chapter appendix, delete exactly five
%  things -- the \documentclass and \usepackage block below, the
%  \title/\author/\date/\maketitle group, \begin{document}, \end{document},
%  and the GENERATED BIBLIOGRAPHY block at the foot of the file -- and change
%  \section/\subsection to the manuscript's appendix-level commands.
%
%  ON THE BIBLIOGRAPHY.  Every \cite key here is the GLOSSARY's key: b5 is Dinh
%  in this appendix exactly as it is in A.1-A.12, b4 is Lezzar, b7 is COMSOL.
%  The KEY is the invariant; the NUMBER printed in a stand-alone build is local,
%  because this appendix cites a subset of the glossary's works and so cannot
%  reproduce the glossary's numbering without listing works it never cites.
%  Delete the block on integration and every \cite resolves against the
%  manuscript's own bibliography, which is the only arrangement in which the
%  printed numbers cannot lie.  Four keys are Appendix B's own -- bwrm (Walton,
%  Reed & Macagno, this appendix's scheme authority), bpre (the companion
%  preprint), bbrody (Brody et al. 1996, cited in C2) and bmort (Mortensen
%  2007, cited in C1) -- and the manuscript's bibliography must carry them;
%  each is declared with its reason inside the block.
%
%  ON CROSS-REFERENCES.  "Appendix~A" in this file always means the derivations
%  appendix; the glossary is called "the glossary" (ledger D73).  Its section
%  pointers ("glossary A.7") go through \xref, so they print the glossary's own
%  number stand-alone and the manuscript's once the glossary is loaded (D74).
%
%  Nothing else needs touching, because:
%
%    * every macro here is \providecommand, not \newcommand, so the frozen
%      PREAMBLE_FINAL.tex definitions win wherever the two overlap
%      (eleven names overlap: \fgap, \dcrit, \kapp, \Kin, \Kdepth, \Kdth,
%      \Kdcfd, \acal, \ud, \kster and \xref; the bodies are identical, so this
%      is a collision of names only, not of meaning); and
%    * \um is the one exception, set with \renewcommand: siunitx, which the
%      manuscript loads for the glossary, reserves the name, and a
%      \providecommand would leave every micrometre here an error (ledger
%      D76; appendixA/integrate_check.py builds the three appendices
%      together and checks it); and
%    * the column type L is guarded, so a manuscript definition of L is kept
%      rather than silently overridden.  (PREAMBLE_FINAL defines Y, not L, so
%      today the guard is a no-op; it is there so that a later clash is a
%      no-op too rather than a silent redefinition.); and
%    * the generated VALUE block is doubly guarded -- each carried file is
%      wrapped in \@ifundefined on a sentinel that file alone defines, and every
%      definition inside is \providecommand -- so a manuscript that has loaded
%      some or all of tex/notation_values, tex/kernel_pairing_values,
%      tex/validation_values and tex/appendix_B_stats keeps its own copies and
%      this appendix adds only what is missing.  Delete nothing there on
%      integration.  The block is machine-carried from PREAMBLE_FINAL.tex; edit
%      it by regenerating (gen_B_preamble_values.py), never by hand.
%
%  The manuscript preamble must, however, add three packages this appendix
%  needs and PREAMBLE_FINAL.tex does not load:
%
%      \usepackage{tikz}\usetikzlibrary{arrows.meta,fit,positioning}
%      \usepackage{enumitem}
%      \usepackage{textcomp}
%
%  Note also that PREAMBLE_FINAL sets margin=1in against the 0.9in used here;
%  every figure in this file is drawn to fit the narrower 6.5in text block.
% ===========================================================================
\documentclass{article}
\usepackage{graphicx} % Required for inserting images
\usepackage[letterpaper,margin=0.9in]{geometry}
\usepackage{amsmath,amssymb}
\usepackage{textcomp}
\usepackage{array}
\usepackage{booktabs}
\usepackage{longtable}
\usepackage{tikz}
\usepackage{enumitem}
% url: the carried COMSOL entry (b7) prints \url{www.comsol.com}. PREAMBLE_FINAL
% loads url for the same entry, so this is not one of the packages the manuscript
% must add on integration -- it is here only so the stand-alone build has it too.
\usepackage{url}
\usetikzlibrary{arrows.meta,fit,positioning}
% --- notation -------------------------------------------------------------
% \providecommand throughout: see the INTEGRATION NOTE above.
\providecommand{\fgap}{f_{\mathrm{gap}}}
\providecommand{\dcrit}{d_{\mathrm{crit}}}
\providecommand{\kapp}{\kappa_{\mathrm{app}}}
\providecommand{\ksteric}{\kappa_{\mathrm{steric}}}
% Glossary macro NAMES, not just glyphs. PREAMBLE_FINAL defines \Kin as
% K_{\text{in-plane}} (text-mode hyphen) and \kinplane does not exist there, so a
% locally-named \kinplane would print K_{inplane} in this appendix and
% K_{in-plane} in the glossary for one quantity, with \providecommand unable to see
% the clash. Using the glossary's names lets the frozen definitions win.
\providecommand{\Kin}{K_{\text{in-plane}}}
\providecommand{\kinplane}{\Kin}   % back-compatible alias for existing call sites
\providecommand{\Kdepth}{K_{\mathrm{depth}}}
\providecommand{\Kdth}{K_{\mathrm{depth}}^{\mathrm{th}}}    % ANALYTIC depth factor
\providecommand{\Kdcfd}{K_{\mathrm{depth}}^{\mathrm{CFD}}}  % MEASURED depth factor
\providecommand{\acal}{\alpha_{\mathrm{cal}}}
% \ud and \kster are the glossary's spellings, provided here ONLY so that the
% correspondence stated in the aspect-ratio convention can print them. Do not
% use them elsewhere in this appendix; \dwc and \ksteric are its own.
\providecommand{\ud}{\widetilde{d}}
\providecommand{\kster}{\kappa_{\mathrm{steric}}}
\providecommand{\feff}{f_{\mathrm{eff}}}
% \um: siunitx, which the glossary preamble loads, reserves \um (an error outside
% \si), so a \providecommand alone would yield to it and every micrometre here
% would fail in the manuscript. Provide, then set; \si{\um} is unaffected (D76).
\providecommand{\um}{}\renewcommand{\um}{\ensuremath{\,\mu\mathrm{m}}}
\providecommand{\indep}{\mathrel{\perp\!\!\!\perp}}
\providecommand{\field}[1]{\smallskip\noindent\textbf{#1}\ }
\providecommand{\wc}{w_{c}}
\providecommand{\dskim}{\delta_{\mathrm{skim}}}
\providecommand{\prov}[1]{{\footnotesize\ttfamily[#1]}}
% \xref{label}{fallback}: the glossary's macro, same body as PREAMBLE_FINAL.tex. A
% pointer into the glossary prints its fallback ("A.7") stand-alone and the
% manuscript's number once the glossary's labels exist (ledger D74).
\makeatletter
\providecommand{\xref}[2]{\@ifundefined{r@#1}{#2}{\ref{#1}}}
\makeatother
\providecommand{\fnom}{f_{\mathrm{nom}}}
\providecommand{\dwc}{d/\wc}
\providecommand{\chired}{\chi^{2}_{\mathrm{red}}}
\providecommand{\wref}{w_{\mathrm{ref}}}
\providecommand{\fref}{f_{\mathrm{ref}}}
\providecommand{\wcref}{\wc(\mathrm{ref})}
\providecommand{\wci}{\wc(i)}
\providecommand{\fgapref}{\fgap(\mathrm{ref})}
\providecommand{\kplat}{\kappa_{\mathrm{plateau}}}
\providecommand{\Rc}{R_{c}}
\providecommand{\Rs}{R_{s}}
\providecommand{\Qc}{Q_{c}}
\providecommand{\Qs}{Q_{s}}

% --- generated values ------------------------------------------------------
% VALUES ARE GENERATED, NOT TYPED (PREAMBLE_FINAL.tex's rule).  Every number
% this appendix prints that the derivation record defines is cited through a
% macro from the block below, never re-typed.  The block is carried from
% PREAMBLE_FINAL.tex mechanically -- see gen_B_preamble_values.py -- so B and
% the manuscript cannot drift apart without the gate firing.
% >>> BEGIN GENERATED VALUE BLOCK -- gen_B_preamble_values.py
% Carried from PREAMBLE_FINAL.tex and tex/appendix_B_stats.tex. DO NOT EDIT:
% regenerate with gen_B_preamble_values.py; validate.sh gate 15h diffs this
% region against a fresh generation and fails on any hand edit.
\makeatletter
% tex/notation_values.tex -- skipped when the manuscript has loaded it
\@ifundefined{kappVal}{%
\IfFileExists{tex/notation_values.tex}{\input{tex/notation_values}}{%
  \providecommand{\kappVal}{4.9138}\providecommand{\kappFit}{4.9137}%
  \providecommand{\kappFitFull}{4.913717814706706}%
  \providecommand{\kappCIlo}{4.75}\providecommand{\kappCIhi}{5.08}\providecommand{\kappCV}{4.0}%
  \providecommand{\kstericVal}{2}\providecommand{\KinVal}{1.000}%
  \providecommand{\kappOverRootThree}{2.8370}%
  \providecommand{\plateauKappa}{4.91}\providecommand{\footKappa}{1.34}%
  \providecommand{\plateauN}{8}\providecommand{\footN}{3}%
  \providecommand{\kappaSpanLo}{1.20}\providecommand{\kappaSpanHi}{5.20}%
  \providecommand{\KdepthAnalyticShallow}{0.8879}\providecommand{\KdepthAnalyticMid}{0.9500}%
  \providecommand{\KdepthAnalyticDeep}{0.9794}%
  \providecommand{\KdepthCfdShallow}{0.891}\providecommand{\KdepthCfdMid}{0.945}%
  \providecommand{\KdepthCfdDeep}{0.981}%
  \providecommand{\KdepthCalNarrow}{0.888}\providecommand{\KdepthCalMid}{0.849}%
  \providecommand{\KdepthCalWide}{0.811}%
  \providecommand{\residualFolded}{2.4569}\providecommand{\residualSepShallow}{2.7574}%
  \providecommand{\residualSepMid}{2.5999}\providecommand{\residualSepDeep}{2.5044}%
  \providecommand{\stericFloorShallow}{1.782}%
  \providecommand{\FclosedVone}{0.5592}\providecommand{\FclosedVthree}{0.7189}%
  \providecommand{\FcomsolVone}{0.5070}\providecommand{\FcomsolVthree}{0.6281}%
  \providecommand{\FexpVone}{0.505}\providecommand{\FexpVthree}{0.649}%
  \providecommand{\FsdVone}{0.033}\providecommand{\FsdVthree}{0.036}%
  \providecommand{\alphaThreeDVone}{0.811}\providecommand{\alphaThreeDVthree}{0.660}%
  \providecommand{\ReWorstChannel}{5.2}\providecommand{\ReGapScale}{5.2\times10^{-3}}%
  \providecommand{\fgapBandLo}{3.63\times10^{-5}}\providecommand{\fgapBandHi}{1.93\times10^{-3}}%
  \providecommand{\fgapDesignCap}{5.76\times10^{-4}}%
}
}{}%

% tex/kernel_pairing_values.tex -- skipped when the manuscript has loaded it
\@ifundefined{xiaAgreePct}{%
\IfFileExists{tex/kernel_pairing_values.tex}{\input{tex/kernel_pairing_values}}{%
  \providecommand{\ctwobar}{4.177}%
  \providecommand{\kFval}{5.798}%
  \providecommand{\depthEff}{0.8475}%
  \providecommand{\kFmeanK}{5.7862}%
  \providecommand{\kFmeanKerr}{0.21}%
  \providecommand{\kFoneAlpha}{5.5342}%
  \providecommand{\kFoneAlphaErr}{4.6}%
  \providecommand{\ctwobarDecl}{4.182}%
  \providecommand{\kFdecl}{5.802}%
  \providecommand{\transferExact}{19.9}%
  \providecommand{\transferDecl}{20.3}%
  \providecommand{\dcritGGleadIn}{2.170}%
  \providecommand{\dcritGGleadOut}{1.504}%
  \providecommand{\choiceLeadGap}{0.2}%
  \providecommand{\crossUp}{18.0}%
  \providecommand{\crossDown}{-15.3}%
  \providecommand{\alphaDIn}{1.867}%
  \providecommand{\dcritGGIn}{2.174}%
  \providecommand{\dcritFFIn}{2.482}%
  \providecommand{\dcritFgIn}{2.565}%
  \providecommand{\dcritgFIn}{2.103}%
  \providecommand{\choiceIn}{14.1}%
  \providecommand{\alphaDOut}{2.692}%
  \providecommand{\dcritGGOut}{1.507}%
  \providecommand{\dcritFFOut}{1.743}%
  \providecommand{\dcritFgOut}{1.779}%
  \providecommand{\dcritgFOut}{1.477}%
  \providecommand{\choiceOut}{15.6}%
  \providecommand{\KdepthDevShallow}{0.349}%
  \providecommand{\KdepthDevMid}{0.530}%
  \providecommand{\KdepthDevDeep}{0.160}%
  \providecommand{\KdepthDevMidTh}{0.527}%
  \providecommand{\KdepthDevMax}{0.531}%
  \providecommand{\kappSuppliedGap}{1.7\times10^{-5}}%
  \providecommand{\dinhBandLo}{1.25}%
  \providecommand{\dinhBandHi}{2.17}%
  \providecommand{\dinhKappaLocal}{4.08}%
  \providecommand{\dinhCollapseLo}{1.80}%
  \providecommand{\dinhKneeRatio}{1.21}%
  \providecommand{\dinhBandWidth}{1.74}%
  \providecommand{\kappCIhalfPct}{3.4}%
  \providecommand{\xiaKappaA}{5.03}%
  \providecommand{\xiaKappaB}{5.16}%
  \providecommand{\xiaKappaC}{5.32}%
  \providecommand{\xiaKappaMean}{5.17}%
  \providecommand{\xiaAgreePct}{5.2}%
  \providecommand{\depthCorrPct}{5.3}%
  \providecommand{\chiGlobal}{30.6}%
  \providecommand{\chiPlateau}{3.8}%
  \providecommand{\chiRatio}{8.0}%
  \providecommand{\etaShortfallPct}{9.7}%
  \providecommand{\etaVal}{0.903}%
  \providecommand{\runoutFmatched}{0.519}%
  \providecommand{\runoutFmismatched}{0.265}%
  \providecommand{\residPctVOne}{9.3}%
  \providecommand{\residPctVThree}{12.6}%
  \providecommand{\squareDuctR}{28.4542}%
  \providecommand{\skimErrDesign}{0.47}%
  \providecommand{\skimErrBand}{0.86}%
  \providecommand{\DhOutlet}{75.8}%
  \providecommand{\DhInlet}{97.7}%
  \providecommand{\Uoutlet}{0.069}%
  \providecommand{\ReHighFlow}{16}%
  \providecommand{\RePRP}{2.6}%
  \providecommand{\LeInlet}{60.5}%
  \providecommand{\LeOverLamInlet}{1.23}%
  \providecommand{\LeOverLamGInlet}{1.31}%
  \providecommand{\LeOutlet}{46.9}%
  \providecommand{\LeOverLamOutlet}{0.96}%
  \providecommand{\LeOverLamGOutlet}{1.02}%
  \providecommand{\aTwelveFL}{1.42}%
  \providecommand{\DTwelveFL}{2.84}%
  \providecommand{\QgapDinh}{3.12}%
  \providecommand{\dldExcessLo}{1.23}%
  \providecommand{\dldExcessHi}{1.33}%
  \providecommand{\ihRatioA}{1.44}%
  \providecommand{\ihRatioB}{1.48}%
  \providecommand{\ihRatioC}{1.21}%
  \providecommand{\ihRatioLo}{1.21}%
  \providecommand{\ihRatioHi}{1.48}%
  \providecommand{\ihRatioMean}{1.4}%
  \providecommand{\davisKappaHi}{2.659}%
  \providecommand{\davisKappaLo}{2.459}%
  \providecommand{\inglisKappaLo}{2.040}%
  \providecommand{\inglisKappaHi}{2.449}%
  \providecommand{\dldKappaLo}{2.040}%
  \providecommand{\dldKappaHi}{2.659}%
  \providecommand{\dldConvergePct}{0.4}%
  \providecommand{\expBeta}{0.51}%
  \providecommand{\expGamma}{1.02}%
  \providecommand{\plateauUwc}{0.035}%
  \providecommand{\footUwc}{0.012}%
  \providecommand{\betaCIlo}{0.41}%
  \providecommand{\betaCIhi}{0.61}%
  \providecommand{\gammaCIlo}{0.84}%
  \providecommand{\gammaCIhi}{1.20}%
  \providecommand{\hillKmax}{5.568}%
  \providecommand{\hillUhalf}{0.01651}%
  \providecommand{\hillN}{1.791}%
  \providecommand{\dinhSaddleFgap}{1.041\times10^{-4}}%
  \providecommand{\dinhSaddleU}{0.02894}%
  \providecommand{\sibALo}{1.02}%
  \providecommand{\sibAHi}{1.47}%
  \providecommand{\sibBLo}{1.26}%
  \providecommand{\sibBHi}{1.81}%
  \providecommand{\dinhBracketLo}{1.0}%
  \providecommand{\dinhBracketHi}{4.7}%
}
}{}%

% tex/validation_values.tex -- skipped when the manuscript has loaded it
\@ifundefined{EsdClosedVOne}{%
\IfFileExists{tex/validation_values.tex}{\input{tex/validation_values}}{%
  \providecommand{\EsdClosedVOne}{+1.64}%
  \providecommand{\EsdComsolVOne}{+0.06}%
  \providecommand{\EsemComsolVOne}{+0.12}%
  \providecommand{\uvalExcessVOne}{0.048}%
  \providecommand{\EsdClosedVThree}{+1.94}%
  \providecommand{\EsdComsolVThree}{-0.58}%
  \providecommand{\EsemComsolVThree}{-1.16}%
  \providecommand{\uvalExcessVThree}{0.011}%
  \providecommand{\uDoverUnumVOne}{660}%
  \providecommand{\uDoverUnumVThree}{720}%
  \providecommand{\uDoverUnum}{660}%
  \providecommand{\uvalExcessMax}{0.05}%
  \providecommand{\FexcessPctVOne}{10.7}%
  \providecommand{\FexcessPctVThree}{10.8}%
  \providecommand{\dFrecurConst}{2.35\times10^{-4}}%
  \providecommand{\dFrecurVar}{1.97\times10^{-5}}%
  \providecommand{\peMarginMin}{48}%
  \providecommand{\peHalfAgreeVOne}{2.3}%
  \providecommand{\peHalfAgreeVThree}{4.7}%
  \providecommand{\netEmergentPct}{0.12}%
  \providecommand{\siblingSqrtcPct}{5}%
}
}{}%

% tex/appendix_B_stats.tex -- Appendix B's own graph counts
\@ifundefined{cifOpenCQ}{%
\IfFileExists{tex/appendix_B_stats.tex}{\input{tex/appendix_B_stats}}{%
\providecommand{\cifNodes}{68}
\providecommand{\cifOpenCQ}{10}
\providecommand{\cifNodesWithOpenCQ}{7}
\providecommand{\cifProven}{22}
\providecommand{\cifWarranted}{33}
\providecommand{\cifDefeated}{4}
\providecommand{\cifSchemes}{13}
\providecommand{\cifCrosswalk}{49}
\providecommand{\cifTierAnalogical}{1}
\providecommand{\cifTierDeductive}{19}
\providecommand{\cifTierExperimental}{18}
\providecommand{\cifTierNumerical}{12}
\providecommand{\cifTierTestimonial}{9}
\providecommand{\cifTierWeakCorroboration}{9}
}
}{}%

\makeatother
% <<< END GENERATED VALUE BLOCK
\setlength{\emergencystretch}{3em}  % absorb minor overfulls in narrow columns / long math spans
\makeatletter
\@ifundefined{NC@rewrite@L}%
  {\newcolumntype{L}[1]{>{\raggedright\arraybackslash}p{#1}}}{}
\makeatother
\setlist[description]{leftmargin=1.6em, style=nextline, itemsep=3pt, topsep=4pt}
\title{Appendix~B --- Argumentation}
\author{Austin Routt}
\date{June 2026}
\begin{document}

\maketitle

% NOTE: in the integrated manuscript this becomes Appendix B; as a stand-alone
% file it is typeset with ordinary \section headings.

\section{Introduction}

This appendix records the argument of the main text as an explicit,
\emph{defeasible} argument map. The claim being audited is the thesis itself, as
the argument graph states it: that CIF can be engineered from its
\emph{performance targets} instead of from a tuned per-gap schedule; that its
yield follows a closed-form law with no fitted constant and its selectivity a
derived law with a single calibrated prefactor; that the two are \emph{provably
decoupled}, so the size threshold and the filtrate fraction target are
independent design \emph{inputs} rather than coupled outcomes; and that all of
this holds inside a validity envelope the design tool checks: it refuses a
design that leaves the plateau or the Reynolds gates, and it reports, without
refusing, a design over its length budget. The
plain-English rendering compresses those four to three (engineerable from
targets, on decoupled yield and cutoff laws, inside a checked envelope), and
this appendix argues them in that order.
Both performance axes are therefore in scope, not the cutoff alone. The map has
two storeys accordingly: the lower one is the cutoff law (C1--C3$'$, with C5--C7
corroborating into it and the yield-side separation lemma C4 drawn in but
deliberately \emph{disconnected} from its spine), and the upper one is what that
disconnection buys: the design map the separation, with the plateau schedule, makes locally
invertible (C8), and the
five-leg accrual over the operating law, the separation lemma, the map's
inversion, its realization and the declared scope, that is the thesis. C4 is the
hinge between them: drawn low, joined to nothing in the cutoff chain, and a
premise of everything above it. C9, the transport lemma, belongs to neither
storey: it upgrades C1's steric premise and, within the model's premises, closes
C3d's residual list, without lying on any path to the thesis.

This appendix's purpose is auditing rather than exposition:
each load-bearing claim is stated as a separate node, together with the support
it rests on, the named reasoning scheme that licenses it, and, most
importantly, the residual risk, i.e.\ the critical questions that remain
unanswered. The closed-form derivations the chain invokes are proved line by
line in Appendix~A; the symbolic results are reproduced by the reference
code \texttt{cif\_cutoff.py} (its \texttt{verify\_all()} self-tests)
and the empirical and CFD results by that program's frozen Part~II regression
stamps, both recorded in the derivation record \cite{b12a}. This appendix is the connective tissue between them: it does not
re-prove the mathematics, it certifies that every link in the argument has been
examined and states exactly what is, and is not, yet settled. In brief: the
cutoff law is \emph{proven in form} and \emph{empirically solid in its apparent
prefactor}; the separation $F \indep \fgap$ is proven outright and $F$'s absolute
magnitude independently validated against 3-D simulation; and the target-to-%
geometry map inverts, inside a declared and machine-checked envelope. What is
\emph{not} settled is stated node by node rather than summarised away. The
argument graph carries \cifOpenCQ\ undischarged critical questions distributed
over \cifNodesWithOpenCQ\ of its \cifNodes\ argument nodes. Two of those deserve
naming here. The lumped residual $\sqrt{c}\,C$, tied to the unpublished gap count
$N$, is the one the \emph{interpretation} of the cutoff prefactor turns on
(C3d): it is one item on that list rather than the whole of it. And exactly one
of the ten is a \emph{live rival} rather than a residual: the $\feff(d)$ account
of the cutoff knee, which C3c argues against and does not eliminate, and which
the graph marks with its own status for that reason. Ten undischarged questions
and one live rival are two different counts of two different things; both are
carried below. The full justification of that assessment is the business of the
table below.

\subsection{How to read the table}
\label{sec:b-howtoread}

Each node is a \emph{defeasible} argument in the sense of Walton, Reed \& Macagno
\cite{bwrm} (hereafter WRM~2008): a claim follows from its premises \emph{unless}
one of the governing scheme's critical questions (CQs) is answered against it.
The appendix instantiates \cifSchemes\ distinct schemes (twelve of WRM's, plus
\emph{accrual} as a local tag the thesis section declares as such), of which
defeasible \emph{modus ponens} is the most common and Argument from Expert
Opinion governs C1. Each node names its own scheme. A node is
therefore never simply ``true'' or ``false''; it carries a calibrated status
and an explicit list of what could still defeat it. The statuses below are the
ones the argument graph defines, and the table uses no others; each says for
itself whether it leaves a critical question open.

\begin{itemize}
  \item \textbf{proven.} A deductive identity (symbolic or circuit-theoretic);
        no open critical question.
  \item \textbf{warranted.} Strong defeasible support, with every critical
        question material to the claim answered and residual risk minor or none.
        A qualifier (``on CIF data,'' ``on the form,''
        ``uniform-$\feff$-robust,'' ``conditionally'') names the scope within
        which the warrant holds; it narrows a warrant without being a status of
        its own.
  \item \textbf{resolved.} A definitional or adjudicative question, settled.
  \item \textbf{validated.} An empirical \emph{magnitude} confirmed against
        independent data. This is a positive finding, not a weakening: it says
        the number was checked, not that the claim is in doubt. C4 carries it
        beside \emph{proven}, because the separation and the magnitude are two
        different claims.
  \item \textbf{warranted as a band.} Warranted only as a band or consistency
        check, not as a discriminating test (C7).
  \item \textbf{warranted (existence).} The \emph{existence} claim is
        warranted; the mechanism is contested separately (C3c's existence half).
  \item \textbf{warranted; rival open.} Warranted, with one named rival left
        explicitly open (C3c's mechanism half).
  \item \textbf{partially closed.} Some factors of the claim are settled
        (often by direct simulation) while at least one named quantity or
        critical question is still open (C3b, C3d, C8).
  \item \textbf{defeated.} A rival kept visible so that its defeat is
        explicit, rather than deleted. Four nodes carry it.
  \item \textbf{open rival.} A rival that is \emph{not} defeated: the
        appendix's one live open question, and the only status the graph marks
        toggleable. Exactly one node carries it: the $\feff(d)$ account of the
        knee, which C3c argues against and does not eliminate. It is what the
        Dung projection reports as the single two-way fork under full stress.
\end{itemize}

A status is \emph{qualified} whenever something remains open, and this appendix
never prints a bare ``warranted'' or ``proven'' over an open critical question;
where it once did, the table now carries the qualifier and the per-node section
names the open item by the identifier the graph uses for it. The partition is the
graph's: a node with no undischarged exception carries \emph{proven},
\emph{warranted}, \emph{resolved} or \emph{defeated}; one with an undischarged
exception must carry a qualified status.

\textbf{CQ numbering is local.} The CQ numbers in this appendix are \emph{this
appendix's}, not WRM~2008's, and the two differ. For Argument from Expert Opinion
this appendix numbers CQ1 expertise, CQ2 assertion/backup, CQ3 consistency, CQ4
trustworthiness/independence, CQ5 the steric criterion; WRM~2008 (p.~15) numbers
CQ1 expertise, CQ2 field, CQ3 opinion, CQ4 trustworthiness (bias as~4.1), CQ5
consistency, CQ6 backup evidence. For C2's analogy CQs the mapping is
local~CQ1 (relevant differences) $=$ Walton~CQ3, local~CQ2 (base-case truth) $=$
Walton~CQ1, local~CQ3 (contrary instances) $\approx$ Walton~CQ4
(counter-analogy); Walton~CQ2 has no local-numbered counterpart because it is
discharged by the similarity premises rather than by a numbered question. The
local numbering is retained (``CQ5'' is the caption adjudication and is cited a
dozen times), and the mapping is stated here instead.

Where a node bundles a proven identity with an empirical calibration, so that
part of it is deductive and part defeasible (e.g.\ C2's exact cubic together with
its fitted prefactor, or C6's algebraic exponents together with their measured
values), the status names the \emph{weaker}, defeasible component; the proven
part is recorded in the per-node detail. That rule is about the parts of a
\emph{single} node. It does not carry to an \emph{accrual}, whose premises are
other nodes: an accrual adds no inference of its own beyond the conjunction, so
it inherits a weaker premise's conditions as conditions \emph{of its statement}
rather than as its own status. The thesis is the case in point, and it says so
where it is argued.

The four columns are the node label; the claim in one line; the status, with the
governing Walton scheme and the convergent support; and the open critical
questions that constitute the node's residual risk.

\subsection{The backbone}

The argument has two storeys. The lower one is a single chain with corroborating
side-branches and one deliberately disconnected lemma; the upper one spends that
chain on an engineering claim:
\[
\begin{gathered}
  \text{C1 (mechanism)} \;\longrightarrow\; \text{C2 }(\delta=w_c\sqrt{\fgap/3})
  \;\longrightarrow\; \text{C3}'\,(\dcrit=\kapp\,\delta), \\[2pt]
  \text{with C5, C6, C7 corroborating} \quad\text{and}\quad
  \text{C4: } F \indep \fgap \ \text{a separation lemma;} \\[4pt]
  \{\text{C2},\,\text{C3a},\,\text{C4},\,\text{S1}\} \;\longrightarrow\;
  \text{C8 (design map)} \;\longrightarrow\;
  \{\text{C3}',\,\text{C4},\,\text{C8},\,\text{scope}\} \;\longrightarrow\;
  \text{THESIS},\\[2pt]
  \text{with C9 (transport lemma) underpinning C1 and C3d.}
\end{gathered}
\]
Mechanism (C1) fixes that the cutoff is set by the diverted-lamina width; the
hydrodynamic conversion (C2) turns that width into a closed form; the operating
law (C3$'$, decomposed into C3a--C3d) promotes it to a particle-diameter cutoff
and calibrates its prefactor; the Xia scaling (C5), the exponents (C6), and the
Dinh consistency band (C7) corroborate; and the filtrate-fraction identity (C4)
is proved \emph{independent} of the gap fraction, so device yield carries no
cutoff information. On the upper storey, that independence is what makes the
target-to-geometry map locally invertible (C8), and the thesis is an accrual over five
legs: the operating law, the separation lemma, the locally invertible map, its
realization, and the declared scope; the display above writes C8 once for the
last two, which the graph carries as two premises. C9 is the odd one out: it derives
$\ksteric=2$ from the flow model rather than from citation, which \emph{upgrades}
C1's premise without adding a branch, and it is therefore attached to the chain
without lying on any path to the thesis.

Figure~\ref{fig:argmap} renders both storeys in one view, which is the map to
read first. Each of the upper-storey nodes then has a detailed map of its own:
Figure~\ref{fig:c9-argmap} for the transport lemma,
Figure~\ref{fig:c8-argmap} for the design map, and
Figure~\ref{fig:thesis-argmap} for the thesis and its five legs, which is also
where the corroborating branches are shown for what they are: nodes with no
support path to the thesis at all.

\begin{figure}[t]
\centering
% D69: the map is scaled to 0.95, so that it fits a 6in text block, and its long caption
% is set in \footnotesize, so that the float fits a 9in text height. (\caption resets
% the size to \normalsize, so \normalsize is redirected inside this float only.)
\scalebox{0.95}{%
\begin{tikzpicture}[
  font=\scriptsize,
  cl/.style={draw, rounded corners, align=center, inner sep=2.5pt,
             text width=2.3cm, minimum height=8mm},
  warr/.style={cl},
  prov/.style={cl, double, double distance=0.8pt},
  part/.style={cl, dashed},
  fl/.style={-{Stealth[length=2mm]}, thick},
  co/.style={-{Stealth[length=2mm]}, dashed},
  nt/.style={font=\scriptsize\itshape, text=black!60, align=center},
]
% --- nodes ---
\node[warr] (c1) at (0,0)        {\textbf{C1}\\ cutoff $=$ lamina width};
\node[warr] (c2) at (0,-1.8)     {\textbf{C2}\\ $\delta=w_c\sqrt{\fgap/3}$};
\node[warr] (c3a) at (0,-3.6)    {\textbf{C3a}\\ $\dcrit\propto\delta$};
\node[part] (c3b) at (-3.3,-5.6) {\textbf{C3b}\\ plateau $\kapp=\plateauKappa$};
\node[part] (c3c) at (0,-5.6)    {\textbf{C3c}\\ knee is real};
\node[part] (c3d) at (3.3,-5.6)  {\textbf{C3d}\\ decompose $\kapp$};
\node[warr] (c5) at (6.9,-1.8)   {\textbf{C5}\\ Xia $\fgap\!\propto\!w_c^{-2}$};
\node[warr] (c6) at (6.9,-3.6)   {\textbf{C6}\\ $\beta=\tfrac12,\ \gamma=1$};
\node[warr] (c7) at (6.9,-5.6)   {\textbf{C7}\\ Dinh band $\subset$ bracket};
\node[prov] (c4) at (0,-8.1)
  {\textbf{C4}\\ $F\indep\fgap$ separation\\ lemma \emph{(proven)}; $F$'s\\
   magnitude \emph{validated}};
% --- ABOVE THE CUTOFF CHAIN: the design map, the scope and the thesis are the
% upper storey proper; C9 belongs to neither storey and is drawn here only
% because there is room. See the Introduction's two-storey paragraph.
\node[prov] (c9) at (-4.6,0)     {\textbf{C9}\\ transport lemma\\ $\ksteric=2$};
\node[part] (c8) at (4.8,-8.1)   {\textbf{C8}\\ design map\\ targets $\to$ geometry};
\node[prov] (sc) at (-6.3,-8.1)  {\textbf{Scope}\\ rigid spheres,\\ Stokes flow,\\
   const.\ viscosity};
\node[warr, text width=11.0cm, minimum height=9mm] (th) at (0,-10.6)
  {\textbf{THESIS.}\ CIF is engineerable from its two performance
   targets, cutoff and yield, because their laws are \emph{provably decoupled},
   making the two independent design inputs,
   \emph{inside a declared scope and a checked envelope}};
% --- grouping box for the operating law ---
\node[draw, dotted, rounded corners, fit=(c3a)(c3b)(c3c)(c3d),
      inner sep=3mm] (box) {};
% label moved to the box's top-right: the bottom-left corner is now the origin of
% the operating-law -> thesis edge, and the two collided there.
\node[nt, anchor=south east] at (box.north east) {C3$'$ operating law (C3a--C3d)};
% --- backbone (deductive / hydrodynamic chain) ---
\draw[fl] (c1)--(c2);
\draw[fl] (c2)--(c3a);
\draw[fl] (c3a)--(c3b);
\draw[fl] (c3a)--(c3c);
\draw[fl] (c3a)--(c3d);
% --- corroboration ---
\draw[co] (c5)--(box.east);
\draw[co] (c6)--(box.east);
\draw[co] (c7)--(box.east);
% --- upper-storey edges ---
% The two edges leaving the C3' cluster are LABELLED, because not every member of
% C3' feeds every consumer: C3d interprets the prefactor's magnitude and is a
% premise of neither C8 nor the thesis. An unlabelled cluster arrow would assert
% that it is.
\draw[fl] (c9) -- (c1);
\draw[fl] (box.south east) -- node[nt, pos=0.42, right=1mm]{C2, C3a, S1} (c8);
\draw[fl] (c4) -- (c8);
\draw[fl] (box.south west) -- node[nt, pos=0.40, right=1mm]{C2, C3a--C3c}
          ([xshift=-2.8cm]th.north);
\draw[fl] (c4) -- (th);
\draw[fl] (c8) -- (th);
\draw[fl] (sc) -- (th.west);
% --- annotations ---
\node[nt, anchor=west] at (1.8,0)     {Dinh caption\\ deselected};
\node[nt, anchor=east] at (-4.6,-1.5) {underpins C1 here, and\\ C3d's decomposition\\ (edge not drawn);\\ no path to the thesis};
\node[nt, anchor=west] at (2.1,-6.9)  {decoupled:\\ $\partial F/\partial \fgap=0$};
\node[nt, anchor=north, text width=12.6cm] at (1.5,-11.6)
  {\textbf{Open critical questions (residual risk).}\ \,
   C3b: the cross-paper anchor is circular \emph{and} its $\xiaAgreePct\%$ agreement is
   numerically the depth correction, so it discriminates nothing and is reported
   as a \emph{lead}; the fit itself is unaffected.\ \,
   C3c: the $\feff(d)$ alternative is bounded, not eliminated, and the knee
   \emph{mechanism} is deferred to future finite-particle work.\ \,
   C3d: the lumped residual, depth-separated
   $\sqrt{c}\,C=\residualSepShallow/\residualSepMid/\residualSepDeep$ over
   $\alpha=0.80/1.40/2.69$ (a triplet, not a range), tied to the
   unpublished gap count $N$.\ \,
   C4: a $\residPctVOne\%$ (V1) / $\residPctVThree\%$ (V3) 1-D-vs-3-D residual on $F$'s
   \emph{accuracy} (it does not
   couple $F$ to $\fgap$).\ \,
   C8: three undercutters, \emph{withheld} rather than excluded:
   \texttt{below\_knee}, \texttt{runout\_mismatch} and
   \texttt{pinned\_row\_count}, which is what the dashed border on C8 records,
   and which Figure~\ref{fig:c8-argmap} argues in full.};
\end{tikzpicture}}
\let\normalsize\footnotesize
\caption{Argument map for the whole argument: the visual companion to the
backbone above. The \emph{lower storey} is the cutoff law: C1$\to$C2$\to$C3$'$
down the spine, with C5, C6 and C7 corroborating into it and C4 drawn
deliberately disconnected from it. The \emph{upper storey} spends that
disconnection: because $F$ carries no cutoff information, the two design targets
are independently specifiable, which is what makes the design map C8 locally invertible,
and the thesis is the accrual of the operating law, the separation lemma, the
design map (its inversion C8a and its realization C8b, drawn here as one
node) and the declared scope: the thesis's five legs, in four arrows. The two
edges leaving the C3$'$ cluster are labelled, because C3d interprets the
prefactor's magnitude and is a premise of neither consumer; an unlabelled cluster
arrow would assert that it is. C9 stands to one side of the spine on purpose: it
derives $\ksteric=2$ and so \emph{underpins} C1, and C3d's decomposition, of
which only the C1 edge is drawn; no support path runs from it to the thesis. Each upper-storey node has a detailed map of its own
(Figures~\ref{fig:c9-argmap}, \ref{fig:c8-argmap} and
\ref{fig:thesis-argmap}); this one shows how they connect.
Border style encodes status: a single border marks a \emph{warranted} node, a
dashed border (C3b, C3c, C3d, C8) a node with something still open, and a doubled
border (C9, C4, Scope) a \emph{proven} identity or declaration. C4's doubled
border covers the separation itself; $F$'s absolute \emph{magnitude} is a
separate, empirical claim, which is where C4's residual at the foot of this figure
attaches. Those three words are this map's convention and this map's alone; the
per-node figures that follow use the dashed border for whatever \emph{that} figure
declares in its own caption (an adjudicated question in
Figure~\ref{fig:c1-argmap}, an excluded rival in Figure~\ref{fig:c5-argmap}, a
withheld undercutter in Figure~\ref{fig:c8-argmap}), and each says which. Solid
arrows are support; dashed arrows are corroboration, with one exception the
spine makes unavoidable. C3a$\to$C3b, C3a$\to$C3c and C3a$\to$C3d are premise
edges the graph carries; C1$\to$C2 and C2$\to$C3a are \emph{reading order}, not
graph edges: \texttt{mechanism} is not a premise of \texttt{c2\_skim}, whose
premise is the parabolic gap profile; \texttt{skim\_law} is not a premise of
\texttt{c3a}, whose premise is expert proportionality; and C2 and C3a enter the
operating law as siblings rather than in series. The chain is how the argument is read and how the physics composes;
the graph's support structure branches at C3$'$. Unlike Appendix~A's
Figure~1, which is a derivation \emph{dependency} graph, this map shows the
\emph{support and defeat} structure: each node's grounds and the open critical
questions that could still defeat it.}
\label{fig:argmap}
\end{figure}

\subsection{Notation and conventions}
\label{sec:b-notation}

The subject law is
\[
  \dcrit = \kapp\!\left(d/w_c,\ \mathrm{depth}\right)\, w_c \sqrt{\fgap/3},
  \qquad
  \delta = w_c\sqrt{\fgap/3},
\]
where $\delta$ is the deep-channel skim width: the leading term of the exact
dividing-streamline inversion of the skim cubic $\fgap = 3x^2 - 2x^3$
(the Inglis $\varepsilon$-cubic, \cite[p.~656]{b17}; $x=\delta/\wc$). Here
$\kapp$ is the \emph{apparent} prefactor, which maps that idealized skim width
onto the measured cutoff diameter; it is defined against the \emph{nominal}
design $\fgap$ and the parabolic $\delta$. \textbf{$\kapp$ is a plateau-region
constant, not a universal one.} It is constant only for
$\dwc \gtrsim \plateauUwc$, where $\kapp = \plateauKappa$; below that the
prefactor becomes a function of the cutoff itself and the law above becomes
\emph{implicit} in $\dcrit$, which is why C7's bracket device, below the knee at
its initial gap fraction, has no stable fixed point above the sub-steric foot to
solve for. The prefactor ramps through that knee
to a \emph{sub-steric (recovery-suppressed)} foot $\kapp \approx \footKappa$
($\dwc < \footUwc$), which is \emph{below} even the \emph{depth-corrected} steric
floor $\ksteric\Kdcfd \approx \stericFloorShallow$ at the shallowest probed aspect
ratio (the bare steric floor is $\ksteric = 2$; the two are distinguished
wherever both appear), because small-particle recovery is incomplete
($\delta_{\mathrm{eff}}<\delta_{\mathrm{nominal}}$; C3c). Every numerical use of
$\kapp$ in this appendix is a plateau-region use unless it says otherwise. The
plateau prefactor factors as
\[
  \kapp = \ksteric \cdot \kinplane \cdot \Kdepth \cdot \sqrt{c}\cdot C
        = 2 \cdot 1 \cdot R(\alpha)\cdot \sqrt{c}\cdot C ,
\]
with $\ksteric = 2$ (cutoff radius $=$ lamina width, hence $D_c = 2b$, where
$b$ is Inglis's first-streamline width, re-lettered in the second convention
below; corroborated by \cite[p.~656]{b17} and \cite[p.~917]{b5} main text, and
\emph{derived} from the flow model in C9 rather than taken on citation),
$\kinplane = 1$ (single-gap CFD, \cite{b7}), and $R(\alpha)=\Kdepth$ the
finite-depth factor \cite[p.~875]{b6}. These three, and the residual
$\sqrt{c}\,C$, are each carried in the derivation record \cite{b12a}, which is
the numerical authority for their values.

Five conventions are fixed here to forestall confusion.

\begin{itemize}
  \item \textbf{C7 and C8 are the \emph{graph's} numbers, not the prospectus's.}
        The labels C1--C9 in this appendix are the argument graph's node
        identifiers: the graph took them from this appendix, and holds the
        canonical \texttt{Graph-*} numbering. Two of them collide
        with prospectus contribution numbers: in \emph{this appendix} C7 is the
        Dinh held-out consistency band and C8 is the design map, whereas the
        \emph{prospectus} C7 is the targets-to-photomask tool and its C8 the
        machine-enforced validity envelope. Different objects, the same
        numerals. The graph resolves this by namespacing rather than
        renumbering, and asks for explicit \texttt{Graph-}/\texttt{Prospectus-}
        prefixes wherever either is cited. This appendix adopts the resolution
        and narrows the prefix: the two section headings that carry the collision
        say ``(graph)'' in the heading itself, and everywhere else here an
        unqualified C7 or C8 is the graph's, because a reader of this appendix
        meets a prospectus number in exactly one place, the thesis section's
        scope note, which is also where the reason the prospectus pair carries no
        argument node at all is recorded.
  \item \textbf{First-streamline width vs.\ the size-scaling exponent.} Two
        distinct quantities are written with the letter $\beta$/$b$ across these
        literatures and must not be conflated. Inglis (2006) denotes the
        first-streamline width by the Greek letter $\beta$ (his Fig.~1 caption:
        ``$\beta$ is the width of the first streamline,''
        \cite[p.~655, Fig.~1 caption]{b17}), with critical diameter
        $D_c = 2\beta$ \cite[p.~656, eqn~(1)]{b17}. To avoid collision with the
        CIF size-scaling exponent
        $\beta = \tfrac12$ ($=\partial\ln\dcrit/\partial\ln\fgap$) used throughout
        this manuscript, we render Inglis's width as lowercase Latin $b$ wherever
        it appears; thus his $D_c = 2\beta$ is written $D_c = 2b$. The
        re-lettering is purely notational and leaves $\ksteric = 2$ unaffected.
  \item \textbf{The lumped residual, two equivalent conventions.} The data-fixed
        excess of the apparent prefactor over the steric value can be written
        with the depth factor $R(\alpha)$ either pulled out or absorbed; both
        denote the same number. Keeping $R(\alpha)$ explicit (the convention of
        this appendix and Appendix~A),
        $\Kdepth\,\sqrt{c}\,C = \kapp/\ksteric = \residualFolded$ (depth-folded,
        and $\alpha$-independent), so that dividing out the \emph{measured}
        $\Kdcfd$ leaves the depth-separated
        $\sqrt{c}\,C = \kapp/(\ksteric\Kdcfd) =
        \residualSepShallow/\residualSepMid/\residualSepDeep$ at
        $\alpha = 0.80/1.40/2.69$: a triplet, not a range, and each value
        belongs to its own $\alpha$. The reference code
        \texttt{cif\_cutoff.py} \cite{b12a} instead \emph{folds} $\Kdepth$ into
        the residual
        and reports $\sqrt{c}\,C = \kapp/\ksteric = \residualFolded$; the two are
        the same quantity with $\Kdepth$ either pulled out (here) or absorbed
        (there), and the glossary's rule applies: \emph{never print the two
        conventions as one range}.
  \item \textbf{Aspect ratio.} Throughout this appendix $\alpha \equiv h/\wc$,
        depth over central-channel width, and it is \emph{not} sorted to
        $\alpha \ge 1$. The reference mathematics takes the same ratio as a
        parameter literally named \texttt{beta} (\texttt{c2(beta)},
        \texttt{K\_depth(beta)}, evaluated at $0.80/1.40/2.69$): when reading
        code, \texttt{beta} is this appendix's $\alpha$, the same quantity
        under a different name, \emph{not} a reciprocal. When reading this
        appendix, $\beta$ is never an aspect ratio; it is the cutoff-law scaling
        exponent. One further correspondence with the glossary, stated so that a
        reader moving between them is not misled: this appendix's $\dwc$ is
        the glossary's $\ud$, the same quantity under a different spelling. The
        steric prefactor, the glossary's $\kster$, is printed the same way here.
  \item \textbf{Depth factor: three triplets, kept apart.} Appendix~A carries
        \emph{three} depth-factor triplets and this appendix follows its symbols.
        Two are taken over the aspect ratios probed by direct simulation
        ($\alpha = 0.80/1.40/2.69$): $\Kdth$, the \emph{analytic}
        rectangular-duct factor \cite[p.~875]{b6}
        ($\KdepthAnalyticShallow/\KdepthAnalyticMid/\KdepthAnalyticDeep$), and
        $\Kdcfd$, the \emph{measured} single-gap CFD factor
        ($\KdepthCfdShallow/\KdepthCfdMid/\KdepthCfdDeep$). They agree to at
        most $\KdepthDevMax\%$, and where only the band matters this appendix
        writes $R(\alpha)\in[0.89,0.98]$ for either. The \emph{third} triplet is
        different in kind, and it is where this appendix's one acknowledged
        assumption lives: the analytic factor evaluated not at those aspect ratios
        but at the three Gifford Fig.~3 \emph{calibration widths}. That figure's
        caption reports the widths and nothing else: ``Effect of central channel
        width ($w_c$ = 100 \um, 125 \um, 150 \um)''
        \cite[p.~4500, Fig.~3 caption]{b1}. So the calibration array has
        \emph{no stated depth}. The $80\um$ used to form $\alpha$ is the depth of
        the \emph{separate} Fig.~4 demonstration device: ``The depth of the device
        was 80 \um\ $\pm$ 3 \um; the width of the central channel was
        100 \um'' \cite[p.~4500, Fig.~4 caption]{b1}: the caption specifically,
        because the body sentence on that same page states the depth without the
        tolerance. Carried across the three widths it gives
        $\alpha = 0.8000/0.6400/0.5333$ and
        $\Kdth = \KdepthCalNarrow/\KdepthCalMid/\KdepthCalWide$, an
        $11$--$19\%$ correction, \emph{not} negligible, and largest at the widest
        channel. The true array depth is not known and is not inferred here
        (the glossary, ``Calibration depth''), and $80\um$ is not the only depth
        the paper reports: the four high-ratio devices of its Fig.~5 are
        ``Four 125 \um-deep CIF devices with values of $\fgap$ ranging
        $2.2$--$10.0\times10^{-4}$ were tested''
        \cite[p.~4501, Fig.~5 caption]{b1}, at the \emph{same} $100\um$ central
        channel as the Fig.~4 device, and the Fig.~6 platelet-enrichment device is ``The
        depth of the device was 150 \um\ $\pm$5 \um; the width of the central
        channel was 125 \um'' \cite[p.~4502]{b1}. Because $\Kdepth$ rises
        monotonically towards $1$ with $\alpha$, either of those deeper
        assumptions would make the correction both \emph{smaller} and
        \emph{more uniform} across the three widths. Taking the shallowest
        reported depth is therefore the \emph{conservative} reading, and, if the
        array is at least that deep, the $11$--$19\%$ above is an upper bound rather than an estimate. A number from one triplet must never
        be quoted against another. Throughout, $\Kdepth$ is folded into the
        \emph{apparent} prefactor as a per-device interpretive factor; it is
        \emph{not} applied to the data, which collapse more tightly on the clean
        parabolic $\delta$, so it does not shift the fitted $\kapp$ and enters
        only the residual decomposition.
\end{itemize}

\section{Claim-by-claim argumentation table}

Table~\ref{tab:argmap} is the proof summary: every node of the backbone, its
status under its Walton scheme, and its residual risk. The per-node premises,
verbatim source quotations, and full critical-question lists are collected in
the detailed treatment that follows this summary.

{\footnotesize
\setlength{\tabcolsep}{4pt}
\renewcommand{\arraystretch}{1.15}
\begin{longtable}{@{}L{0.05\textwidth}L{0.21\textwidth}L{0.32\textwidth}L{0.32\textwidth}@{}}
\caption{Claim-by-claim argumentation map for the whole argument: the cutoff law
(C1--C7), the transport lemma (C9), the design map (C8) and the thesis. Status
levels are defined in \S\ref{sec:b-howtoread} and taken from the argument
graph's own legend:
\emph{proven}, \emph{warranted}, \emph{resolved} and \emph{defeated} carry no open
critical question; \emph{validated}, \emph{warranted as a band}, \emph{warranted
(existence)}, \emph{warranted; rival open}, \emph{partially closed} and
\emph{open rival} are the graph's \emph{qualified} statuses, each of which may
stand over an undischarged exception, and the right-hand column says which. That
partition is the graph's own (\texttt{meta.convention}), not this appendix's. A
status may pair two of these where a node bundles two claims, as C4 does:
\emph{proven} for the separation and \emph{validated} for the magnitude, the
latter carrying the node's one open exception.}
\label{tab:argmap}\\
\toprule
\textbf{Node} & \textbf{Claim} & \textbf{Status (Walton scheme; support)} &
\textbf{Open critical question (residual risk)}\\
\midrule
\endfirsthead
\multicolumn{4}{@{}l}{\itshape Table~\ref{tab:argmap} (continued)}\\[2pt]
\toprule
\textbf{Node} & \textbf{Claim} & \textbf{Status (Walton scheme; support)} &
\textbf{Open critical question (residual risk)}\\
\midrule
\endhead
\midrule
\multicolumn{4}{r@{}}{\itshape continued on next page}\\
\endfoot
\bottomrule
\endlastfoot

C1 &
The cutoff is set by the diverted-lamina width, not the gap width
($a \lesssim w$, with Dinh's $w$ being our $\delta$; $\dcrit \approx 2\delta$). &
\textbf{Warranted.} Expert Opinion (WRM~2008), convergent over five
peer-reviewed sources spanning \emph{three} independent groups (the graph's own
\texttt{accrual\_over} list): the originating CIF lab (Gifford \cite{b1};
Lezzar \cite{b4}; Dinh \cite{b5}: all from the Shevkoplyas lab) plus two independent
confirmations (Yamada \& Seki \cite{b15}; Inglis \& Herman 2013 \cite{b18}). &
$a<w$ vs $a<w/2$ \emph{resolved} $\to \ksteric = 2$ (Dinh \emph{main text},
\cite[p.~917]{b5}, $+$ Inglis's $D_c = 2\beta$, our $2b$, \cite[p.~656]{b17}); the Fig.~1
``about half'' caption \cite[p.~916]{b5} is \emph{deselected}, not refuted.
Residual: none on the mechanism. \\
\midrule

C2 &
$\delta = w_c\sqrt{\fgap/3}$, the leading term of the exact cubic
$f = 3x^2 - 2x^3$. &
\textbf{Warranted on the form.} A proven, code-checked identity (deductive),
reinforced by Analogy (DLD, \cite[p.~656]{b17}) and Expert Opinion (Inglis \&
Herman's identical near-wall-lamina integral, \cite[p.~1726]{b18}); exact,
Inglis-identical, CFD-confirmed ($\delta_{\mathrm{skim}}\propto\sqrt{f}$,
$R^2\!\ge\!0.9999$, and agreement with the exact rectangular-duct field to
$\le\!0.6\%$, both \cite{b7}). &
None on form. The cutoff is fit on the parabolic $\delta$; the finite-depth
factor $R(\alpha)$ is folded into the apparent $\kappa$, not applied to the data.\\
\midrule

C3a &
$\dcrit \propto$ lamina width ($\kappa \equiv \dcrit/\delta$). &
\textbf{Warranted.} Expert Opinion over four sources (Xia assumption~(b),
\cite{b2}; Inglis $D_c = 2b$, \cite[p.~656]{b17}; Dinh's C1 criterion, \cite{b5};
Inglis \& Herman's design assumption, \cite{b18}) \emph{and} Argument from Example:
the measured exponents select the form, $\beta = \expBeta$ excluding the Hele-Shaw
$\beta = 1$. The second leg is what keeps the operating law, C8 and the design map
from resting on expert testimony alone: they survive the withdrawal of every
source. &
The \emph{value} of the constant $\to$ C3b--C3d. \\
\midrule

C3b &
Plateau $\kapp = \plateauKappa\ [\kappCIlo, \kappCIhi]$, $\beta = \expBeta$,
$\gamma = \expGamma$. &
\textbf{Partially closed: warranted on CIF data; the cross-paper anchor
is a lead, not a corroboration.} Argument from Example (the error-weighted fit
generates the premise):
CV $\kappCV\%$; CIs inflated by $\sqrt{\chired}$ ($\chired = \chiPlateau$), hence
conservative; survives a 2000-draw digitization Monte Carlo
($P(0.4\!<\!\beta\!<\!0.7)=1.00$) and is invariant across three independent
digitizations (the primary and two re-digitizations), which return $\kapp$ in $[4.91, 4.92]$ and reject the global law
every time. None of that depends on the Xia comparison. &
\texttt{xia\_anchor\_unverified}. The recalculation is real and its inputs are
quoted, but it is \emph{not} a held-out test, for the two reasons CQ2 separates, after a
caveat that is the graph's own:
the $\sim\!140\um$ figure is the channel \emph{depth} while the central-channel
\emph{width} tapers $300 \to 150\um$, so the literal device aspect ratio is
uncertain, which is why \texttt{cif\_cutoff\_future\_work.py} flags the anchor
\textsc{unverified}. It is also circular: Xia \cite{b2} takes its $7\um$ cutoff
from \emph{the same Gifford Fig.~3} \cite[p.~4500]{b1} this plateau was fitted
from, and Gifford is a Xia co-author. Decisively, the $\xiaAgreePct\%$ agreement
\emph{is} the depth correction: $1/\Kdepth(1.40) - 1 = \depthCorrPct\%$, and
$\kappVal/\KdepthAnalyticMid = \xiaKappaMean$, so the comparison is
equally consistent with a depth-robust prefactor and with a depth correction of
that size and discriminates between them not at all. Reported as a \textbf{lead};
the fit and the exponents stand on the Gifford three-width data alone. \\
\midrule

C3c &
The cutoff knee is real. &
\textbf{Warranted on CIF data (existence); mechanism rival open.} Argument from
Sign for the \emph{existence}, Abduction for the \emph{mechanism} (Doyeux,
\cite[p.~359]{b21}); the digitization undercutter is closed separately, by
Evidence-to-Hypothesis in falsification form. The global single power law is
rejected at reduced $\chi^2 = \chiGlobal$, and $\kapp$ ramps from $\plateauKappa$
to $\footKappa$ as $d/w_c\!\to\!0$. Both figures are \emph{binned means}: the
$n=\plateauN$ plateau bin ($\dwc \ge \plateauUwc$) and the $n=\footN$ foot bin
($\dwc < \footUwc$); the
fifteen individual points span $\kappaSpanLo$--$\kappaSpanHi$, which is the span
and never the ramp. Foot and rejection robust across all three digitizations. &
\texttt{more\_reliable} (existence) and \texttt{more\_coherent\_explanation}
(mechanism). \emph{Mechanism}: the knee is attributed to a finite-particle capture-regime change, deferred
to future finite-particle work; the $\feff(d)$ alternative (the content of
\texttt{more\_reliable}) is bounded but not eliminated; DLD is silent, its
\emph{apparent} prefactor $\kapp$ running $\dldKappaLo$--$\dldKappaHi$ with no knee
and never below the steric floor $\ksteric=2$. Those two bounds are \emph{two
computations of the same DLD prefactor, not two laboratories}, and they run in
opposite senses: Davis's eqn~2-22 falls with $\varepsilon$ from $\davisKappaHi$ to
$\davisKappaLo$ \cite[eqn~(2-22)]{b22}, while the exact Inglis cubic rises from
$\inglisKappaLo$ to $\inglisKappaHi$ \cite[p.~656]{b17}. \\
\midrule

C3d &
Physical decomposition of $\kapp$. &
\textbf{Partially closed.} Numerical verification (CFD, \cite{b7}; deductive
computation) $+$ Established Rule $+$ Expert Opinion (Doyeux, \cite[p.~359]{b21})
$+$ Abduction. Settled: $\ksteric = 2$, $\kinplane = 1$ (CFD, Runs A/B/C: no
in-plane funnelling), $R(\alpha)\in[0.89,0.98]$. Data-fixed residual
$\kapp/\ksteric = R(\alpha)\,\sqrt{c}\,C = \residualFolded$ (depth-folded);
dividing out the \emph{measured} $\Kdcfd$ leaves the depth-separated
$\sqrt{c}\,C = \residualSepShallow/\residualSepMid/\residualSepDeep$ at
$\alpha = 0.80/1.40/2.69$ open, \emph{a triplet and not a range}
(\S\ref{sec:b-notation}), each value belonging to its own $\alpha$
(per-gap over-extraction $\times$ cumulative finite-particle displacement, the
Doyeux mechanism), tied to the unpublished gap count $N$. Every factor is carried
in the derivation record \cite{b12a}. &
\texttt{more\_coherent\_explanation} (the mechanism) and
\texttt{degenerate\_from\_published\_data} (the split: single-lineage data fix the
\emph{product} $\sqrt{c}\,C$ and not its division into $c$ and $C$). Needs
finite-particle work $+$ a micrograph gap count. \emph{Not} the single-gap
solve (done). \emph{Not} the $F$-identity (decoupled, C4): reading $\kappa$ out
of $F$ inverts the C4 lemma. \\
\midrule

C5 &
The skim law derives the corrected Xia scaling $\fgap \propto w_c^{-2}$. &
\textbf{Warranted.} Evidence-to-Hypothesis (verification form): the law predicts
the \emph{corrected} direction and excludes the published original \cite{b2},
confirmed by the 2022 author correction \cite[p.~1]{b3}. &
Shares premises with C2/C3a: a consilience (one relationship), not a fully
independent test. \\
\midrule

C6 &
$\beta = \tfrac12$, $\gamma = 1$. &
\textbf{Warranted; uniform-$\feff$-robust.} Defeasible Modus Ponens on the C2
identity and Argument
from an Established Rule for the shift-invariance (both deductive, Appendix~A,
Result~6), corroborated by Argument from Example twice over: the CIF plateau fit
($\beta = \expBeta$, $\gamma = \expGamma$) and Davis's independent DLD device
family (exponent $0.48$, P-C6.5). The shift-invariance is \emph{algebraically exact, not
approximate}: identical at $c = 1, 2, 3.5$. &
--- \\
\midrule

C7 &
The prediction for Dinh's $\fgap^{*}=1.04\times10^{-4}$ device sits inside the
measured bracket: band $[\dinhBandLo, \dinhBandHi]\um \subset (\dinhBracketLo, \dinhBracketHi)\um$
\cite[ESI, Fig.~S1]{b5}. &
\textbf{Warranted as a band.} Numerical computation for the band (deductive),
checked against data that was not fitted (Evidence to a Hypothesis). Corroboration
only, not a discriminating test. &
None \emph{open}: the graph records no open exception here and every critical
question is answered. What the answers establish is how little the band can be
asked to carry: it is doubly sensitive to depth and $\feff$ (at its initial
gap fraction that device sits below the knee, where
$\dcrit = \kappa\,\delta$ has no stable fixed point \emph{above the
sub-steric foot}), and it is weak ($4.7\times$ bracket). The interval
$c_{\mathrm{Dinh}}/c_{\mathrm{Gifford}}\in[0.64, 4.69]$ over which the band stays
inside $(\dinhBracketLo, \dinhBracketHi)\um$ is a \emph{tolerance}, not a plausibility range: it is
$7.3\times$ wide because the bracket is, and CQ4 reads the same insensitivity as
weakness. Anchor on plateau devices, not Dinh. \\
\midrule

C4 &
$F \indep \fgap$ (separation lemma): $F = 2R_c/(R_s + 2R_c)$ depends only on
outlet geometry. &
\textbf{Proven} (the separation); \textbf{validated} ($F$'s magnitude). The
decoupling $F \indep \fgap$ is a deductive current-divider identity, carried by
Argument from an Established Rule (the hydraulic--electric analogy), and carries
no open critical question; $F$'s
absolute magnitude is a separate, empirical claim, confirmed to within $1\sigma$
by COMSOL \cite{b7} on Dinh V1/V3 in the companion preprint \cite{bpre}, and it is that claim the
residuals at right attach to. Neither bears on the decoupling. &
The $\residPctVOne\%$ (V1) / $\residPctVThree\%$ (V3) 1-D-vs-3-D residual is on
$F$ itself (closed by 3-D outlet-bifurcation physics) and does \emph{not} couple
$F$ to $\fgap$;
\texttt{sibling\_3d\_\allowbreak{}junction\_\allowbreak{}uncertainty\_\allowbreak{}on\_\allowbreak{}sqrt\_c}, a
sibling 3-D-junction effect, contributes a $\sim\!5\%$-scale uncertainty on
$\sqrt{c}$, distinct from $F$'s own residual and the node's one open critical
question. \\
\midrule

C9 &
Transport lemma: $\varphi'=(\varphi-f)/(1-2f)$ makes the centreline a
\emph{repelling} fixed point, so selection is a hard wall constraint rather than a
statistical criterion; with centre exclusion (P-C9.1$'$) that gives
$\ksteric = 2$ exactly. &
\textbf{Proven.} Defeasible Modus Ponens on a closed-form migration map, with
Argument from an Established Rule for its three further consequences; the
critical position reproduces to six decimals
($\varphi_{\mathrm{crit}} = 0.216643$). Deliberately \emph{not} load-bearing:
no support path runs from $\ksteric = 2$ to the thesis. &
None. Within the model's premises, the lemma also \emph{closes} C3d's residual list to three
model gaps (the
flat-lamina idealisation, the wall porosity $\phi$, and finite-particle
displacement), which answers C3d's exhaustiveness question while leaving its
identifiability question open. \\
\midrule

C8 &
The design map: the target Jacobian is lower-triangular
($\partial\dcrit/\partial N = 0$) with non-zero diagonal, so targets invert
locally to geometry and forward substitution realizes both on a device inside the
envelope. &
\textbf{Partially closed.} The \emph{inversion} (C8a) is proven, resting on C2,
C3a, C4 and the S1 schedule, and is reached a second time, by construction rather
than by argument, through an executable
forward march ($N\!\cdot\!D^{2}$ holds to $0.29\%$ across five cutoff targets while
$N$ spans $1446$--$6562$). The \emph{realization} (C8b) is Argument from an
Established Rule and rests on the inversion plus three granted assumptions; the
$1.73$-decade decoupling sweep supports C4, not C8b. &
Three undercutters, all \emph{withheld} rather than excluded:
\texttt{below\_knee}, \texttt{runout\_mismatch}, \texttt{pinned\_row\_count}.
The first has published counterexamples: at their wide ends the plateau law puts
Dinh's three designs at $u=\dwc=0.029$--$0.032$ (the bracket device at $0.029$), and
Xia's RS-7 lower still at $0.023$ on the taper Xia prints (C5),
against the floor $u_{\mathrm{lo}}$ set here to the plateau-bin edge $\plateauUwc$, so it
is the guard that holds the defeater dormant. The \emph{need} for the guard is
derived: below the plateau edge the operating law is \emph{implicit} in $\dcrit$
(C7), so a target no longer fixes the cutoff. Its \emph{value} is not: $u_{\mathrm{lo}}$
is an empirical bin edge, an external input (CQ7). \\
\midrule

Thesis &
CIF is engineerable from its two performance targets rather than from a tuned
per-gap schedule: yield follows a closed-form law with no fitted constant and
selectivity a derived law with a single calibrated prefactor, and the two are
\emph{provably decoupled}, so cutoff and yield are independent design inputs
rather than coupled outcomes. &
\textbf{Warranted, conditionally.} Accrual over five legs: the operating law
(C2, C3a--C3c), the separation lemma (C4), the locally invertible map (C8a),
realizability (C8b) and the declared scope. C1 and C9 establish $\ksteric = 2$,
which underpins the mechanism and C3d's decomposition; neither is a
premise of the operating law, and neither reaches this node. Evaluated
under preponderance of the evidence at the weakest leg's tier. &
Inherited, not new: C8's three envelope boundaries, plus the open evidentiary
questions of C3b--C3c and C4. Scope is carried as a premise (rigid spheres in
Stokes flow), so platelet discoidness, deformability, cell crowding and whole
blood are \emph{declared} out of scope, not overlooked. \\
\end{longtable}
}

\section{Net assessment}

The cutoff law is \textbf{proven in form}: the exact skim cubic, its identity
with the Inglis critical-streamline cubic, the $\beta=\tfrac12$ / $\gamma=1$
exponents, and the recovery of the corrected Xia scaling are all established
symbolically and reproduced byte-for-byte by the reference code \cite{b12a}, with
the Xia recovery entering as a \emph{consilience} rather than an independent
test, since it shares its premises with C2 and C3a (C5). It is
\textbf{empirically solid in its apparent plateau prefactor}
($\kapp = \plateauKappa$, $95\%$ CI $[\kappCIlo, \kappCIhi]$, CV $\kappCV\%$),
robust to a 2000-draw digitization Monte Carlo, and returning $\kapp$ in
$[4.91, 4.92]$ across three independent digitizations of the calibration
figure (the primary and two re-digitizations). The cross-paper comparison with Xia \cite{b2} ($\kapp = \xiaKappaMean$, a
$\xiaAgreePct\%$ agreement) once stood here as corroboration and is now reported
as a \textbf{lead}: its $\sim\!140\um$ figure is the channel \emph{depth} while
the central-channel \emph{width} tapers $300 \to 150\um$, so the device aspect
ratio is uncertain; Xia reads its $7\um$ cutoff from the same Gifford Fig.~3
\cite[p.~4500]{b1} this plateau was fitted from; and, decisively, that
$\xiaAgreePct\%$ is numerically the $140/80\um$ depth correction itself, so the
agreement cannot discriminate between a depth-robust prefactor and a depth
correction of exactly that size (C3b, CQ2).
It carries no weight in this assessment, which is why the sentence above rests on
the CIF calibration data alone. Its physical
decomposition is \textbf{partially closed}: $\ksteric = 2$, $\kinplane = 1$
(CFD-settled, \cite{b7}), and $R(\alpha)\in[0.89,0.98]$ over the three
CFD-probed aspect ratios (the analytic duct law \cite[p.~875]{b6} extending to
$\KdepthCalWide$ at the widest calibration channel, $\wc = 150\um$, which is the
lowest aspect ratio of the calibration triplet) are settled, so the data-fixed
residual is $\kapp/\ksteric = R(\alpha)\,\sqrt{c}\,C = \residualFolded$. Dividing
out the measured $\Kdcfd$ leaves a single open quantity, the depth-separated
triplet $\sqrt{c}\,C = \residualSepShallow/\residualSepMid/\residualSepDeep$ at
$\alpha = 0.80/1.40/2.69$, a triplet and not a range
(\S\ref{sec:b-notation}). That quantity is per-gap over-extraction times
cumulative finite-particle displacement, tied to the unpublished gap count $N$. Every
constant here is carried in the derivation record \cite{b12a}. The exponents (by shift-invariance, C6), the
knee's \emph{existence} (a size-dependent $\feff$ is itself size-dependence and
so cannot defeat the variation it would merely reattribute), and the Dinh
consistency band (which stays inside the bracket across
$c_{\mathrm{Dinh}}/c_{\mathrm{Gifford}}\in[0.64, 4.69]$) all survive the
$\feff$ confound; only the absolute prefactor's \emph{decomposition} and the knee
\emph{mechanism} depend on it, and closing it is the task of
finite-particle simulation plus a micrograph gap count, \emph{not} of the
single-gap solve (done) and \emph{not} of the $F$-identity, which is provably
decoupled (C4).

Two scope conditions frame the whole assessment. First, $\kapp$ is
\emph{apparent}: it is defined against the \emph{nominal} design $\fgap$ and the
parabolic $\delta$, so the open residual $\sqrt{c}\,C$ rescales only the
decomposition, never any reported result. Second, the law is derived for \emph{rigid} spheres, and the graph carries that
restriction as a premise: deformability, cell crowding and whole blood are
declared out of scope. The restriction is not a formality. Davis reports that
deformable red cells ``do not always behave as if they have a clear hard sphere
threshold'' \cite[\S5.2.5]{b22}, and observes departures in \emph{both}
directions and with flow speed. Cells fast enough to ``fit'' into the first
streamline stop displacing \cite[\S6.1.1]{b22}, while at lower speeds cells
displace in devices designed to retain them \cite[\S6.1.2]{b22}. He declines to
attribute either to deformability alone, writing of the spherical white cells
that ``the cells and beads also differ in density, electrostatic interactions and
surface chemistry'' \cite[\S6.2.1]{b22}. No signed correction is therefore
available to carry deformable cells inside this law, which is precisely why they
are excluded rather than corrected for.


\clearpage
\section*{C1 --- Mechanism: the cutoff is set by the diverted-lamina width}

\noindent\textbf{Summary.} Five peer-reviewed sources, spanning \emph{three}
independent research groups, assert that the CIF size cutoff is set by the width
of the near-wall lamina diverted into the side channels, not by the gap width. A
standard geometric-exclusion premise (a particle's centre is barred from a
wall strip equal to its radius) converts ``the lamina reaches the centre'' into
the capture criterion $a\lesssim\delta$, so the cutoff diameter is
$\dcrit\approx2\delta$ and the steric prefactor is $\ksteric=2$. The corpus holds
one internal discrepancy, between Dinh's figure caption (``about half'') and its
main text ($a<w$); it is adjudicated to $\ksteric=2$ by the main text's reasoning
and by the Inglis-2006 $D_c=2b$ foundation, with the caption \emph{deselected}
rather than overruled. The node is \textbf{warranted}: the mechanism carries no
residual; only the \emph{value} of the prefactor (settled here at $\ksteric=2$,
its excess deferred to C3d) is argued downstream.

% ----------------------- ARGUMENT GRAPH ---------------------
\begin{figure}[t]
\centering
\begin{tikzpicture}[
  font=\scriptsize,
  src/.style={draw, rounded corners, align=center, inner sep=3pt,
              text width=3.0cm, minimum height=10mm},
  claim/.style={draw, rounded corners, align=center, inner sep=4pt,
                text width=6.6cm, minimum height=10mm},
  ded/.style={draw, double, double distance=0.7pt, rounded corners,
              align=center, inner sep=3pt, text width=4.7cm, minimum height=11mm},
  adj/.style={draw, dashed, rounded corners, align=center, inner sep=3pt,
              text width=4.7cm, minimum height=11mm},
  conv/.style={-{Stealth[length=2mm]}, thick},
  up/.style={-{Stealth[length=2mm]}, thick},
  link/.style={-{Stealth[length=1.8mm]}, dashed},
  nt/.style={font=\scriptsize\itshape, text=black!60, align=center},
]
% --- tier 1: three independent groups ---
\node[src] (g1) at (-5.2,0)
  {\textbf{CIF lineage} (Shevkoplyas)\\ Gifford, Dinh, Lezzar\\ cutoff $\propto$ lamina width};
\node[src] (g2) at (0,0)
  {\textbf{Yamada \& Seki}\\ mechanism $+$ exclusion\\ (not the side-channel size)};
\node[src] (g3) at (5.2,0)
  {\textbf{Inglis \& Herman}\\ mechanism + exclusion\\ (same hedge)};
% --- tier 2: the claim ---
\node[claim] (cl) at (0,-2.5)
  {\textbf{C1.}\ \ $a\lesssim\delta \Rightarrow \dcrit\approx2\delta\ (\ksteric=2)$;\\
   the cutoff is set by the diverted-lamina width, \emph{not} the gap width.
   \ \ \textsf{[warranted]}};
% --- tier 3: deductive strand + CQ5 adjudication ---
\node[ded] (ks) at (-3.1,-5.1)
  {\textbf{Defeasible Modus Ponens.} exclusion ($a$-strip) $+$ capture iff $a<\delta$
   $\Rightarrow$ $a=\delta$ $\Rightarrow$ $\dcrit=2\delta$ $\Rightarrow$ $\ksteric=2$;\\
   \textbf{Established Rule.} Inglis (2006) $D_c=2b$};
\node[adj] (cq5) at (3.1,-5.1)
  {\textbf{CQ5 adjudication.} \underline{selected} main text $a<w$ $+$ Inglis
   $D_c=2b$; \underline{deselected} Fig.~1 ``about half'' $(\ksteric{=}1)$};
% --- edges ---
\draw[conv] (g1) -- (cl);
\draw[conv] (g2) -- (cl);
\draw[conv] (g3) -- (cl);
\draw[up]   (ks) -- (cl);
\draw[link] (cq5) -- (cl);
% --- annotations ---
\node[nt, anchor=south, text width=6.5cm] at (0,0.72)
  {Argument from Expert Opinion: \emph{convergent} over three independent groups};
\end{tikzpicture}
\caption{Argument graph for C1. Three independent groups converge on the
mechanism (Expert Opinion); a short deductive strand fixes $\ksteric=2$ by
Defeasible Modus Ponens from geometric exclusion plus the capture criterion, with
Argument from an Established Rule on Inglis's $D_c=2b$; CQ5 adjudicates
the lone caption-vs-main-text discrepancy. Border style follows the overview map:
single border = warranted claim, doubled border = deductive identity, dashed box
= an open/adjudicated critical question.}
\label{fig:c1-argmap}
\end{figure}

% ----------------------- CLAIM ------------------------------
\field{Claim.} A particle is carried into the side (filtrate) channel iff its
effective radius is at most the diverted-lamina width, $a\lesssim\delta$ (Dinh's
$w$; Inglis's first-streamline width $b$); the cutoff diameter is
$\dcrit\approx2\delta$. The cutoff is governed by the lamina width, not the gap
width.

% ----------------------- SCHEME -----------------------------
\field{Scheme.} Argument from Expert Opinion (WRM~2008,
\S\,Argument from Expert Opinion), \emph{convergent over five peer-reviewed sources
spanning three independent groups}, with a deductive strand beneath it: Defeasible
Modus Ponens on geometric exclusion plus the capture criterion, and Argument from an
Established Rule on Inglis's $D_c=2b$.

% ----------------------- PREMISES ---------------------------
\field{Premises (paper-stated, verbatim).}
\begin{description}[leftmargin=1.6em, style=nextline, itemsep=3pt, topsep=4pt]
  \item[P-C1.1 (Gifford 2014, {\cite[p.~4503]{b1}}).] ``the size cutoff is determined
    primarily by the width of the streamlines diverted into the side (filtrate)
    channels.'' \ $\Rightarrow$ cutoff $\propto$ lamina (streamline) width, not
    gap width.
  \item[P-C1.2$'$ (Dinh 2024, {\cite[p.~917]{b5}}; lower branch).] ``For particles with a
    radius smaller than the width of the extracted flow lamina ($a<w$), we would
    expect their distribution between the filtrate and retentate to approximate
    that of the fluid volume.'' \ $\Rightarrow$ $a<\delta$ $\Rightarrow$
    tracer-like.
  \item[P-C1.2$''$ (Dinh 2024, {\cite[p.~917]{b5}}; upper branch).] ``we would expect
    particles with radii exceeding the width of the extracted flow lamina
    ($a>w$) to remain in the middle channel.'' \ $\Rightarrow$ $a>\delta$
    $\Rightarrow$ retained. With P-C1.2$'$ the threshold sits at $a=\delta$,
    hence $\dcrit=2\delta$.
  \item[P-C1.3 (Yamada \& Seki 2005, {\cite[p.~1234]{b15}}; geometric exclusion).] ``the
    center position of a particle cannot be present within a certain distance
    from the sidewalls, which is equal to the particle radius.'' \ $\Rightarrow$
    the centre is excluded from a wall strip of thickness $a$; capture requires
    the lamina to reach the centre, i.e.\ $a\lesssim\delta$.
  \item[P-C1.3$'$ (Yamada \& Seki 2005, {\cite[p.~1234]{b15}}; mechanism).] ``When
    the relative flow rates distributed into side channels at a branch point are
    sufficiently low \ldots\ only a small portion of the liquid near the sidewalls
    is withdrawn from the main stream. In this flow state, particles whose diameter
    is larger than a specific value would never go through the side channels, even
    if they are flowing near the sidewalls in the main channel, and even if the
    particle size is smaller than the cross sectional size of the side channel.''
    \ $\Rightarrow$ the third independent group asserts the mechanism itself and
    negates the gap explicitly: retention is set by the narrow withdrawn lamina,
    not by the side channel's own cross-section. The condition Yamada states, a
    sufficiently small relative flow rate into each side channel, is the regime CIF
    is built for: its per-gap flow fraction $\fgap$ is small by design, which is
    why a CIF device needs many gaps (P-C1.4). In the source this is the sentence
    that P-C1.3 immediately explains, on the same page.
  \item[P-C1.4 (Lezzar 2022, {\cite[p.~2 (Sci Rep PDF pagination)]{b4}}).] ``It is the magnitude of this flow
    fraction---not the width of the gaps---that determines the size cutoff''
    (and the design reports a 6\um{} critical diameter, requiring ``thousands of
    gaps''). \ $\Rightarrow$ an explicit statement that the flow fraction (hence
    the lamina), not the gap, sets the cutoff.
  \item[P-C1.5 / IH-Q1 (Inglis \& Herman 2013, {\cite[p.~1724]{b18}}; independent).]
    ``Assuming that particles \ldots\ are geometrically restricted away from
    walls by a distance equal to their radius, the cutoff size is determined by
    the width of the boundary layer in the main channel that is drawn up into
    each branch'', carrying the same hedge CIF carries, ``a frequent assumption
    in the microfluidic particle separation literature, that is not always
    true.'' \ $\Rightarrow$ an independent branch-flow group states the C1
    mechanism and the geometric-exclusion premise in one sentence.
\end{description}

% ----------------------- INFERENCE --------------------------
\field{Logical inference.} P-C1.3 (exclusion strip of thickness $a$) together
with P-C1.2$'$/2$''$ (capture iff $a<\delta$) gives the marginal captured
particle $a=\delta$, hence $\dcrit=2a=2\delta$; i.e.\ the steric prefactor
relating cutoff diameter to lamina width is $\ksteric=2$. This is the object
proved as Result~5 of Appendix~A and held in \texttt{steric.KAPPA\_STERIC} in the
reference code; it feeds CQ5 below (where it is defended against the caption) and
C3d (where its excess $\kapp/\ksteric=\residualFolded$ is decomposed).

% ----------------------- CRITICAL QUESTIONS -----------------
\field{Critical questions.}
\begin{description}[leftmargin=1.6em, style=nextline, itemsep=3pt, topsep=4pt]
  \item[CQ1: Expertise (Position-to-Know).] All sources are domain experts on
    this device class: Gifford, Dinh, Lezzar are primary CIF/plasma-skimming
    papers; Inglis \& Herman (2013) and Inglis (2006) are an independent
    branch-flow/DLD group; Yamada \& Seki (2005) originated hydrodynamic
    filtration. \ \emph{Answered: yes.}
  \item[CQ2: Assertion (Evidence).] Each source actually asserts the
    mechanism; verbatim P-C1.1--P-C1.5. \ \emph{Answered: yes.}
  \item[CQ3: Consistency among experts.] No primary source contradicts the
    lamina-width (not gap-width) mechanism: all three groups state it directly
    (Gifford/Lezzar/Dinh, with Lezzar contrastive against the gap width;
    Inglis \& Herman; and Yamada \& Seki, contrastive against the side channel's
    own cross-section, P-C1.3$'$). Yamada \& Seki additionally supply the
    geometric-exclusion sub-premise the deductive strand rests on (P-C1.3).
    \ \emph{Answered: yes.}
  \item[CQ4: Trustworthiness / independence (bias CQ).] The convergence is not
    an artifact of a single laboratory. Three of the five mechanism sources
    (Gifford, Dinh, Lezzar) share the Shevkoplyas group, as does Xia in the wider
    corpus, so within that lineage they \emph{corroborate but do not
    independently establish the mechanism}; it is \emph{independently} asserted by
    two further groups, Yamada \& Seki (2005) and Inglis \& Herman (2013), and the
    geometric-exclusion sub-premise is a standard result in the size-separation
    literature (stated crisply by Yamada \& Seki; the same centre-position
    picture underlies pinched-flow fractionation, cf.~\cite{bmort}). The two
    independent groups, not the lineage, carry the claim. \ \emph{Answered: yes.}
    \emph{(Inglis 2006 and Inglis \& Herman 2013 share a lead author and are one
    group, not two; the three independent groups are the Shevkoplyas lineage,
    Yamada \& Seki, and Inglis.)}
  \item[CQ5: Internal consistency of the steric criterion (RESOLVED $\to
    \ksteric=2$).] The criterion is $a=\delta$ ($\ksteric=2$), not $a=\delta/2$
    ($\ksteric=1$). Dinh states it two ways: the Fig.~1 caption (verbatim)
    ``Particles with an effective radius smaller than about half the width of the
    extracted flow lamina $w$ are carried into the side channels by the filtrate;
    larger particles remain in the middle channel and become progressively more
    concentrated''
    \cite[p.~916, Fig.~1 caption]{b5}, which taken literally gives
    $\ksteric=1$; versus the main text (P-C1.2$'$/2$''$, $a<w$ / $a>w$), which
    gives $\ksteric=2$. The discrepancy is settled in favour of the main text
    because (i) the main text is the substantive discussion and supplies the
    reasoning, whereas the caption is a one-line paraphrase that even hedges
    (``\emph{about} half''); and (ii) the main text matches the theoretical
    foundation CIF is built on: Inglis (2006) defines the critical radius as
    equal to the first-streamline width and hence $D_c=2b$ (C2), and the CIF
    ``extracted lamina'' \emph{is} Inglis's ``first streamline.'' The caption is
    therefore \emph{deselected} (it may describe a stricter ``fully entrained''
    threshold; the marginal particle is never drawn in Dinh's schematic), not
    declared wrong. This \emph{confirms}, rather than moves, the $\ksteric=2$
    baseline: had the caption been adopted, the C3d residual would sit at
    $\kapp/\ksteric=\kappVal$ rather than $\residualFolded$. \ \emph{Answered: $\ksteric=2$.}
\end{description}

% ----------------------- RESIDUAL / STATUS ------------------
\field{Residual risk.} None on the mechanism. The $a<w$ vs $a<w/2$ wording is
resolved to $\ksteric=2$ (CQ5), now with independent third-party support
(Inglis \& Herman, P-C1.5). The \emph{magnitude} of the prefactor in excess of
$\ksteric$ is not a C1 question; it is decomposed at C3d.

\field{Status.} \textbf{Warranted.}


\clearpage
\section*{C2 --- Hydrodynamic conversion: $\delta = \wc\sqrt{\fgap/3}$}

\noindent\textbf{Summary.} Integrating the parabolic gap profile over the
near-wall lamina yields the exact flow-fraction cubic $f = 3x^{2} - 2x^{3}$
($x=\delta/\wc$), whose small-$f$ inverse is $\delta = \wc\sqrt{f/3}$ to within
$0.5\%$ across the Gifford calibration/design band ($\fgap\le5.76\times10^{-4}$;
worst case $\skimErrDesign\%$) and within $0.9\%$ across the wider numerical sweep band
($\fgap = 3.63\times10^{-5}$ to $1.93\times10^{-3}$; worst case $\skimErrBand\%$). The
band is quoted with the figure throughout, because the design-band and
sweep-band worst cases differ by nearly a factor of two
($\skimErrDesign\%$ against $\skimErrBand\%$), and each figure is the small-$f$
inversion error with the parabola itself taken exact. This cubic is \emph{proved} identical to the
Inglis (2006) DLD critical-streamline cubic, unifying the two device classes.
The same near-wall-lamina construction is \emph{independently applied}
(numerically, with no closed form) by Inglis~\&~Herman (2013) over the same
rectangular-duct field, and the cubic's deep-channel form is recovered by an
exact rectangular-duct calculation and by single-gap CFD. The only correction at
real, finite aspect ratios is the bounded depth factor $\Kdepth = R(\alpha)$,
CFD-measured to be $[0.89,0.98]$ over the probed range $\alpha=0.8$--$2.69$ and
extended by the analytic duct law to a $0.81$ floor at $\alpha = 80/150 = 0.5333$; it is
folded into the apparent prefactor rather than applied to the data (the cutoff
data collapse more tightly on the clean parabolic $\delta$). The node is
\textbf{warranted on the form}: the identity is the load-bearing claim, the
cross-geometry transfer is the only defeasible step, and the single-gap CFD
closes it ($\kinplane=1$).

% ----------------------- ARGUMENT GRAPH ---------------------
\begin{figure}[t]
\centering
\begin{tikzpicture}[
  font=\scriptsize,
  ded/.style={draw, double, double distance=0.7pt, rounded corners,
              align=center, inner sep=3pt, text width=3.55cm, minimum height=12mm},
  src/.style={draw, rounded corners, align=center, inner sep=3pt,
              text width=3.55cm, minimum height=12mm},
  claim/.style={draw, rounded corners, align=center, inner sep=4pt,
                text width=7.2cm, minimum height=11mm},
  adj/.style={draw, dashed, rounded corners, align=center, inner sep=3pt,
              text width=4.9cm, minimum height=12mm},
  conv/.style={-{Stealth[length=2mm]}, thick},
  up/.style={-{Stealth[length=2mm]}, thick},
  link/.style={-{Stealth[length=1.8mm]}, dashed},
  nt/.style={font=\scriptsize\itshape, text=black!60, align=center},
]
% --- tier 1: deductive core + corroborators ---
\node[ded] (core) at (-5.1,0)
  {\textbf{Deductive core.} closed-form cubic $f=3x^{2}-2x^{3}$ (P-C2.1)
   \emph{proved} $\equiv$ Inglis DLD cubic (P-C2.4)};
\node[src] (num) at (0,0)
  {exact duct field $+$ single-gap CFD (P-C2.5): $\dskim\!\propto\!\sqrt{f}$,
   $R^{2}\!\ge\!0.9999$; deep-limit $0.99913$};
\node[src] (ih) at (5.1,0)
  {\textbf{Inglis \& Herman} (independent): same lamina construction,
   numerical (P-C2.6)};
% --- tier 2: claim ---
\node[claim] (cl) at (0,-2.6)
  {\textbf{C2.}\ \ $\delta = \wc\sqrt{\fgap/3}$: leading term of the exact
   cubic (deep-channel, small-$f$).\ \ \textsf{[warranted on the form]}};
% --- tier 3: analogy CQ (settled) + depth caveat ---
\node[adj] (an) at (-3.1,-5.2)
  {\textbf{Analogy (cross-geometry).} one cubic governs the CIF aperture \&
   the DLD post-gap. \underline{CQ}: aperture funnelling? \underline{settled}:
   CFD $\kinplane=1$ (line-sink)};
\node[adj] (dep) at (3.1,-5.2)
  {\textbf{Depth caveat.} $\Kdepth=R(\alpha)\in[0.81,0.98]$ folded into
   $\kapp$ (C3d), \emph{not} applied to data; collapse tighter on
   parabolic $\delta$};
% --- edges ---
\draw[conv] (core) -- (cl);
\draw[conv] (num)  -- (cl);
\draw[conv] (ih)   -- (cl);
\draw[up]   (an)   -- (cl);
\draw[link] (dep)  -- (cl);
% --- annotation ---
\node[nt, anchor=south, text width=7cm] at (0,0.85)
  {Proven identity (load-bearing) $+$ Analogy (the only
   defeasible step) $+$ independent Expert Opinion};
\end{tikzpicture}
\caption{Argument graph for C2. The deductive core, the closed-form cubic and
its \emph{proved} identity to the Inglis DLD cubic (doubled border), is the
load-bearing claim; the exact-duct/CFD recovery and the independent
Inglis~\&~Herman construction corroborate it. The single defeasible step is the
\emph{cross-geometry} transfer of one cubic between a wall-aperture and a
post-gap; its critical question (aperture funnelling) is settled by the
single-gap CFD ($\kinplane=1$). The depth factor is a bounded caveat folded into
$\kapp$, not a defeater of the form. Border style follows the overview map:
doubled = deductive identity, single = warranted/corroborating, dashed =
critical question or caveat.}
\label{fig:c2-argmap}
\end{figure}

% ----------------------- CLAIM ------------------------------
\field{Claim.} In the deep-channel, small-$f$ limit the diverted lamina width is
$\delta = \wc\sqrt{\fgap/3}$, the leading term of the exact cubic
$f = 3x^{2} - 2x^{3}$ ($x = \delta/\wc$).

% ----------------------- SCHEME -----------------------------
\field{Scheme.} A deductive core, \emph{a proven, code-checked identity}
(defeasible modus ponens in the umbrella sense; no critical questions of its
own), comprising (i) the closed-form cubic and its inversion (P-C2.1)
and (ii) the \emph{proved} algebraic identity of the CIF and DLD cubics
(P-C2.4). This is reinforced by an \emph{Argument from Analogy} for the physical
transfer of the \emph{single} cubic across two different geometries (a
wall-aperture vs.\ the space between two posts). That is the step the
Analogy-CQs interrogate and the single-gap CFD settles ($\kinplane=1$). It is
reinforced further by an independent \emph{Argument from Expert Opinion} (Inglis~\&~Herman's
identical near-wall-lamina construction). The identity is therefore the
load-bearing claim; the analogy bears only on whether one equation governs both
geometries, not on the algebra itself.

% ----------------------- NOTATION ---------------------------
\field{Notation.} Inglis (2006) prints the first-streamline width as Greek
$\beta$ throughout: in the defining sentence (``which we define as~$\beta$''),
in eqn~(1) ($D_c = 2\beta$), in the cubic eqn~(5), and in the Fig.~1 caption and
in-figure label; Latin~``$b$'' does \emph{not} appear as that symbol anywhere in
the paper (confirmed against the typeset page, not the OCR, which silently
flattens $\beta\!\to\!b$). Per the convention fixed in the overview's Notation
section, this manuscript renders the symbol as Latin~$b$; quoted Inglis phrases
keep his glyph, and our restatements use~$b$. The identity $D_c = 2b$ (Inglis's
$D_c = 2\beta$) gives $\ksteric = 2$, independent of the glyph.

% ----------------------- PREMISES ---------------------------
\field{Premises.}
\begin{description}

\item[P-C2.1 (exact cubic: symbolic, code-checked).]
Integrating the width-wise parabolic velocity profile $u \propto s(\wc - s)$
across the gap gives the cumulative near-wall flow fraction
$f(x) = 3x^{2} - 2x^{3}$ exactly ($x = \delta/\wc$), with total flux
$Q_{\mathrm{tot}} \propto \wc^{3}/6$. Inverting in the small-$f$ limit,
\[
  \delta \;=\; \wc\sqrt{f/3} \;+\; \tfrac{1}{9}\,\wc f \;+\;\cdots,
\]
with coefficients $a_1 = 1/\sqrt{3} = 0.577350$ and $a_2 = 1/9 = 0.111111$ (the
program asserts both exactly; Appendix~A also lists $a_3 = 0.053458$). The
leading term is \emph{exact as $f\!\to\!0$}; its relative error grows as
$\tfrac{1}{3\sqrt{3}}\sqrt{f} \approx 0.192\sqrt{f}$, reaching $\approx 0.2\%$ at
$f = 1.04\times10^{-4}$, a worst case of $\skimErrDesign\%$ across the Gifford
calibration/design band ($\fgap\le5.76\times10^{-4}$; Gifford's own per-gap design
bound is $\fgap < 5\times10^{-4}$ \cite[p.~4498]{b1}), and a worst case of
$\skimErrBand\%$ across the wider numerical sweep band, which reaches
$\fgap = 1.93\times10^{-3}$.
\prov{code: skim.cubic, skim.delta\_leading, skim.a1/a2,
skim.verify\_symbolic/\_numeric; Appendix~A Results~1--2}

\item[P-C2.2 (Xia 2016, verbatim).]
The centre-channel cross-sectional velocity profile is ``approximately
parabolic'' \cite[art.~35943, p.~10 (Sci Rep PDF pagination)]{b2}, which licenses the parabolic premise
of P-C2.1. On the same page Xia add that the diverted streamline width is
``roughly proportional to the effective c.d.,'' the proportionality C3a relies
on.

\item[P-C2.3 (Inglis 2006: the DLD construction, faithfully rendered).]
The critical particle radius ``is equal to the width of the first streamline''
\cite[p.~656]{b17}; Inglis denotes that width $\beta$ (rendered $b$ here; see
\emph{Notation}) and writes $D_c = 2\beta$ [eqn~(1)], i.e.\ $D_c = 2b$. The
streamline-width construction is, in our notation,
\[
  \int_0^{b} u(x)\,\mathrm{d}x \;=\; \varepsilon \int_0^{g} u(x)\,\mathrm{d}x,
  \qquad u(x) = \frac{g^2}{4} - \Bigl(x - \frac{g}{2}\Bigr)^{\!2},
\]
which reduces to the cubic
$\bigl(b/g\bigr)^{3} - \tfrac{3}{2}\bigl(b/g\bigr)^{2} + \varepsilon/2 = 0$
\cite[eqns~(3)--(5)]{b17}. Inglis obtains this ``by assuming a conventional
parabolic flow profile through the gap, with zero velocity at the post
sidewalls'' \cite[p.~656]{b17}.

\item[P-C2.4 (identity: symbolic, code-checked).]
Computing the Inglis integral $\int_0^{b} u = \varepsilon\!\int_0^{g} u$ (with
$u = g^2/4 - (x-g/2)^2$) yields $\varepsilon = 3y^2 - 2y^3$ ($y = b/g$),
\emph{algebraically identical} to the CIF cubic $f = 3x^2 - 2x^3$.
\prov{code: steric.eps\_of\_y, asserted equal to skim.cubic in
steric.verify\_symbolic; Appendix~A Results~3--4}
Hence the CIF skim cubic and the DLD critical-streamline cubic are the same
equation: a CIF gap and a DLD gap obey one streamline-width law, and the leading
inversion gives $b = g\sqrt{\varepsilon/3}$, identical in form to
$\delta = \wc\sqrt{f/3}$. This cross-domain identity, which the prior DLD and
CIF literatures do not state, is the structural keystone that lets results
proven for one device transfer to the other.

\item[P-C2.5 (exact duct field: numerical, code-checked).]
The exact rectangular-duct (depth-integrated) velocity field
\cite[eqn~(15)]{bbrody} recovers $\delta$ to
$0.09\%$ in the deep limit
($\delta_{\mathrm{general}}/\delta_{\mathrm{parab}} \to 0.99913$ at $\alpha=50$,
confirmed by re-running the routine), and single-gap CFD confirms
$\dskim \propto \sqrt{f_{\mathrm{phys}}}$ to $R^{2} \ge 0.9999$ with $\dskim$ on
the exact duct law to $\le 0.6\%$ (the validator's guard tolerance is $0.8\%$;
the \emph{achieved} agreement is $\le 0.6\%$, max $0.59\%$ on Run~B for $\dskim$
itself). The
CFD-measured depth factors are $R(0.80/1.40/2.69) =
\KdepthCfdShallow/\KdepthCfdMid/\KdepthCfdDeep$ against the analytic duct law's
$\KdepthAnalyticShallow/\KdepthAnalyticMid/\KdepthAnalyticDeep$; the analytic
triple is $\sqrt{3/c_2(\alpha)}$ evaluated by the kernel's \texttt{K\_depth}.
\prov{code: f\_general, duct.K\_depth, verify\_cfd; Methods \S4.3}

\item[P-C2.6 / IH-Q2 (Inglis \& Herman 2013, verbatim: independent use of the
construction).]
The cutoff flux is ``calculated by integrating the analytical solution to the
Navier--Stokes equation \dots\ from $W/2 - R_c$ to $W/2$ \dots\ on the
depth-averaged velocity, $\tilde{u}(x)$, in a rectangular channel
($d\times w$)'' \cite[p.~1726]{b18}. An independent group thus
integrates the same near-wall lamina over the same rectangular-duct field to
relate flow fraction to cutoff (numerically, with no closed form). They
attribute that duct field to the \emph{same} source this node cites,
\cite{bbrody}: Brody's eqn~(15) is their ref.~17, so the duct-field provenance
is shared, not merely parallel.

\end{description}

% ----------------------- DISTINCTION ------------------------
\field{Distinction from Inglis \& Herman (sets up the novelty statement).}
Inglis~\&~Herman integrate the analytical depth-averaged duct field for the
cutoff flux and solve the device numerically (COMSOL), with no closed-form
cutoff law, and report that analytical resistor-network
approaches ``have so far failed us'' \cite[p.~1725]{b18}. CIF's new
contribution is the closed-form parabolic-limit cubic (P-C2.1), its exact
identity to the Inglis-2006 DLD cubic (P-C2.4), the cross-domain unification the
prior literatures do not state, and (separately, C4) the closed-form
$F$-identity that \emph{succeeds} where their network route failed.

% ----------------------- EMPIRICAL SUPPORT ------------------
\field{Empirical support for the $\sqrt{f}$ law.}
Inglis's own illustrative values confirm the $\sqrt{f}$ scaling this node
asserts: ``For $\varepsilon = 0.1$ the predicted critical particle diameter is
$0.4g$ \dots\ For $\varepsilon = 0.01$ \dots\ $0.12g$'' \cite[p.~656]{b17}.
These match the \emph{exact} cubic root ($\dcrit/g = 0.392$ at $\varepsilon=0.1$,
$0.118$ at $0.01$, which Inglis rounds to $0.4$ and $0.12$); the leading law
$2\sqrt{\varepsilon/3} = 0.365$ and $0.115$ captures the $\sqrt{\varepsilon}$
scaling but runs a few percent low at these DLD-scale shift fractions ($6.8\%$
at $\varepsilon=0.1$, $2.0\%$ at $0.01$); the $+f/9$ term of P-C2.1 removes most of that
(leaving $1.1\%$ and $0.10\%$), and the gap shrinks to sub-percent across the much
smaller CIF $f$-range. Inglis attributes
the ``surprisingly large critical size at small $\varepsilon$'' to the slow
near-wall fluid forcing a wider first streamline \cite[p.~656]{b17}: the
parabolic-vs-plug enhancement that distinguishes the $\sqrt{f}$ regime from a
flat profile. Davis's empirical $\dcrit = 1.4\,G\varepsilon^{0.48}$ \cite[eqn~(2-22)]{b22} (his $G$ is the DLD gap, Inglis's $g$; he quotes $0.46\,G$ at $\varepsilon = 0.1$) is the same
square-root-like law fit to data ($\varepsilon^{0.48}\!\approx\!\sqrt{\varepsilon}$).
Because the prefactor these imply is constant across $\varepsilon =
0.01$--$0.5$ (no regime change), this evidence supports the $\sqrt{f}$ law and,
secondarily, the $\ksteric = 2$ floor, and \emph{not} a knee (C3c).

% ----------------------- CRITICAL QUESTIONS -----------------
\field{Critical questions.}
The cubic identity itself (P-C2.1, P-C2.4) is deductive and carries no critical
questions; the questions below test the single \emph{defeasible} step, the
cross-geometry transfer of the cubic, and the bounded depth caveat, not the
algebra.
\begin{description}

\item[Verification-CQ: does the cubic hold against data?]
The cubic is a proven identity on a stated premise: P-C2.1 (symbolic proof)
$+$ P-C2.4 (Inglis-identical) $+$ P-C2.5 (exact-duct $+$ CFD recovery). The
construction was validated two ways by Davis. His Femlab finite-element
simulations find the parabolic profile ``a reasonable estimation for the
velocity profile between two walls'' \cite[\S2.3.2]{b22}, and his measured
critical diameters over ``about 20 different devices'' ($G = 1.3$--$38\um$,
$\varepsilon = 0.005$--$0.5$, beads $0.9$--$22\um$) \cite[\S2.4]{b22} fit the
square-root-like law. It was validated experimentally by Inglis (particles
$2.3$--$22\um$ \cite[p.~655]{b17}), and it is now reproduced by the CIF
single-gap CFD in full 3-D geometry. Scheme: Argument from Evidence to
Hypothesis, verification form (the graph's \texttt{c2\_davis\_fem} and
\texttt{c2\_cfd}). $\rightarrow$ \emph{answered: yes.}

\item[Bounds: how large are the deviations?]
Deviations are bounded. The leading-$\delta$ form is \emph{exact as
$f\!\to\!0$} and its relative error \emph{grows} with $f$ (it is $\approx
0.192\sqrt{f}$, hence largest at the top of the range), but stays
bounded: $\approx 0.2\%$ at $f = 1.04\times10^{-4}$, worst case $\skimErrDesign\%$ across
the design band and $\skimErrBand\%$ across the numerical sweep band; the only other
deviation is from non-parabolic (finite-depth)
profiles, bounded by $\Kdepth = R(\alpha) \in [0.81, 0.98]$ (CFD-anchored at the upper end, analytic at the lower).
Scheme: none. This is the deductive core's own error budget, not a critical
question; the core carries none, as stated above.
$\rightarrow$ \emph{answered: bounded.}

\item[Analogy-CQ1: relevant differences (DLD $\leftrightarrow$ CIF).]
The single cubic transfers across the two geometries despite the difference: a
CIF gap is a finite aperture in a channel wall, not the space between two posts,
but the single-gap CFD shows this aperture is a line-sink with no funnelling
($\kinplane = 1$), so the cubic transfers intact and there is no
aperture-funnelling enhancement of the prefactor, which is what settles the
funnelling question at C3d. Scheme: Argument from Analogy, difference CQ.
$\rightarrow$ \emph{answered: the form transfers; no funnelling enhancement.}

\item[Analogy-CQ2: base-case truth.]
The DLD base case is true, triply: Inglis experimental data $+$ Davis FEA $+$
CIF CFD all confirm the parabolic streamline-width cubic. Scheme: Analogy,
base-truth CQ. $\rightarrow$ \emph{answered: yes.}

\item[Analogy-CQ3: contrary instances.]
The analogy excludes the right cases: plug / electroosmotic flow is correctly
excluded because CIF is pressure-driven with no-slip walls. Davis quantifies the
contrast through a \emph{variable} enhancement parameter $\zeta$, ``a function of
$\varepsilon$ and post geometry'' \cite[eqn~(2-19)]{b22}, obtained as the
cube-root of his eqn~(2-21): at $\varepsilon = 0.1$ the cubic has a single real
root $\zeta = 1.96$, giving $\dcrit = 3.92\,G\varepsilon \approx 0.39\,G$,
``about twice as wide as'' the plug-flow baseline $\dcrit = 2\,G\varepsilon$ for
``even width stream lines'' \cite[p.~22]{b22}. Crucially $\zeta$ is
$\varepsilon$-dependent (it $\to 1$ in the plug limit), which is \emph{why} the
parabolic law is $\sqrt{\varepsilon}$-like rather than linear, consistent with
this node's $\delta \propto \sqrt{f}$ and with Davis's own
$\varepsilon^{0.48}$ fit. Note that $\zeta$ is Davis's
\emph{streamline-width enhancement}, \emph{not} $\ksteric$: the steric factor
$\ksteric = 2$ is pure $D = 2\delta$ geometry and is profile-\emph{independent};
what the profile changes is the $f \leftrightarrow \delta$ map (parabolic
$\delta \propto \sqrt{f}$ vs.\ plug $\delta \propto f$). Scheme: Analogy,
contrary-case CQ. $\rightarrow$ \emph{answered: yes.}

\item[Expert-CQ (Inglis \& Herman: the independent strand).]
The corroborating expert source is credible, independent, and consistent:
Inglis~\&~Herman 2013 is peer-reviewed (\emph{Lab on a Chip}), an independent
(non-Shevkoplyas) branch-flow group, and applies the \emph{identical}
near-wall-lamina construction over the same duct field (P-C2.6); their result
(no closed form) does not conflict with CIF's: it is the numerical counterpart
of the same integral. Scheme: Expert Opinion (expertise, independence,
consistency CQs). $\rightarrow$ \emph{answered: yes.}

\end{description}

% ----------------------- RESIDUAL / STATUS ------------------
\field{Residual risk.} None on the form (exact, Inglis-identical, CFD-confirmed).
One caveat: ``deep parabolic'' is a slight stretch at real aspect ratios. CIF
runs at $\alpha = 0.5$--$2.7$, below Davis's stated validity for ignoring depth:
``\emph{a ratio of depth $E$ to gap size $G$ (aspect ratio) of at least 5}
\dots\ \emph{the depth dimension can be ignored, resulting in a two dimensional
calculation}'' \cite[\S2.3.2]{b22}. Both ratios are depth/(parabola-span), so the
comparison is like-for-like. That is precisely
the small $\Kdepth = R(\alpha)$ correction the duct law gives and the CFD
confirms ($[0.89, 0.98]$ over $\alpha = 0.8$--$2.69$; analytic floor $0.81$ at
$\alpha = 80/150 = 0.5333$). The cutoff still empirically prefers the parabolic $\delta$
(C3b, Methods \S4.2): the collapse is tighter on the clean $\delta$, so
$R(\alpha)$ is folded into the apparent $\kapp$ rather than applied to the data.

\field{Status.} \textbf{Warranted on the form.}


\clearpage
\section*{C3 --- The operating law: $\dcrit=\kapp(\dwc,\,\text{depth})\,\delta$}

\noindent\textbf{Summary.} The operating law converts the diverted-lamina width
$\delta$ (fixed by C2 as $\delta=\wc\sqrt{\fgap/3}$) into a particle-diameter
cutoff through a single prefactor, $\dcrit=\kapp\,\delta$. The central empirical
finding is that this prefactor is \emph{not} the universal constant the early
literature assumed. On the large-particle plateau ($\dwc\gtrsim0.035$) it is
tightly determined at $\kapp=\plateauKappa$ on the CIF calibration data, and it carries the
$\tfrac12$ and $1$ scaling exponents the law predicts (C3b; exponents derived in
Appendix~A, Result~6). The cross-paper comparison that once stood beside those
numbers is now reported as a \emph{lead}: it is same-figure, and its agreement is
numerically the depth correction (C3b, CQ2). Below a knee at
$\dwc\approx0.03$ it ramps down to a \emph{sub-steric (recovery-suppressed)} foot of ${\approx}\footKappa$,
a size dependence that a single power law cannot fit and that, correctly
normalised, deterministic lateral displacement (DLD) theory does not share,
marking it CIF-specific and finite-particle in origin (C3c). Decomposing the
plateau value settles three of its five factors: the steric factor
($\ksteric=2$, Appendix~A Result~5), in-plane funnelling ($\kinplane=1$, by
single-gap CFD), and a depth factor ($\Kdepth=R(\alpha)$); it isolates a
single lumped residual $\sqrt{c}\,C$ that the published data cannot split and
that is provably \emph{not} recoverable from the device-level filtrate fraction
$F$ (C3d).

% ----------------------- ARGUMENT GRAPH 1 (structural) ------
\begin{figure}[t]
\centering
\begin{tikzpicture}[
  font=\scriptsize,
  claim/.style={draw, rounded corners, align=center, inner sep=4pt,
                text width=10.4cm, minimum height=11mm},
  sub/.style={draw, rounded corners, align=center, inner sep=3pt,
              text width=3.05cm, minimum height=15mm},
  part/.style={draw, dashed, rounded corners, align=center, inner sep=3pt,
               text width=3.05cm, minimum height=15mm},
  up/.style={-{Stealth[length=2mm]}, thick},
  nt/.style={font=\scriptsize\itshape, text=black!60, align=center},
]
\node[claim] (cl) at (0,0)
  {\textbf{C3.}\ \ $\dcrit=\kapp(\dwc,\text{depth})\,\delta$, with
   $\delta=\wc\sqrt{\fgap/3}$ (C2). \ $\kapp$ is \emph{not} a universal constant:
   $\plateauKappa\,[\kappCIlo,\kappCIhi]$ on the plateau, ramping to ${\approx}\footKappa$ below a knee
   at $\dwc\!\approx\!0.03$.};
\node[sub] (a) at (-5.55,-2.7)
  {\textbf{C3a}: form\\ $\dcrit\propto\delta$\\ ($\kappa\!\equiv\!\dcrit/\delta$;
   4 sources)\\ \textsf{[warranted]}};
\node[part] (b) at (-1.85,-2.7)
  {\textbf{C3b}: value\\ plateau $\kapp=\plateauKappa$,\\ $\beta=\tfrac12$,
   $\gamma=1$;\\ cross-paper anchor\\ a \emph{lead}, not evidence\\
   \textsf{[partially closed]}};
\node[part] (c) at (1.85,-2.7)
  {\textbf{C3c}: knee\\ $\kapp$ varies; global\\ law rejected
   ($\chired=30.6$)\\ \textsf{[warranted: existence;}\\ \textsf{rival open]}};
\node[part] (d) at (5.55,-2.7)
  {\textbf{C3d}: magnitude\\ $\kapp=2\!\cdot\!1\!\cdot\!R\!\cdot\!\sqrt{c}\,C$;\\
   3 settled, $\sqrt{c}\,C$ open\\ \textsf{[partially closed]}};
\draw[up] (a) -- (cl);
\draw[up] (b) -- (cl);
\draw[up] (c) -- (cl);
\draw[up] (d) -- (cl);
\node[nt, anchor=north, text width=11.5cm] at (0,-4.05)
  {C3a fixes the \emph{form}, C3b the plateau \emph{value} and exponents, C3c
   shows the value \emph{varies} (the knee), C3d \emph{interprets} the value.
   Single border = warranted; dashed = at least one item still open.};
\end{tikzpicture}
\caption{Structure of node C3. The operating law rests on four sub-arguments:
proportionality (C3a), the plateau value and exponents (C3b), the reality of the
knee (C3c), and the decomposition of the prefactor's magnitude (C3d). Only C3a is
closed. The other three are drawn dashed, and for three different reasons: C3b's
cross-paper anchor is circular \emph{and} non-diagnostic, so it is reported as a
lead; C3c's \emph{existence} claim stands while its \emph{mechanism} has a
live rival; and C3d's lumped residual $\sqrt{c}\,C$ is a named open quantity.
None of the three touches the plateau fit itself, which is settled on the CIF
calibration data alone.}
\label{fig:c3-structure}
\end{figure}

\field{Claim.} $\dcrit=\kapp\,\delta$, where $\kapp$ is not a universal
constant: $\kapp=\plateauKappa\ [\kappCIlo,\,\kappCIhi]$ on the plateau ($\dwc\gtrsim0.035$),
ramping to ${\approx}\footKappa$ below a knee at $\dwc\approx0.03$. (This replaces the
refuted ``$\kappa\approx5$ constant, globally $\sqrt{f}$'' claim.) The lamina
width $\delta=\wc\sqrt{\fgap/3}$ is fixed by C2; $\ksteric=2$ and the $\tfrac12/1$
exponents are derived in Appendix~A (Results~5,~6). Four sub-arguments follow,
each a defeasible argument in the sense of WRM~2008.

% ===========================================================================
\subsection*{C3a --- Proportionality ($\dcrit\propto$ lamina width)}

\field{Summary.} Four peer-reviewed sources state the cutoff scales linearly
with the diverted-lamina width, so the operating law can be written
$\dcrit=\kappa\,\delta$ with $\kappa$ a (possibly size-dependent) proportionality
constant; the dispute is confined to $\kappa$'s value.

\field{Claim.} The cutoff diameter scales linearly with the lamina width;
$\kappa\equiv\dcrit/\delta$.

\field{Scheme.} Argument from Expert Opinion (WRM~2008), convergent over four
primary sources, $+$ Argument from Example (the measured exponents select the
form). The expertise, assertion, and
backup-evidence critical questions for the expert-opinion scheme are discharged
at C1 (the same convergent peer-reviewed base) and are inherited here; the live
critical question at C3a is consistency.

\field{Premises.}
\begin{description}
\item[P-C3a.1 (Xia 2016, verbatim).]
``the width of the streamline that is diverted into a side channel at each
(filtration) gap is roughly proportional to the effective c.d.\ of the device at
that gap'' \cite[Methods, p.~10 (Sci Rep PDF pagination), assumption~(b)]{b2}. Xia defines ``c.d.'' as
the \emph{critical diameter}: ``particles above a certain size (or `critical
diameter', c.d.)'' \cite[p.~2 (Sci Rep PDF pagination)]{b2}. So
assumption~(b) states $\delta\propto\dcrit$,
equivalently $\dcrit\propto\delta$, the claim.
\item[P-C3a.2 (Inglis 2006).]
$\dcrit=2b$ \cite[eqn~(1)]{b17}: the critical diameter is twice the
first-streamline width $b$, hence linear in it. (Inglis prints this width as
Greek~$\beta$; we render it $b$ throughout; see the C2 Notation paragraph. This
premise is P-C2.3; with the CIF extracted lamina identified with Inglis's first
streamline, $b\leftrightarrow\delta$.)
\item[P-C3a.3 (Dinh 2024).]
The C1 criterion, particles with $a<w$ ``approximate \dots\ the fluid volume''
while those with $a>w$ ``remain in the middle channel'', is a linear-in-$w$
threshold \cite[p.~917]{b5} (P-C1.2$'$/2$''$). (Dinh's $w$ is our $\delta$.)
\item[P-C3a.4 / IH-Q2$'$ (Inglis \& Herman 2013, verbatim).]
The design assumption ``a boundary stream width of $R_{c}$ \dots\ gives a cutoff
size of $R_{c}$'' \cite[p.~1726]{b18}: cutoff radius linear in
lamina width. This is the design-assumption sentence; it is distinct from the
adjacent integration sentence used at C2 (registry tag IH-Q2) and carries its
own registry number, IH-Q2$'$.
\end{description}

\field{Critical question.}
\begin{description}
\item[CQ: Consistency.] The four sources agree the cutoff scales with the
lamina width (premises P-C3a.1--4); all four state a linear cutoff--lamina-width
relation and no primary source contradicts it. Only the value of the
proportionality constant is disputed ($\rightarrow$ C3b--d). Scheme: Expert
Opinion, consistency CQ.
$\rightarrow$ \emph{answered: yes.}
\end{description}

\field{Status.} \textbf{Warranted.}

% ===========================================================================
\subsection*{C3b --- Plateau value and exponents}

\field{Summary.} A conservative weighted fit to three-width calibration data
places the plateau prefactor at $\kapp=\plateauKappa$ (CV~$\kappCV\%$) with the predicted
$\tfrac12$ and $1$ exponents inside their confidence intervals, and the result is
invariant to digitization noise and to three independent digitizations (the primary and two re-digitizations). This
node also carries a \emph{downgrade}. The cross-paper comparison with Xia's
devices, long reported here as a held-out corroboration at $\kapp = 5.17$, is a
real recalculation from quoted data but it is not held out and it is not
diagnostic: Xia takes its cutoff from the same Gifford figure this plateau was
fitted from, and (the stronger reason) the $5.2\%$ by which the recomputed mean
exceeds the plateau \emph{is} the $140/80\um$ depth correction. It is reported as
a \emph{lead}. The fit stands on the CIF calibration data alone, and it does so
without help.

\field{Claim.} On the large-particle plateau, $\kapp=\plateauKappa\ [\kappCIlo,\,\kappCIhi]$
(CV~$\kappCV\%$), with cutoff exponents $\beta=\expBeta\ [\betaCIlo,\,\betaCIhi]$ (model $\tfrac12$)
and $\gamma=\expGamma\ [\gammaCIlo,\,\gammaCIhi]$ (model $1$). The model exponents $\tfrac12$ and
$1$ are derived in C6 / Appendix~A (Result~6). The cross-paper figure
$\kapp\approx\xiaKappaMean$ (Xia) is reported below as a \emph{lead} and is not part of
this claim (CQ2).

\field{Scheme.} Argument from Example. (The weighted least-squares fit is
premise-generation, a deterministic computation, not a defeasible scheme; the
generalisation from the calibration sample is the Argument from Example.)

\field{Premises (derived).}
\begin{description}
\item[P-C3b.1 (weighted fit).]
A weighted least-squares fit (error-bar-weighted, with confidence intervals
scaled by $\sqrt{\chired}=\sqrt{\chiPlateau}\approx1.95$, hence conservative) on the
verified Gifford Fig.~3 data (15 points across three width curves; plateau
subset $\dwc\ge0.035$, $n=\plateauN$) gives $\beta=\expBeta\ [\betaCIlo,\,\betaCIhi]$,
$\gamma=\expGamma\ [\gammaCIlo,\,\gammaCIhi]$, plateau $\kapp=\plateauKappa\ [\kappCIlo,\,\kappCIhi]$, CV~$\kappCV\%$. A
threshold-sensitivity sweep gives $\beta\in[0.51,\,0.59]$ across cuts
$0.030$--$0.045$, so the result is not an artifact of the plateau definition.
\prov{code: cif\_cutoff.py B3/B4}
\item[P-C3b.2 (cross-paper consistency check: \emph{not} a held-out test).]
\emph{Recalculation} from three Xia devices, using their stated cutoff and
$\fgap$ values, gives $\kapp = \xiaKappaA/\xiaKappaB/\xiaKappaC$, mean $\xiaKappaMean$, at $140\um$ depth
against the Gifford plateau $\kappVal$ fitted at an assumed $80\um$: an
agreement of $\xiaAgreePct\%$ across papers and depths. All three device triples are
quoted data: OR-7 at $\fgap = 5.8\times10^{-4}$ for a $7\um$ cutoff, and RS-8 and
RS-9 at $7.2\times10^{-4}$ and $8.6\times10^{-4}$ for $8$ and $9\um$, all at
$\wcref = 100\um$, are each stated verbatim in the source \cite[pp.~6 and~10 (Sci Rep PDF pagination)]{b2}. \textbf{The recalculation is real and rests on quoted data.} What it is
\emph{not} is a held-out prediction, for two reasons that CQ2 separates, and
the appendix states them rather than leaving them to be found. \prov{code: B5;
glossary \xref{app:units_selval}{A.7}}
\item[P-C3b.3 (digitization robustness).]
The plateau $\beta/\gamma$ survive a 2000-draw digitization Monte Carlo ($\beta$
median $0.51$, $P(0.4<\beta<0.7)=1.00$) and are invariant across three
independent digitizations of Fig.~3 (the primary and two re-digitizations) ($\kapp\in[4.91,\,4.92]$, $<1\%$ spread;
$\beta\in[0.50,\,0.51]$; $\gamma\in[1.01,\,1.02]$; the digitizations agree to a
median of $0.4\%$ and a worst-point $2.1\%$ on $f^{*}$). \prov{code: B2b/B2c}
\end{description}

\field{Critical questions.}
\begin{description}
\item[CQ1: Sample / computation adequacy.] The plateau $\kappa$ and exponents
are well-determined and not digitization artifacts: CV~$\kappCV\%$ across three
widths; conservative ($\chi^{2}$-inflated) confidence intervals; a threshold
sweep that keeps $\beta\in[0.51,0.59]$ across cuts $0.030$--$0.045$; survival of
the Monte Carlo and of three independent digitizations; and all numbers
reproducing byte-for-byte on execution. None of this depends on the Xia anchor.
Scheme: Argument from Example. $\rightarrow$ \emph{answered: yes.}

\item[CQ2: Is the anchor independent, and is the agreement diagnostic?
\textup{(\texttt{xia\_anchor\_unverified})}]
Neither, and this is the node's one open question. There are two separate
reasons, \textbf{and the second is the stronger one}. It is worth saying at the
outset what is \emph{not} in question: the recalculation is real and its inputs
are quoted, so nothing here impugns the arithmetic or the data.

\emph{First, it is not independent.} Xia reads its cutoff from the same Gifford
figure the plateau was fitted from, and Gifford is a Xia co-author, so the anchor
and the fit share a \emph{figure} rather than merely a literature. Xia's own
sentence supplies the chain: ``\emph{OR-7 was patterned using $\wc = 100\um$ and
$\fgap = 5.8\times10^{-4}$, a combination of parameters we found corresponding to
a c.d.\ of 7}''~$\mu$m ``\emph{previously}'', citing Gifford~2014 and pointing at
the specific figure: ``\emph{e.g.\ see Fig.~3 of Gifford et al.}'' That is the
figure this plateau is fitted from. Reproducibility is not independence.

\emph{Second, and decisively: the agreement is numerically the depth correction
itself.} Xia's devices are $140\um$ deep against the plateau's assumed $80\um$,
and $1/\Kdepth(1.40) - 1 = \depthCorrPct\%$, the same correction as the $\xiaAgreePct\%$
agreement above, the two printed figures differing only because they fall either
side of a rounding boundary at full precision. Equivalently,
$\kappVal/\KdepthAnalyticMid = \xiaKappaMean$ reproduces the recomputed mean to three significant
figures. The observed agreement is therefore \emph{equally consistent} with a
depth-robust prefactor and with a depth correction of exactly that size: it
discriminates between them not at all. This is the reason that matters, and note
that it does not depend on settling where $\xiaKappaMean$ came from. Whether the number is
a genuine recalculation that happens to land on the depth correction, or was
synthesised from it, the comparison carries the same information either way,
which is none.

\emph{What survives.} The plateau value, the exponents and their intervals are
fixed by the Gifford three-width fit and were never computed from this
comparison, so none of them moves. What cannot be claimed is that they are
corroborated \emph{across laboratories} or \emph{across depth}. The exponents'
independent corroboration is C6's Davis cross-check (an unrelated laboratory
and device family), and the depth factor's is the single-gap CFD of C3d.
$\rightarrow$ \emph{open: reported as a \textbf{lead}, not as a cross-paper
agreement, and carrying no evidential weight in this appendix.}
\end{description}

\field{Residual risk.} One open item, named: \texttt{xia\_anchor\_unverified}.
It bears on the cross-paper comparison only. The fitted plateau value, the
exponents and their intervals do not rest on it, and nothing downstream (the
operating law, C6's exponents, C7's band) is recomputed as a result, because
none of them ever took an input from it. Closing it would need a cutoff
determination for Xia's devices made independently of Gifford Fig.~3, which the
primary paper does not report; it will not be closed from the literature.

\field{Status.} \textbf{Partially closed: warranted on CIF data; the cross-paper
anchor is a lead, not a corroboration.}

% ===========================================================================
\subsection*{C3c --- The knee is real (defended on CIF data alone)}

\field{Summary.} A single power law cannot fit the data ($\chired=\chiGlobal$) and
$\kapp$ swings $3.7\times$ from plateau to foot, so the prefactor is genuinely
size-dependent: a knee. DLD theory, correctly re-normalised, has no knee
(constant $\ksteric$), so the feature is CIF-specific. Because a size-independent
Stokes-flow solve cannot produce a size-dependent knee, and the CFD finds no
flow funnelling, the knee is attributed to a finite-particle capture-regime
change (mechanism class: Doyeux's finite-sphere displacement). Existence is
warranted now; the mechanistic shape is deferred to finite-particle work.

\field{Claim.} $\kapp$ is not constant: a single global power law is rejected,
and $\kapp$ falls from $\plateauKappa$ (plateau) to ${\approx}\footKappa$ (foot) through a knee
at $\dwc\approx0.03$. The knee is attributed to a finite-particle capture-regime change; its
existence is warranted on CIF data, its mechanism deferred to future
finite-particle work.

\field{Scheme.} Argument from Sign for the \emph{existence} $+$ Abduction for
the \emph{mechanism}; the digitization undercutter is closed separately, by
Evidence-to-Hypothesis in falsification form.

\field{Premises (derived).}
\begin{description}
\item[P-C3c.1 (global misfit).]
The global single power law is rejected at $\chired=\chiGlobal$ (frozen regression
test: it must be rejected); a constant plateau-$\kappa$ mispredicts the foot by
$269\%$ (frozen test~2). \prov{code: B3/B8}
\item[P-C3c.2 (the ramp).]
$\kapp$ ramps from $\plateauKappa$ (plateau) to $\footKappa$ (foot, $\dwc<\footUwc$), a
$3.7\times$ swing; the foot $\kappa$ is $1.33$--$1.34$ across all three
independent digitizations. \prov{code: B4}
\item[P-C3c.3 (DLD is silent).]
In the CIF normalisation $\kapp=\dcrit/\!\left(g\sqrt{\varepsilon/3}\right)$,
DLD theory gives a near-constant \emph{apparent} prefactor, Davis
$\davisKappaLo$--$\davisKappaHi$ (recomputed from eqn~2-22), Inglis $\inglisKappaLo$--$\inglisKappaHi$ (exact cubic
root), across $\varepsilon=0.01$--$0.5$: no knee, and never below the steric
floor. \emph{These are $\kapp$, not $\ksteric$}, which is exactly $2$ (C9); the
distinction matters because the CIF knee drops \emph{below} $2$ in the apparent
prefactor, a regime DLD never probes, and it is the apparent prefactor that both
literatures report. The Davis/Inglis small-$\varepsilon$ values sit just above
$\ksteric=2$, supporting the steric floor of C2 and C9, not the knee.
\prov{code: B7}
\item[P-C3c.4 (the finite-particle mechanism class).]
Doyeux 2011 exhibits the relevant mechanism: a finite sphere is displaced
relative to the fluid dividing line, and the attraction they measure is
``directed towards the low flow rate branch'' \cite[abstract]{b21}, i.e.\ toward
the (low-flow) side channel. (The classical Zweifach--Fung reading points the
other way; it is stated, and answered from the same abstract, at P-C3d.6.) This
displacement is a finite-size effect, absent for a point tracer and growing
with particle radius, so it raises the effective cutoff most for the largest
particles (high $\kapp$ on the plateau) and vanishes for the smallest, where
$\kapp$ falls toward and then below the steric floor to the sub-steric foot
(P-C3c.5), reproducing the knee's direction; the
precise shape $\kapp(\dwc)$ is deferred to future finite-particle simulation.
\item[P-C3c.5 (a known size-dependent extraction effect: bounded, and not the
knee's cause).]
Dinh reports that ``recovery of the smallest 2~fL PLTs ($a\approx1\um$) in the
filtrate was consistently lower than the percent of extracted fluid volume''
\cite[p.~917]{b5}, so size-dependent departures from tracer behaviour do
occur. Two considerations bear on whether this effect could be the knee. \emph{First, its
sign would lower, not raise, $\kapp$ at small $d$:} suppressing the recovery of
the smallest particles narrows the effective lamina they sample
($\delta_{\mathrm{eff}}<\delta_{\mathrm{nominal}}$), and since
$\kapp=\ksteric\,(\delta_{\mathrm{eff}}/\delta_{\mathrm{nominal}})$, the apparent
prefactor falls. Crowding therefore \emph{shares} the knee's direction and could
be confounded with it, so its sign alone cannot rule it out. \emph{Second, and
decisively, it is a concentration-/flow-dependent many-body effect,
not an intrinsic cutoff property:} Dinh's own sweep shows the 2~fL
filtrate-recovery difference varying with operating conditions
($23.2\rightarrow9.8$ percentage points as the flow rate rises from $30$ to
$90\,\mu$L/min) \cite[p.~917]{b5}, so it is minimal in the dilute,
single-particle regime the rigid-bead calibration assumes. Crowding thus neither
produces the calibration knee nor supplies the size-dependent $\feff$ the
alternative hypothesis (CQ2) would need.
\end{description}

\field{Critical questions.}
\begin{description}
\item[CQ1: Strength of the sign \textup{(\texttt{more\_reliable})}.] The misfit
is strong and quantified:
$\chired=\chiGlobal$ ($30.6$--$31.4$ across digitizations); plateau/foot differ by
$3.7\times$; no plausible digitization error bridges this (foot $\kappa$
$1.33$--$1.34$, global rejection $30.6$--$31.4$ in every digitization). Scheme:
Argument from Sign, strength CQ. The graph carries this as an \emph{open}
exception on \texttt{c3c\_exists}, not because the sign is weak, but because
``is there a more reliable indicator?'' cannot be closed by the same data that
raised it. $\rightarrow$ \emph{answered on the evidence available: strong.}
\item[CQ2: Alternative explanations of the sign
\textup{(\texttt{more\_coherent\_explanation})}.] The live alternative, that
the apparent knee is an effective per-gap fraction $\feff$ that varies (with
device or with particle size) rather than a steric regime change, is bounded
but not eliminated. (i)~The exponents are shift-invariant, so any \emph{uniform}
$c$ leaves $\beta,\gamma$ unchanged (C6 / Appendix~A Result~6), disposing of a
uniform $\feff$ but not a size-dependent one; (ii)~Doyeux's finite-particle
displacement (P-C3c.4) supplies a \emph{positive} steric mechanism that
reproduces the knee's direction, which a size-dependent $\feff$ would have to
mimic without an independent driver; (iii)~DLD is silent (P-C3c.3); (iv)~the one
identified size-dependent extraction effect, crowding (P-C3c.5), is
concentration-gated and minimal in the dilute calibration, so the alternative
lacks an operative driver there. Scheme: Abduction, more-coherent-explanation
CQ (the exception this question is named for, and the scheme of
\texttt{c3c\_mech}). $\rightarrow$ \emph{answered: the size-dependent $\feff$ alternative is bounded,
not ruled out, but it contests only the \textbf{mechanism}; it cannot defeat
the knee's \textbf{existence}, because a size-dependent $\feff$ is itself a form
of size-dependence (see ``Containment'' below). The steric reading of the
mechanism is favoured; the mechanism is deferred to future finite-particle work.}
\item[CQ3: Mechanism honesty.] The knee is not ``3-D funnelling.'' Two
independent reasons: (a)~a Stokes-flow solve is size-independent by construction,
so it cannot produce a size-dependent knee at all; (b)~the CFD measured
$\kinplane=1$, so the flow does not funnel (C3d, Methods~\S4.3). Scheme:
Established Rule (size-independence of Stokes flow) $+$ numerical verification (CFD).
$\rightarrow$ \emph{answered: the mechanism is finite-particle and deferred; the
``funnelling'' reading is retracted.}
\end{description}

\field{Residual risk.} Two open exceptions, named in the graph:
\texttt{more\_reliable} on the knee's \emph{existence} and
\texttt{more\_coherent\_explanation} on its \emph{mechanism}. The mechanism
requires future finite-particle simulation; the size-dependent $\feff$
alternative is bounded but not eliminated.

\field{Status.} \textbf{Warranted on CIF data (existence); mechanism rival open.}

\field{Containment of the $\feff(d)$ alternative.} The only \emph{rival} still
open in this appendix is one hypothesis: a \emph{size-dependent} effective flow
fraction $\feff(d)$ in place of the finite-particle steric account. (An open \emph{rival} is not the same object as an open critical
question: \texttt{feff\_arg} is the single node in the graph whose status is
\emph{open rival}, whereas the \cifOpenCQ\ undischarged critical questions of
\S\ref{sec:b-howtoread} are exception atoms spread over \cifNodesWithOpenCQ\
nodes, two of them on this very node.) Its reach is exactly two nodes. It undercuts (a)~the knee
\emph{mechanism} (this node's abductive attribution) and (b)~the $\sqrt{c}\,C$
\emph{decomposition} (C3d). It does \emph{not} touch the operating law (C3$'$),
the plateau exponents $\beta=\tfrac12,\gamma=1$ (C6, shift-invariant to any uniform $c$;
the size-dependent case is delegated to this node), or the knee's
\emph{existence}. Existence is immune for a structural reason: a size-dependent
$\feff(d)$ is \emph{itself} a form of size-dependence of the apparent prefactor,
so it \emph{co-signs} ``$\kapp$ varies'' rather than defeating it; the contest
is over the \emph{mechanism} of the variation, never its reality. Granting the
exception therefore flips only the mechanism and the decomposition to open; the
law, the lemmas, and existence are unmoved.

% ===========================================================================
\subsection*{C3d --- Interpretation of $\kappa$'s magnitude (partially closed)}

\field{Summary.} The measured plateau prefactor ($\plateauKappa$) is a product of five
interpretable factors. The steric factor ($2$) follows from C1; single-gap CFD
across the full deep-CIF aspect-ratio range settles two more: no in-plane
funnelling ($\kinplane=1$, a one-sided null with zero measured bulge) and a depth
factor matching the analytic duct law to $\le0.6\%$, and in doing so reverses
the earlier ``apertures funnel more'' intuition and retires the ``3-D
funnelling'' reading of the knee. What remains is a single lumped residual, a
degenerate product of per-gap over-extraction and cumulative finite-particle
displacement, both tied to the unpublished gap count; it is explicitly not
recoverable from the device-level filtrate fraction $F$. An independent
branch-flow group likewise finds its measured cutoffs above prediction and splits
the excess into the same two mechanism classes, corroborating the picture.

\field{Claim.} The plateau prefactor decomposes as
\[
  \kapp \;=\; \ksteric\cdot\kinplane\cdot\Kdepth\cdot\sqrt{c}\,C
        \;=\; 2\cdot 1\cdot R(\alpha)\cdot\sqrt{c}\,C .
\]
Three factors are settled: $\ksteric=2$, $\kinplane=1$ (CFD),
$\Kdepth=R(\alpha)\in[0.89,\,0.98]$ (CFD $+$ duct law), leaving the single
lumped product $\sqrt{c}\,C$ as the one open quantity. With $R(\alpha)$ folded
in, the excess over steric is depth-folded and $\alpha$-independent,
$\kapp/\ksteric=\kappVal/2=\residualFolded$;
dividing out the measured $\Kdcfd$, $\sqrt{c}\,C=\kapp/\!\left(\ksteric\Kdcfd\right)=
\residualSepShallow/\residualSepMid/\residualSepDeep$ across $\alpha=0.80/1.40/2.69$ (each
value $\times$ its $\Kdcfd$ returns $\residualFolded$).

% ----------------------- ARGUMENT GRAPH 2 (decomposition) ---
\begin{figure}[t]
\centering
\begin{tikzpicture}[
  font=\scriptsize,
  res/.style={draw, rounded corners, align=center, inner sep=3pt,
              text width=2.0cm, minimum height=13mm},
  ded/.style={draw, double, double distance=0.7pt, rounded corners,
              align=center, inner sep=3pt, text width=2.0cm, minimum height=13mm},
  open/.style={draw, dashed, rounded corners, align=center, inner sep=3pt,
               text width=2.25cm, minimum height=13mm},
  split/.style={draw, dashed, rounded corners, align=center, inner sep=3pt,
                text width=3.0cm, minimum height=10mm},
  corr/.style={draw, rounded corners, align=center, inner sep=3pt,
               text width=3.6cm, minimum height=12mm},
  op/.style={font=\large},
  down/.style={-{Stealth[length=1.8mm]}, thick},
  link/.style={-{Stealth[length=1.8mm]}, dashed},
]
% --- the decomposition spine ---
\node (lhs) at (-6.7,0) {$\kapp=\plateauKappa=$};
\node[ded]  (ks)  at (-5.0,0) {$\ksteric=2$\\ \scriptsize C1 / App.~A R5};
\node       (x1)  at (-3.75,0) {$\times$};
\node[res]  (ki)  at (-2.65,0) {$\kinplane=1$\\ \scriptsize CFD (line-sink)};
\node       (x2)  at (-1.4,0) {$\times$};
\node[res]  (kd)  at (-0.3,0) {$\Kdepth=R(\alpha)$\\ \scriptsize $[0.89,0.98]$};
\node       (x3)  at (0.95,0) {$\times$};
\node[open] (sc)  at (2.35,0) {$\sqrt{c}\,C$\\ \scriptsize \textbf{OPEN}};
% --- settled vs open annotation ---
\node[font=\scriptsize\itshape, text=black!60, align=center, anchor=south]
  at (-2.65,0.95) {three factors settled};
\node[font=\scriptsize\itshape, text=black!60, align=center, anchor=south]
  at (2.35,0.95) {one open};
% --- split of the residual ---
\node[split] (cc) at (0.6,-2.5)
  {$\sqrt{c}$: per-gap over-extraction ($c=\feff/\fnom$)};
\node[split] (CC) at (4.1,-2.5)
  {$C$: cumulative finite-particle displacement};
\node[font=\scriptsize\itshape, text=black!60, align=center] (N)
  at (2.35,-3.6) {both degenerate from published data, tied to gap count $N$};
\draw[link] (sc) -- (cc);
\draw[link] (sc) -- (CC);
% --- corroboration ---
\node[corr] (ih) at (-4.6,-2.5)
  {\textbf{Inglis \& Herman 2013} (p.1729): cutoffs above prediction, same two-part split
   (Doyeux ${\approx}2\times$, inferred $+$ over-extraction)};
\draw[link] (ih.east) -- (cc.west);
\end{tikzpicture}
\caption{Decomposition of the plateau prefactor (C3d). Three factors are
settled: $\ksteric=2$ (doubled border, deductive, from C1 / Appendix~A
Result~5), $\kinplane=1$ and $\Kdepth=R(\alpha)$ (both from single-gap CFD),
leaving a single open quantity, the lumped residual $\sqrt{c}\,C$ (dashed),
which splits into per-gap over-extraction $\sqrt{c}$ and cumulative
finite-particle displacement $C$, both degenerate from the published data and
tied to the unpublished gap count $N$. An independent branch-flow group reports
the same two-part split, corroborating the interpretation.}
\label{fig:c3d-decomp}
\end{figure}

\field{Scheme.} Numerical verification (single-gap CFD: a deductive computation
with no critical questions) for the settled factors $+$ Argument from an
Established Rule $+$ Argument from Expert Opinion (Doyeux) $+$ Defeasible Modus
Ponens closing the Zweifach--Fung undercutter that stands on it
(\texttt{c3d\_zf\_refuted}) $+$ Abduction for the open residual.

\field{Premises.}
\begin{description}
\item[P-C3d.1 ($\ksteric=2$).] From C1/CQ5 and Appendix~A Result~5 (Inglis
$\dcrit=2b$ $+$ Dinh main-text $a<w$ / $a>w$).
\item[P-C3d.2 (plateau $\kapp=\plateauKappa$).] From C3b.
\item[P-C3d.3 ($\kinplane=1$; CFD).] The single-gap Stokes campaign across
$\alpha=0.80/1.40/2.69$ (Runs~C/B/A) measures $\kinplane=1.000$ at every
fraction and station (max bulge $0.0000\um$); the dividing streamsurface
\emph{thins} toward the gap (line-sink behaviour), a one-sided null
($\delta_{\mathrm{bulge}}$ is a maximum, so $K=1$ only if the line never bulges;
funnelling would read $K>1$). One documented exception: Run~A at
$2\times10^{-4}$, $\kinplane=\text{n/a}$ (a low-fraction post-processing limit,
flagged not forced) [Methods~\S4.3]. \prov{code: B9}
\item[P-C3d.4 ($\Kdepth=R(\alpha)$; CFD $+$ duct law).] The CFD depth factor
matches the analytic rectangular-duct law to $\le0.6\%$ (CFD
$\KdepthCfdShallow/\KdepthCfdMid/\KdepthCfdDeep$ vs analytic
$\KdepthAnalyticShallow/\KdepthAnalyticMid/\KdepthAnalyticDeep$; per-aspect-ratio
deviations $\KdepthDevShallow\%/\KdepthDevMid\%/\KdepthDevDeep\%$, max
$\KdepthDevMax\%$ on Run~B); the cross-validation
$\texttt{duct.K\_depth}(100,80)=\KdepthAnalyticShallow$ against
$\text{CFD }R(\alpha{=}0.8)=\KdepthCfdShallow$ has COMSOL independently
confirming the code's duct module at a real device aspect ratio, to
$\KdepthDevShallow\%$.
\prov{code: B9; Appendix~A Result~9; glossary \xref{app:units_cutoff}{A.4}/\xref{app:units_cheatsheet}{A.12}}
\item[P-C3d.5 (the residual: arithmetic).] With $\ksteric=2$ and
$\kinplane=1$, the residual reduces to the single lumped product $\sqrt{c}\,C$,
reportable two equivalent ways: the depth-folded ($\alpha$-independent) excess
$\kapp/\ksteric=R(\alpha)\,\sqrt{c}\,C=\residualFolded$, or, dividing out the measured $\Kdcfd$,
$\sqrt{c}\,C=\kapp/\!\left(\ksteric\Kdcfd\right)=\residualSepShallow/\residualSepMid/\residualSepDeep$
across
$\alpha=0.80/1.40/2.69$. Here $\sqrt{c}$ is per-gap over-extraction
($c=\feff/\fnom$, a per-device constant from the local gap-to-channel resistance
ratio) and $C$ is cumulative finite-particle displacement over many gaps; the
two are degenerate from published data, both tied to the unpublished gap count
$N$. A sobering bound: if $C=1$, then $\sqrt{c}=\sqrt{c}\,C$ and
$c=(\sqrt{c}\,C)^{2}$ is of order $6$--$8$ ($\approx6$ with the depth factor
folded in, up to ${\sim}8$ with it explicit): the true per-gap extraction would
be several-fold the nominal design fraction, which is why the split cannot be
hand-waved. \prov{code: B6}
\item[P-C3d.6 / D-Q1 (Doyeux 2011, verbatim): the direction, and the standing
objection to it.]
This is the premise the Scheme field names as Argument from Expert Opinion, and
the graph carries it as \texttt{c3d\_doyeux\_direction}. Doyeux \emph{et al.}\
set out to test the classical reading of branch partitioning, on which particles
are attracted to the \emph{high}-flow branch, and report the reverse: ``such an
attraction does not occur but is, on the contrary, directed towards the low flow
rate branch'' \cite[abstract]{b21}. In CIF the side channel \emph{is} the
low-flow branch, so the displacement is toward extraction; it therefore raises
the cutoff rather than lowering it, and it grows with particle radius, which is
the size-dependence C3c needs (P-C3c.4).

\emph{The Zweifach--Fung objection, and why it does not bite.} The graph carries
\texttt{zweifach\_fung\_objection} as an undercutter on this premise, and it is
the first question a reader of the blood-flow literature will raise: the
Zweifach--Fung effect is the long-observed \emph{enrichment} of particles in the
high-flow branch, which looks like a direct contradiction. Doyeux answer it in
the same abstract, by separating two effects rather than denying either:
``the enrichment in particles in the high flow rate branch is shown to be mainly
a consequence of the initial distribution in the inlet branch''
\cite[abstract]{b21}. A finite sphere's centre cannot approach a wall closer than
one radius, so the near-wall streamtube is already depleted of centres before the
bifurcation; the enrichment is that inlet depletion, and it is compatible with a
particle-scale attraction of the opposite sign. Both hold at once, and only the
attraction is a force on the particle. The graph closes the undercutter at
\texttt{c3d\_zf\_refuted} (Defeasible Modus Ponens, warranted). Inglis \& Herman
state the same contrast independently, in the sentence quoted next.
\item[P-C3d.7 / IH-Q3 (Inglis \& Herman 2013, p.~1729, quoted continuously):
independent two-part corroboration.]
``The experimentally observed cutoff diameters are significantly higher than the
cutoff sizes determined from our design process. This is in agreement with
Doyeux \emph{et al.}\ who showed that contrary to the situation in the
Zweifach--Fung effect, particles are attracted to the low flow rate side
of a branch. Through experiment and modelling Doyeux \emph{et al.}\ found that
excluding a particular size particle from a branch may require a streamline
width that is as much as 50\% smaller than the particle radius. A second
contributing reason is that the devices, as fabricated, were approximately
48~$\mu$m deep, as opposed to the 50~$\mu$m as designed. This increases the main
channel resistance more than it does the branch channel resistance because the
branches are only 25~$\mu$m wide. This increased main channel resistance will
mean more fluid flow in the branches, increasing the critical size at each
branch'' \cite[p.~1729]{b18}.

This passage is quoted \emph{whole}: no ellipsis, and nothing elided but the two
reference superscripts, which is why no ellipsis marks them. Two glyphs are set
in the appendix's house form because the OCR layer cannot distinguish them: the
depth and width units are the paper's micrometres (the layer flattens $\mu$ to
m), and the dash in Zweifach--Fung is set as an en dash, as the bibliography
sets it. $\rightarrow$ an independent
branch-flow group reports cutoffs above prediction and decomposes the excess
into the same two mechanisms. Their \emph{first} reason is the Doyeux
direction, stated with its own anti-Zweifach--Fung contrast, so the
corroboration at P-C3d.6 is theirs as well as Doyeux's; the finite-particle
factor is ${\approx}2\times$, inferred from their ``50\% smaller than the
particle radius.'' Their \emph{second} is over-extraction from increased
main-channel resistance, the $\sqrt{c}$ factor, and they give the reason for it:
a shallow mould raises the \emph{main} channel's resistance disproportionately
\emph{because the branches are narrow to begin with}, so flow redistributes into
the branches. That is the post-CFD CIF picture in kind, a finite-particle knee
$+$ $\sqrt{c}$ over-extraction, and not flow funnelling.
\end{description}

\field{Critical questions.}
\begin{description}
\item[CQ1: Computation / measurement validity (CFD).] The CFD result is
trustworthy: $\kinplane=1$ is a one-sided null with max bulge $0.0000\um$ across
all runs/stations; $R(\alpha)$ reproduces the analytic duct law to $\le0.6\%$
and reproduces a known code value at a real aspect ratio;
$\dskim\propto\sqrt{f}$ to $R^{2}\ge0.9999$; all wired to frozen regression
stamps against \texttt{REFERENCE\_OUTPUT*.txt}. Scheme: numerical verification (deductive computation).
$\rightarrow$ \emph{answered: yes} (standing caveat: raw COMSOL fields are not
archived; the reference tables $+$ post-processing scripts are the
reproducibility surface; a cheap $0.5\times/2\times$ mesh-convergence point would
close the obvious referee question).
\item[CQ2: Direction (does the aperture funnel more than DLD posts?).] No. The
aperture funnels less. The earlier intuition (CIF wall-apertures should funnel
more than DLD posts, a weak ${\sim}1.2\times$ Davis prior) is reversed by the
CFD: $\kinplane=1$, below $1.2\times$. The aperture is a line-sink and the
captured layer does not bulge. Scheme: numerical verification overruling a prior argument from analogy,
C2's \texttt{c2\_funnel\_rival}, which the graph carries as defeated.
$\rightarrow$ \emph{answered: reversed; no funnelling.}
\item[CQ3: Can the residual be read out of $F$?] No. The over-extraction is
not ``fixed by $F$.'' The rigorous outlet identity $F=2R_{c}/(R_{s}+2R_{c})$
($\FclosedVone$ and $\FclosedVthree$ at Dinh's V1 and V3) is outlet-geometry-only and
$N$-independent (C4); $0.3516$ is a cumulative-yield number (the product $N\fgap$
that $\ln(1-F)=-2N\fgap$ returns at $F\approx0.505$), not $F$, and
is historical rather than a live value (its provenance and arithmetic are given
under C4's cautionary tale). The
cutoff depends on the local per-gap $c$, narrowable by a local
resistance-network calculation (the \texttt{network} module is the seed), not by
$F$ (reading $\kappa$ out of $F$ inverts the C4 separation lemma). Scheme:
Established Rule (C4) blocking a category error. $\rightarrow$ \emph{answered:
no.}
\end{description}

\field{Residual risk.} Two open exceptions, named in the graph. The first is
\texttt{degenerate\_from\_published\_data}: identifiability, meaning that the two
factors cannot be separated from what the papers actually report. The second is
\texttt{more\_coherent\_explanation}, the $\feff$ rival, shared with C3c. The
$\sqrt{c}\,C$ residual is open and tied to $N$; resolved by finite-particle work
(the $C$ factor) $+$ a local resistance-network calculation and/or micrograph gap
count (the $\sqrt{c}$ factor), not by the single-gap solve (now done) and not by
the $F$-identity (decoupled). A sibling ${\sim}\siblingSqrtcPct\%$ 3-D-junction uncertainty rides
on $\sqrt{c}$. \emph{Exhaustiveness}, as distinct from identifiability, is
answered downstream: within the model's premises, C9d closes the residual list to
three model gaps.

\field{Status.} \textbf{Partially closed: $\kinplane$, $\ksteric$, $R(\alpha)$
settled; $\sqrt{c}\,C$ open.}


\clearpage
\section*{C5 --- Consilience (sign-level): the skim law recovers the corrected
Xia (2022) scaling and excludes the 2016 original}

\noindent\textbf{Summary.} On the plateau, where the apparent prefactor $\kapp$
is constant, the skim law forces the centre-channel width into the per-gap
fraction as $\fgap\propto\wc^{-2}$ at fixed cutoff. This is exactly the 2022
author-corrected Xia design equation and the reciprocal of the since-corrected
2016 original. Agreement on the \emph{form} is a shared-premise consistency
check: it reuses the parabolic lamina width of C2 and the proportionality of
C3a, which are Xia's own two assumptions, and the author correction is a
same-lineage self-correction, so neither is independent confirmation. Agreement on the
\emph{direction} is the load-bearing result: the law contradicts the 2016
equation as published and agrees only with the 2022 correction, a retrodiction
(it would have flagged the 2016 error).

% ----------------------- ARGUMENT GRAPH ---------------------
\begin{figure}[t]
\centering
\begin{tikzpicture}[
  font=\scriptsize,
  ded/.style={draw, double, double distance=0.7pt, rounded corners,
              align=center, inner sep=4pt, text width=7.4cm, minimum height=12mm},
  ev/.style={draw, rounded corners, align=center, inner sep=3pt,
             text width=4.0cm, minimum height=12mm},
  excl/.style={draw, dashed, rounded corners, align=center, inner sep=3pt,
               text width=4.0cm, minimum height=12mm},
  concl/.style={draw, rounded corners, align=center, inner sep=4pt,
                text width=9.6cm, minimum height=11mm},
  pred/.style={-{Stealth[length=2mm]}, thick},
  up/.style={-{Stealth[length=2mm]}, thick},
  nt/.style={font=\scriptsize\itshape, text=black!60, align=center},
]
% --- the deductive conditional ---
\node[ded] (cond) at (0,0)
  {\textbf{Conditional (P-C5.1, deductive).}\ skim law
   $\delta=\wc\sqrt{\fgap/3}$ (C2) $+$ constant $\kappa$ (C3a)
   $\Rightarrow$ $\fgap\propto\wc^{-2}$, i.e.\ $\fgap=(\wcref/\wc)^{2}\fgapref$};
% --- two pieces of evidence ---
\node[ev]   (corr) at (-3.5,-2.5)
  {\textbf{Xia 2022 correction}\\ $\bigl(w_c(\mathrm{ref})/w_c(i)\bigr)^{2}$\\
   \textbf{matches} the derived form};
\node[excl] (orig) at (3.5,-2.5)
  {\textbf{Xia 2016 original}\\ $\bigl(w_c(i)/w_c(\mathrm{ref})\bigr)^{2}$\\
   \textbf{inverted} (excluded)};
% --- conclusion ---
\node[concl] (cl) at (0,-5.0)
  {\textbf{C5.}\ The law selects the corrected scaling and excludes the
   published original, a sign-level retrodiction (it would have flagged the
   2016 error). Shared premises (C2, C3a) $\Rightarrow$ corroboration, \emph{not}
   independent confirmation. \ \textsf{[warranted]}};
% --- edges ---
\draw[pred] (cond) -- (corr) node[midway,above left,nt]{predicts $\checkmark$};
\draw[pred] (cond) -- (orig) node[midway,above right,nt]{contradicts $\times$};
\draw[up]   (corr) -- (cl);
\draw[up]   (orig) -- (cl);
% --- caveat ---
\node[nt, anchor=north, text width=11cm] at (0,-6.05)
  {The \emph{form} reuses Xia's own two assumptions (C2 $+$ C3a) and the author
   correction is same-lineage, so the form carries no independent weight; the load-bearing
   residue is the \emph{direction}.};
\end{tikzpicture}
\caption{Argument graph for C5 (verification-form Evidence-to-Hypothesis). The
deductive conditional (doubled border) predicts a single cross-width schedule,
$\fgap\propto\wc^{-2}$. The 2022 author correction matches it; the 2016
published original is its reciprocal and is excluded. Because the \emph{form}
follows from Xia's own two assumptions (C2 $+$ C3a) and the author correction is
a same-lineage self-correction, the confirmation is sign-level (the \emph{direction})
rather than an independent verification. Border style follows the overview map:
doubled = deductive, single = warranted/corroborating, dashed = the excluded
rival.}
\label{fig:c5-argmap}
\end{figure}

% ----------------------- CLAIM ------------------------------
\field{Claim.} With $\kappa$ held constant (the plateau / proportionality
regime),
\[
  \dcrit = \kappa\, \wc \sqrt{\tfrac{\fgap}{3}}
  \quad\Longrightarrow\quad
  \fgap\propto\wc^{-2},
  \qquad\text{i.e.}\qquad
  \fgap(i) = \left(\frac{\wcref}{\wc(i)}\right)^{2} \fgapref ,
\]
which is the 2022 author-corrected Xia design equation (their eqn~(1)).

% ----------------------- SCHEME -----------------------------
\field{Scheme.} Argument from Evidence to a Hypothesis, verification form. The
\emph{conditional} (P-C5.1) is that if the skim law holds, the cross-width
schedule must be $\fgap\propto\wc^{-2}$. The \emph{evidence} (P-C5.2, P-C5.3) is
that the corrected literature reports that schedule while the erroneous original
reported its reciprocal, so the evidence discriminates rather than merely
fitting. The \emph{conclusion} is that the skim law is corroborated in direction.

% ----------------------- PREMISES ---------------------------
\field{Premises.}
\begin{description}

\item[P-C5.1 (derived conditional; exact, code-checked).]
Holding the lamina width $\delta = \wc\sqrt{\fgap/3}$ fixed (equivalently,
$\dcrit$ fixed \emph{at constant} $\kappa$) and solving for $\fgap$,
\[
  \wc\sqrt{\tfrac{\fgap}{3}} = \wcref\sqrt{\tfrac{\fgapref}{3}}
  \;\Longrightarrow\;
  \fgap = \frac{\wcref^{2}}{\wc^{2}}\, \fgapref
  \quad\text{(exact).}
\]
The constant-$\kappa$ step is the plateau / proportionality assumption (identical
to C3a and to Xia's assumption~(b)); it is the premise CQ1 and CQ4 address, made
explicit here rather than left tacit.
\emph{Where the $3$ comes from.} $\delta=\wc\sqrt{\fgap/3}$ is the deep-channel
limit of a finite-aspect-ratio form, $\delta=\wc\sqrt{\fgap/c_2(\alpha)}$ with
$c_2\to3$ as $\alpha\to\infty$; the depth factor $\Kdepth=\sqrt{3/c_2(\alpha)}$
is that same rectangular-duct series, evaluated by \texttt{K\_depth} in P-C2.5.
It is folded into the \emph{apparent} prefactor rather than applied to the data
(\S\ref{sec:b-howtoread}), so the depth factor sits \emph{inside} the
constant-$\kapp$ assumption rather than beside it: $\Kdepth$ does vary along a
taper, and that variation shows up as variation in $\kapp$, which is the premise
CQ4 examines. No separate aspect-ratio premise is needed here. \prov{code: scaling.f\_gap\_of\_wc
$\rightarrow$ $\fgapref\wcref^{2}/\wc^{2}$; verify\_symbolic asserts it equals
$(\wcref/\wc)^{2}\fgapref$; verify\_numeric confirms cutoff-invariance from Xia's
OR-7 reference values ($\wcref=100\um$, $\fgapref=5.8\times10^{-4}$), evaluated at
$\wc=150\um$, the RS-7 taper's minimum centre-channel width, to $10^{-18}$;
Appendix~A Result~7}

\item[P-C5.2 (corrected Xia scaling; verbatim).]
The 2022 author correction states that in eqn~(1) ``the terms $w_c(i)$ and
$w_c(\mathrm{ref})$ are inverted,'' and that it ``should read''
\[
  \fgap(i) = \left(\frac{w_c(\mathrm{ref})}{w_c(i)}\right)^{2}
    \fgap(\mathrm{ref}),
\]
exactly the direction the derived conditional (P-C5.1) forces \cite{b3}.

\item[P-C5.3 (the 2016 original was the inverted form; verbatim).]
The published 2016 eqn~(1) reads
\[
  \fgap(i) = \left(\frac{w_c(i)}{w_c(\mathrm{ref})}\right)^{2}
    \fgap(\mathrm{ref}),
\]
the wrong direction as $\wc$ narrows along the taper
($300\to150\um$, with $\wcref=100\um$): at fixed cutoff a narrower channel must
divert a \emph{larger} fraction, so $\fgap$ must rise as $\wc$ falls, whereas the
2016 form makes it fall. The 2022 author correction repairs it in exactly the
direction P-C5.1 predicts; the erroneous form is the published eqn~(1)
\cite[eqn~(1)]{b2}, and the correction reproduces it \cite[p.~1]{b3}.

\end{description}

% ----------------------- CRITICAL QUESTIONS -----------------
\field{Critical questions.}
\begin{description}

\item[CQ1: Could the evidence be explained without the hypothesis?
(alternative-explanation).]
On the \emph{form}, yes: $\fgap\propto\wc^{-2}$ follows directly from Xia's two
assumptions, the parabolic profile (C2) and the proportionality of cutoff to
lamina width (C3a), so the form agreement restates one relationship rather than
supplying a second, independent witness. Answer (by established rule /
derivation): shared-premise consistency, disclosed; the form carries no
independent confirmatory weight.

\item[CQ2: Is the evidence reliable and independent of the hypothesis's
source? (evidence-reliability).]
Reliable: a published author correction, transcribed from the original. But
\emph{not} third-party independent: it is a self-correction by the
Gifford/Xia/Shevkoplyas lineage whose devices and data this analysis builds on.
The confirmatory reading is therefore \emph{the law is consistent with the
field's accepted (corrected) form and inconsistent with the since-corrected 2016
original},
not \emph{an independent group verified the law}. Answer: reliable but
same-lineage; weight bounded accordingly.

\item[CQ3: Was the prediction risky; does the evidence discriminate?
(severity).]
Yes, in direction: the law's $\wc^{-2}$ contradicted the \emph{published} 2016
eqn~(1) and matched only the later correction, so the evidence excludes the
then-published rival. The risk is against the literature; given the shared
premises the direction is determinate, so the confirmation is modest but real.
Answer: discriminates on sign, a successful retrodiction of the corrected
direction; equivalently, the law would have flagged the 2016 error.

\item[CQ4: Does the conditional hold without hidden assumptions?
(conditional-robustness).]
The conditional assumes $\kapp$ independent of $\wc$ at fixed cutoff: the
plateau, where $\kapp$ is flat (CV~$\kappCV\%$). RS-7, the one RS device whose
taper Xia prints, meets that only where it is narrow: across its $300\to150\um$
taper the cutoff-to-width ratio runs $\dwc=0.023$--$0.047$ (RS-8 and RS-9 would
run $0.027$--$0.053$ and $0.030$--$0.060$ on the same taper; their starting
widths are not printed), and $0.07$--$0.09$ is the ratio at the $100\um$
reference width. Off the plateau, on C7's fit of $\kapp(\dwc)$, the implicit law
holds the cutoff only if the gap fraction falls more slowly than $\wc^{-2}$: at
$300\um$ the corrected form supplies about half of what a $7\um$ cutoff requires,
and RS-7's sections wider than about $236\um$ fall below the saddle-node C7
derives. The direction does not depend on the plateau: the required gap fraction
still falls as $\wc$ widens, and the inverted form errs the other way, about
$42\times$ too much at $300\um$. Answer: the direction holds across the taper; the
exact $\wc^{-2}$ form needs the plateau, which RS-7 reaches only on its narrow
sections.

\end{description}

% ----------------------- STATUS -----------------------------
\field{Status.} \textbf{Warranted} as a sign-level cross-check: the law selects
the corrected scaling and excludes the published original. It is \emph{not}
independent confirmation, because the form agreement reuses the premises of C2
and C3a and the author correction is a same-lineage self-correction. The
confirmatory residue is therefore the \emph{direction} (the law would have flagged the 2016 error), which
holds across RS-7's whole taper; the exact $\wc^{-2}$ form needs the plateau, which
RS-7 reaches only on its narrow sections (CQ4).


\clearpage
\section*{C6 --- Cutoff-law exponents: $\beta=\tfrac12$ (in $\fgap$), $\gamma=1$
(in $\wc$)}

\noindent\textbf{Summary.} The cutoff law's exponents are $\tfrac12$ on the
per-gap fraction $\fgap$ and $1$ on the channel width $\wc$. They are not
free-standing empirical findings: they are the direct corollary of the
skim-width identity proved in C2 ($\delta=\wc\sqrt{\fgap/3}$) together with the
proportionality of C3a ($\dcrit\propto\delta$). The plateau fit returns
$(\expBeta,\,\expGamma)$; the combined width--fraction scaling is corroborated by the
Xia 2022 corrected relation $\fgap\propto\wc^{-2}$ (C5); and $\beta$ is
recovered independently, in another laboratory and another device family, by
Davis's DLD correlation (P-C6.5). Because $\beta$ and
$\gamma$ are slopes in a log--log regression, they are \emph{exactly} invariant
to any \emph{uniform} mis-estimate of the per-gap fraction; this is the one
confound the prefactor's physical decomposition
($\kapp=\ksteric\cdot\kinplane\cdot\Kdepth\cdot\sqrt{c}\,C$, C3d) does
\emph{not} survive (its $\sqrt{c}$ absorbs a uniform $c$), whereas the exponents
do.

% ----------------------- ARGUMENT GRAPH ---------------------
\begin{figure}[t]
\centering
\begin{tikzpicture}[
  font=\scriptsize,
  ded/.style={draw, double, double distance=0.7pt, rounded corners,
              align=center, inner sep=4pt, text width=9.0cm, minimum height=11mm},
  claim/.style={draw, rounded corners, align=center, inner sep=4pt,
                text width=7.0cm, minimum height=10mm},
  sup/.style={draw, rounded corners, align=center, inner sep=3pt,
              text width=2.7cm, minimum height=18mm},
  up/.style={-{Stealth[length=2mm]}, thick},
  nt/.style={font=\scriptsize\itshape, text=black!60, align=center},
]
% --- tier 1: deductive core ---
\node[ded] (core) at (0,0)
  {\textbf{Deductive core (P-C6.1).}\ $\delta=\wc\sqrt{\fgap/3}$ (C2) $+$
   $\dcrit\propto\delta$ (C3a) $\Rightarrow$ $\beta=\tfrac12$, $\gamma=1$
   \emph{exactly} (log-derivatives, $\kappa$ held fixed)};
% --- tier 2: claim ---
\node[claim] (cl) at (0,-2.25)
  {\textbf{C6.}\ $\beta=\tfrac12$ (in $\fgap$), $\gamma=1$ (in $\wc$).
   \ \textsf{[warranted]}};
% --- tier 3: four supports ---
\node[sup] (fit) at (-5.4,-4.7)
  {\textbf{Plateau fit (P-C6.2)}\\ $\beta=\expBeta\,[\betaCIlo,\betaCIhi]$\\
   $\gamma=\expGamma\,[\gammaCIlo,\gammaCIhi]$\\ ($\chired=\chiPlateau$)};
\node[sup] (inv) at (-1.8,-4.7)
  {\textbf{Invariance (P-C6.3)}\\ uniform $\feff=c\,\fgap$:\\ slopes unchanged
   (exact);\\ the prefactor's $\sqrt{c}$ does not};
\node[sup] (xia) at (1.8,-4.7)
  {\textbf{Consilience (P-C6.4)}\\ Xia 2022:\\ $\fgap\propto\wc^{-2}$\\
   (sign-level, C5)};
\node[sup] (dav) at (5.4,-4.7)
  {\textbf{Independent lab (P-C6.5)}\\ Davis 2008, DLD:\\ ${\sim}20$ devices,\\
   exponent $0.48$\\ (inside $[\betaCIlo,\betaCIhi]$)};
% --- edges ---
\draw[up] (core) -- (cl);
\draw[up] (fit) -- (cl);
\draw[up] (inv) -- (cl);
\draw[up] (xia) -- (cl);
\draw[up] (dav) -- (cl);
% --- discrimination note ---
\node[nt, anchor=north, text width=12cm] at (0,-6.35)
  {Law discrimination (CQ2--CQ3): the data select $\beta=\tfrac12$ (a lamina thinner than
   the depth); the linear law ($\beta=1$, a lamina wider than the depth) lies outside
   $[\betaCIlo,\betaCIhi]$, and the
   single global power law ($\beta\approx1.21$) is rejected at $\chired=\chiGlobal$.};
\end{tikzpicture}
\caption{Argument graph for C6. The exponents are a deductive corollary of the
C2 skim-width identity (doubled border), confirmed by the three-width plateau
fit, corroborated in sign by the Xia 2022 cross-width scaling and in exponent by
Davis's independent DLD device family, and, uniquely
among the law's parameters, \emph{exactly invariant} to a uniform per-gap
over-extraction $c$ (the confound the prefactor absorbs through $\sqrt{c}$). The
measured $\beta\approx\tfrac12$ positively selects the thin-lamina $\sqrt{\fgap}$ law over the
linear law and over the rejected single global fit.}
\label{fig:c6-argmap}
\end{figure}

% ----------------------- CLAIM ------------------------------
\field{Claim.} In the operating law
\[
  \dcrit = \kapp(\dwc,\,\mathrm{depth})\,\wc\,\sqrt{\fgap/3},
  \qquad \delta=\wc\sqrt{\fgap/3},
\]
the exponent of the per-gap fraction $\fgap$ is $\beta=\tfrac12$ and the exponent
of the channel width $\wc$ is $\gamma=1$.

% ----------------------- SCHEME -----------------------------
\field{Scheme.} Defeasible Modus Ponens on the C2 hydrodynamic derivation
(deductive), with Argument from an Established Rule for the log--log
shift-invariance, combined with Argument from Example twice over (the weighted fit
to the Gifford Fig.~3 data, and Davis's independent DLD device family) and
cross-checked against C5, whose own scheme is Evidence to a Hypothesis in
verification form, over the Xia 2022 corrected cross-width scaling. The named schemes fix the
matching critical questions below.

% ----------------------- PREMISES ---------------------------
\field{Premises.}
\begin{description}

\item[P-C6.1 (derivation; inherited from C2 $+$ C3a).]
Holding $\kappa$ fixed (the plateau), the $\sqrt{\fgap}$ dependence of the law
gives $\partial\ln\dcrit/\partial\ln\fgap=\tfrac12$ and the linear $\wc$
dependence gives $\partial\ln\dcrit/\partial\ln\wc=1$, both \emph{exactly}. The
$\sqrt{\fgap}$ is itself the leading term of the exact skim cubic
$\fgap=3x^{2}-2x^{3}$ ($x=\delta/\wc$) established in C2, so $\beta$ and $\gamma$
are corollaries of an already-proven identity rather than independent claims.
\prov{code: scaling.verify\_symbolic() asserts beta == 1/2, gamma == 1;
Appendix~A Result~6}

\item[P-C6.2 (plateau fit).]
The error-bar-weighted log--log fit on the plateau ($\dwc\ge\plateauUwc$,
$n=\plateauN$)
returns
\[
  \beta=\expBeta\;[\betaCIlo,\,\betaCIhi]\ (\text{model }\tfrac12),
  \qquad
  \gamma=\expGamma\;[\gammaCIlo,\,\gammaCIhi]\ (\text{model }1),
  \qquad \chired=\chiPlateau .
\]
Both model values $\tfrac12$ and $1$ lie inside their $95\%$ confidence
intervals. The intervals are conservative: the reported covariance is scaled by
$\chired$, i.e.\ the standard errors are inflated by $\sqrt{\chiPlateau}\approx1.95$,
exactly the factor that would bring the plateau $\chired$ to unity.
\prov{code: B3/B4}

\item[P-C6.3 (invariance to a uniform $\feff$ mis-estimate).]
Write the true per-gap fraction as $\feff=c\,\fgap$ with $c$ a device-level
constant. Then $\dcrit=\kappa\,\wc\sqrt{\feff/3}=\sqrt{c}\,\kappa\,\wc
\sqrt{\fgap/3}$: a uniform $c$ rescales the \emph{prefactor} ($\kappa\to
\sqrt{c}\,\kappa$) and leaves the exponents of $\fgap$ and $\wc$ untouched.
Equivalently, in the log--log regression a factor $c$ shifts only the
\emph{intercept} ($\ln\kappa$) and not the \emph{slopes}; the invariance is
algebraically exact, not approximate, and is confirmed numerically as identical
at $c\in\{1,\,2,\,3.5\}$ (the plateau $\beta=\expBeta$, $\gamma=\expGamma$ in every case).
This is the confound the prefactor's decomposition does \emph{not} survive (its
$\sqrt{c}$ absorbs $c$, C3d), whereas the exponents do. \prov{code: B6, ``EXPONENTS and KNEE robust to $c$''}

\item[P-C6.4 (consilience with the Xia 2022 corrected scaling; shared-premise,
C5).]
The Xia 2022 Author Correction states (verbatim) that the original equation
inverted the width terms and ``should read: $\fgap(i)=(\wcref/\wci)^{2}
\fgapref$.'' With $\beta=\tfrac12$ fixed by the parabolic profile, requiring a
matched cutoff across widths forces $\wc\sqrt{\fgap}=\text{const}$, i.e.\
$\fgap\propto\wc^{-2}$, precisely the corrected relation. This is a consilience,
\emph{not} an independent test: the $\wc^{-2}$ form follows from Xia's own two
assumptions (``approximately parabolic'' profile; diverted-streamline width
``roughly proportional to the effective c.d.''), the same physics C2/C3a rest
on. Its evidential value is in \emph{sign}: the law agrees with the
\emph{corrected} equation and excludes the originally published (inverted) form.
The empirical cross-width measurement of $\gamma$ remains the three-width Gifford
fit of P-C6.2. \prov{Xia 2022 Author Correction; Xia 2016 p.~10 (Sci Rep PDF
pagination); C5; Appendix~A Result~7}

\item[P-C6.5 (independent-lab cross-check; a different device family).]
Davis fits a DLD critical-diameter correlation to data from ``about 20 different
devices'' spanning $G = 1.3$--$38\um$ and $\varepsilon = 0.005$--$0.5$, with
polystyrene beads $0.9$--$22\um$, and obtains $\dcrit = 1.4\,G\varepsilon^{0.48}$
\cite[eqn~(2-22)]{b22}. The exponent $0.48$ lies inside this node's plateau
interval $[\betaCIlo,\,\betaCIhi]$, and it is load-bearing in its own paper rather than
incidental: ``equation 2-22 was used to design all the'' thresholds reported
there. This is the one support for $\beta$ that survives the withdrawal of
the entire CIF derivation, sharing with it neither laboratory, device family nor
data; it does not test $\gamma$, because his correlation is proportional to the gap
$G$ by its form. \emph{Its limit is the source's own, and it is the limit this node would
predict}: Davis records that against his measured points ``the parabolic
relationship underestimates the critical size,'' attributing the gap to ``short,
rounded posts'' rather than to the exponent. The cross-check is therefore an
agreement of \emph{exponent} and not of prefactor (his recomputed
$\kapp = \davisKappaLo$--$\davisKappaHi$ sits above the exact cubic's
$\inglisKappaLo$--$\inglisKappaHi$, C3c P-C3c.3), which is exactly
the split P-C6.3 asserts: a uniform mis-estimate moves the prefactor and never the
slopes. \prov{Davis 2008 eqn~(2-22), \S2.4 (pp.~22--24); graph node
\texttt{c6\_davis}; Appendix~A Result~6}

\end{description}

% ----------------------- CRITICAL QUESTIONS -----------------
\field{Critical questions.}
The exponents $\beta=\tfrac12,\gamma=1$ are a deductive corollary of the C2
identity and carry no critical questions \emph{as identities}; the questions
below test the \emph{regime} in which they apply and the empirical fit, not the
derivation.
\begin{description}

\item[CQ1: Is the rule (the derivation) correctly applied and sound?]
Yes. $\beta=\tfrac12$, $\gamma=1$ are a proven symbolic identity
(\texttt{scaling.verify\_symbolic()}), being the leading-order consequence of the
exact cubic $\fgap=3x^{2}-2x^{3}$ proved in C2; the symbolic and
finite-difference exponents agree and reproduce byte-for-byte.

\item[CQ2: Do the exponents hold globally, or only on a restricted subset?]
Only on the plateau, by construction, and this is a strength, not a concession.
A \emph{single} global power law across all $15$ points is \emph{rejected} by its
misfit ($\chired=\chiGlobal$; cf.\ C3c). The blended slopes that such a fit returns
($\beta\approx1.21$, $\gamma\approx1.55$) are the artifact of forcing one
straight line through a curve with a knee; they are therefore \emph{not adopted}
as the law's exponents. \prov{code: \texttt{global\_powerlaw\_misfit(with\_slopes=True)},
which returns the misfit and these two slopes from one regression} The global fit is reported alongside the plateau fit
precisely so it can be seen to fail: shown to be rejected, not suppressed.

\item[CQ3: Could a different flow regime yield different exponents?]
Yes, and the observed value \emph{discriminates} between the laws they give. In a
shallow, Hele-Shaw channel ($h\ll\wc$) a lamina much wider than the depth
($\delta\gg h$) samples a depth-averaged plug across the width,
$\fgap\propto\delta/\wc$, hence $\delta\propto\fgap$ and $\beta=1$, a
\emph{distinct} law; but $\beta$ is $\tfrac12$ whenever $\delta\ll h$, at any aspect
ratio, and CIF's laminae are orders of magnitude thinner than the channel depth
(Appendix~A Result~6). The plateau fit lands at
$\beta\approx\tfrac12$, with $\beta=1$ far outside its confidence interval
$[\betaCIlo,\betaCIhi]$, so the data positively selects the thin-lamina ($\sqrt{\fgap}$)
law over the linear one. CIF's moderate aspect ratio ($\alpha=h/\wc\approx
0.5$--$2.7$, not asymptotically deep) enters only through the depth factor
$R(\alpha)\in[0.81,0.98]$ (C2) folded into the apparent prefactor, never through the
exponent. For $\beta$ that is exact, since $R$ does not depend on $\fgap$. For
$\gamma$ it is not an identity, since $\alpha$ falls as $\wc$ widens across the fit:
folded into the prefactor, where C2 puts it, $R(\alpha)$ leaves $\gamma=\expGamma$;
applied to the lamina, it would move $\gamma$ to $1.24$ (interval $[1.04,1.45]$,
excluding $1$) at the assumed $80\um$ depth, the shallowest Gifford reports, but keep
$1$ inside the interval at the $125$ or $150\um$ it also reports
(\S\ref{sec:b-notation}).

\item[CQ4: Is $\beta=\tfrac12$ an artifact of the plateau cut or of the
figure digitization?]
No, on three independent counts. \emph{(i) Threshold:} sweeping the plateau cut
over $\dwc = 0.030$, $\plateauUwc$, $0.040$ and $0.045$ keeps $\beta\in[\expBeta,0.59]$
(the first two select the same $\plateauN$ points).
\emph{(ii)
Digitization noise:} perturbing each $f^{*}$ within its digitized $\sigma$ over
$2000$ draws gives $\beta$ median $\expBeta$ with $\mathrm{P}(0.4<\beta<0.7)=1.00$.
\emph{(iii) Independent re-digitization:} three independent digitizations of
Fig.~3 (the primary one and two re-digitizations, whose group labels were discarded
and recovered by sort-and-pair) give $\beta\in[0.50,
\expBeta]$ and reject the global law every time ($\chired\in[\chiGlobal,31.4]$).

\item[CQ5: Is the $n=\plateauN$ plateau sample representative?]
Yes, and the answer no longer leans on the cross-paper comparison it once cited.
The Xia reference triples (OR-7, RS-8, RS-9, all at $\wcref=100\um$) do run
through the plateau law to
$\kapp=\xiaKappaA/\xiaKappaB/\xiaKappaC$, close to the Gifford plateau $\kapp\approx\plateauKappa$,
but C3b's CQ2 establishes that the comparison is circular and, worse,
non-diagnostic (its agreement is the depth correction), so it is not evidence of
anything about the sample. The
answer instead rests on what is internal to this node and on C6's own
corroboration: the threshold sweep (CQ4~(i)), the digitization Monte Carlo and
the re-digitizations (CQ4~(ii)--(iii)), and the fact that these
exponents are established by \emph{shift-invariance} rather than by the plateau's
absolute value, so a representativeness worry about $n=\plateauN$ would have to
move
$\beta$ and $\gamma$, not $\kapp$, and the sweep shows it does not. Above all,
P-C6.5 answers the question from outside the sample altogether: Davis's ${\sim}20$
devices are an unrelated laboratory and a different device family, and they return
an exponent inside this node's interval.

\end{description}

% ----------------------- SCOPE NOTE -------------------------
\field{Scope note.} P-C6.3 establishes robustness to a \emph{uniform}
(device-level) $\feff$ mis-estimate only. A \emph{size-dependent} $\feff(d)$
would shift the slopes and is the separate alternative addressed (bounded but
not eliminated) in C3c, where it would manifest as the knee rather than as a
shift of the plateau exponents. C6 does not claim to exclude it; that is C3c's
burden.

% ----------------------- STATUS -----------------------------
\field{Status.} \textbf{Warranted; robust to a uniform $\feff$ mis-estimate.}
Residual risk: none on the exponents themselves. The size-dependent $\feff(d)$
possibility is delegated to C3c; the prefactor's physical decomposition is the
subject of C3d and is not part of this claim.


\clearpage
\section*{C7 (graph) --- The Dinh held-out consistency band}

\noindent\textbf{Summary.} Applying the fitted $\kapp(\dwc)$ collapse curve and
the skim law $\delta=\wc\sqrt{\fgap/3}$ to Dinh's tapered geometry yields a
$\dinhBandLo$--$\dinhBandHi\um$ cutoff band that lies inside the measured $(\dinhBracketLo,\dinhBracketHi)\um$
rigid-bead bracket. Because Dinh's \emph{bracket device} (the
$\fgap^{*}=1.04\times10^{-4}$ design, the only one with rigid-bead data) operates
\emph{at or below the knee} at its initial gap fraction (held there, the
self-consistent cutoff sits at the sub-steric foot, \emph{beneath} the
bracket; under the schedule Dinh describes but does not print, read with the
Xia scaling Dinh 2024 cites, in its corrected form (C5), it rises into the bracket, so
what is shown is the containment of the plateau-anchored band, not of the device's
own cutoff),
and the bracket is $4.7\times$ wide with soft
edges, this is a
\emph{low-severity, anchor-contingent consistency check, not a discriminating
prediction}. The appropriate anchors for future finite-particle work are
plateau devices, not Dinh.\footnote{Naming caution: this argument node
\textbf{C7} is distinct from \emph{Result~7} of the algebraic Appendix~A (the
cross-width scaling $\fgap\propto\wc^{-2}$, which backs node \textbf{C5});
Appendix~A contains no empirical held-out node.}

% ----------------------- ARGUMENT GRAPH 1 (structure) -------
\begin{figure}[t]
\centering
\begin{tikzpicture}[
  font=\scriptsize,
  prem/.style={draw, rounded corners, align=center, inner sep=3pt,
               text width=5.3cm, minimum height=13mm},
  claim/.style={draw, rounded corners, align=center, inner sep=4pt,
                text width=8.4cm, minimum height=10mm},
  cav/.style={draw, dashed, rounded corners, align=center, inner sep=3pt,
              text width=4.5cm, minimum height=14mm},
  concl/.style={draw, rounded corners, align=center, inner sep=4pt,
                text width=8.8cm, minimum height=11mm},
  up/.style={-{Stealth[length=2mm]}, thick},
  cl/.style={-{Stealth[length=1.8mm]}, dashed},
]
% --- premises ---
\node[prem] (calc) at (-3.6,0)
  {\textbf{Computation (P-C7.1).}\ Hill $\kapp$-curve $+$ skim law on Dinh's
   geometry ($\wc\!:\!75\!\to\!52\um$, initial $\fgap^{*}{=}1.04\!\times\!10^{-4}$)
   $\Rightarrow$ band $[\dinhBandLo,\dinhBandHi]\um$};
\node[prem] (evid) at (3.6,0)
  {\textbf{Evidence (P-C7.2).}\ $1\um$ tracks fluid ($54\%\!\approx\!52\%$);
   $4.7\um$ retained ($1.8$--$6.8\%$ vs ${\sim}50\%$) $\Rightarrow$ bracket
   $(\dinhBracketLo,\dinhBracketHi)\um$};
% --- consistency claim ---
\node[claim] (cl1) at (0,-2.35)
  {\textbf{C7.}\ band $[\dinhBandLo,\dinhBandHi]\um\ \subset\ $ bracket $(\dinhBracketLo,\dinhBracketHi)\um$:
   held-out consistency on a device whose data was not fitted\\
   (held out, but same lineage)};
% --- caveats ---
\node[cav] (cq4) at (-5.0,-4.6)
  {\textbf{Low severity (CQ4).}\ bracket $4.7\times$ wide, over-extraction
   tolerance $7.3\times$ wide, soft edges; almost any admissible parameter
   lands inside};
\node[cav] (cq3) at (5.0,-4.6)
  {\textbf{Anchor-contingent (CQ3).}\ at the initial $\fgap^{*}$ the
   self-consistent fixed point ${\approx}0.4$--$0.6\um$ sits \emph{below} the
   bracket, so the band is plateau-anchored, not the bracket device's cutoff};
% --- conclusion ---
\node[concl] (cc) at (0,-6.85)
  {\textbf{Warranted only as a low-severity, anchor-contingent corroboration},
   not a test that could have discriminated against the law. Clean anchors
   $=$ plateau devices (Xia OR-7 and the RS-8/RS-9 reference triples, published
   $\fgap=5.8$--$8.6\times10^{-4}$ at $\wc(\mathrm{ref})=100\um$), not Dinh.};
% --- edges ---
\draw[up] (calc) -- (cl1);
\draw[up] (evid) -- (cl1);
\draw[up] (cl1) -- (cc);
\draw[cl] (cq4) -- (cc);
\draw[cl] (cq3) -- (cc);
\end{tikzpicture}
\caption{Argument graph for C7. The computed band (Computation, P-C7.1) and the
measured bracket (Evidence, P-C7.2) agree: the band is contained in the
bracket. Two dashed caveats downgrade what that agreement is worth: the test is
low-severity (CQ4, a $4.7\times$ soft bracket) and anchor-contingent (CQ3: at the
device's initial gap fraction the self-consistent cutoff lies below the bracket,
so the band is pinned to the Gifford plateau rather than iterated to the bracket
device's own fixed point). The
conclusion
is corroboration, not a discriminating test. Border language follows the
overview: single $=$ warranted/corroborating, dashed $=$ open/adjudicated
caveat; dashed arrows carry the downgrade.}
\label{fig:c7-argmap}
\end{figure}

% ----------------------- CLAIM ------------------------------
\field{Claim.} The $\kapp(\dwc)$ collapse curve, transferred to Dinh's geometry
through the skim law $\delta=\wc\sqrt{\fgap/3}$, places the bracket device's cutoff in the band
$[\dinhBandLo,\dinhBandHi]\um$, which is contained in the measured rigid-bead bracket
$(\dinhBracketLo,\dinhBracketHi)\um$.

% ----------------------- SCHEME -----------------------------
\field{Scheme.} A deterministic computation for the band itself (no critical
questions), combined with Argument from Evidence to a Hypothesis (verification
form) for the held-out match (WRM~2008). The governing critical question of the verification half is
severity: \emph{could the band have landed outside the bracket?} (CQ4).

% ----------------------- PREMISES ---------------------------
\field{Premises.}
\begin{description}

\item[P-C7.1 (computed band).]
Fitting a Hill form $\kapp(u)=\kappa_{\max}\,u^{n}/(u^{n}+u_{1/2}^{\,n})$
($\kappa_{\max}=\hillKmax$, $u_{1/2}=\hillUhalf$, $n=\hillN$; $u=\dwc$) to the $\kapp$
collapse, and evaluating $\delta=\wc\sqrt{\fgap/3}$ at Dinh's geometry
($\wc:75\to52\um$, initial $\fgap^{*}=1.04\times10^{-4}$, gaps $19\um$; from ``\emph{All filtration gaps in
the linear section were 19 $\mu$m-wide. The middle channel narrowed from 75
$\mu$m near the inlet to 52 $\mu$m near the outlet in all designs. Depth of all
channels was 140 $\mu$m}'' \cite[p.~917, Fig.~2 caption]{b5}), gives, holding $\fgap$
at that value, $\delta=0.442\um$
($\wc=75$) and $0.306\um$ ($\wc=52$), and the band $[\dinhBandLo,\dinhBandHi]\um$
(regression-stamped, and reproduced to the printed precision by the reference
code \texttt{cif\_cutoff.py}).

\emph{The band's width is an ambiguity, not a gradient.} A tapered device is
designed to hold one cutoff, not a range of them; Xia ramps $\fgap$ precisely so
the lamina is taper-invariant (C5; \cite[p.~9 (Sci Rep PDF pagination)]{b2},
corrected in \cite[p.~1]{b3}). Dinh states the same design rule for its own
devices: once the initial values are chosen, ``\emph{the values of $\fgap(i)$ are
iteratively adjusted with each subsequent gap $i$ by incrementally increasing the
fluidic resistance of the middle channel and/or decreasing the fluidic resistance
of the side channels to maintain the desired size cut-off throughout the
device}'' \cite[p.~916]{b5}, citing Xia's 2016 paper for the method (its
ref.~28). What the paper does not print is the schedule itself, only the initial
$\fgap^{*}=1.04\times10^{-4}$ for this device, one of three initial values its
designs span \cite[p.~917, Fig.~2 caption]{b5}. The band is computed holding
$\fgap$ at that initial value, which makes $\delta\propto\wc$ and spreads the
prediction across the taper. Of the band's $\dinhBandWidth\times$ width,
$75/52=1.44\times$ is that spread and only $\dinhKneeRatio\times$ is the near-knee
$\kapp$ ambiguity, where the fitted prefactor has ramped off the plateau to a local
${\approx}\dinhKappaLocal$. Under the Xia-scaled reading, which is the one Dinh
describes, the lamina is the same at every gap, so the cutoff at the plateau
prefactor is a single value and the band collapses to $[\dinhCollapseLo,\,\dinhBandHi]\um$,
the $\kapp$ ambiguity alone; the iterated cutoff is not single-valued (end of the
CQ3 backing). The wider band is retained as the conservative
envelope, and it is the band the graph records: it contains the Xia-scaled
reading, and taking the narrower one would assume that a schedule the paper
describes but does not print meets its stated intent exactly. \emph{The rule is
stated; the schedule is not, and that is why the band is wide.}

\emph{The channel depth
($h=140\um$) is a property of the device but does \textbf{not} enter the band:}
the skim law is depth-free and the $\kapp$ curve is fit on the parabolic
$\delta$; depth enters only the interpretation (CQ2). \prov{code: cif\_cutoff.py
B5 (kcurve); skim.delta\_leading; glossary \xref{app:units_selval}{A.7}}

\item[P-C7.2 (Dinh 2024, {\cite[ESI, Fig.~S1(A) and (B)]{b5}}; bead data, verbatim).]
For the same device ($\fgap^{*}=1.04\times10^{-4}$, $30\,\mu$L/min): $1\um$ beads
give ``\emph{fluid volume recovery in filtrate was 52\% and bead recovery in
filtrate was 54\%}'' (tracer-like, \emph{below} cutoff, which is how Dinh 2024 reads the
same panel: for $a<w$ the split should ``\emph{approximate that of the fluid volume}'',
with Fig.~S1A given as the example \cite[p.~917]{b5}); $4.7\um$ beads give
``\emph{fluid volume recovery in filtrate: 50\%, bead recovery in filtrate:
1.8\%}'' at $106\,\mu$L$^{-1}$, and stay far below the fluid baseline across the
whole sweep ($1.8/3.5/6.8\%$ bead against $50/52/48\%$ fluid at
$106/799/1714\,\mu$L$^{-1}$; \emph{retained}, \emph{above} cutoff).
$\Rightarrow$ the measured cutoff lies in $(\dinhBracketLo,\dinhBracketHi)\um$.
\prov{Dinh 2024, ESI Fig.~S1(A) and (B)}

\end{description}

% ----------------------- CRITICAL QUESTIONS -----------------
\field{Critical questions.}
\begin{description}

\item[CQ1: Is the calculation correct?]
Yes. The executable reference reproduces the band to the printed precision: its
regression registry asserts \texttt{Dinh band low/high} $=\dinhBandLo/\dinhBandHi$ to
the two decimals printed. $\rightarrow$
\emph{answered: yes.}

\item[CQ2: Are there interfering factors?]
Two, of unequal magnitude. \emph{(a) Per-gap over-extraction (dominant).} If the
true per-gap fraction is $\feff=c\,\fnom$, then $\delta\propto\sqrt{c}$ and the
band scales by $\sqrt{c_{\mathrm{Dinh}}/c_{\mathrm{Gifford}}}$. \emph{(b) Depth
(secondary).} The fitted $\kapp$ absorbs Gifford's finite-depth factor
$R(\alpha_{\mathrm{Gifford}})$; transferring it to Dinh's higher aspect ratio
carries a mismatch $R(\alpha_{\mathrm{Dinh}})/R(\alpha_{\mathrm{Gifford}})$ of
$1.09$--$1.21$ across Dinh's two stations ($\alpha=1.87$ and $2.69$) and the
calibration array's three widths ($\alpha=0.53$, $0.64$, $0.80$), and (because
the band uses the depth-free parabolic $\delta$) this
does not enter the band formula, only its interpretation. The interval over which
the band remains inside $(\dinhBracketLo,\dinhBracketHi)$ is $c_{\mathrm{Dinh}}/c_{\mathrm{Gifford}}\in
[0.64,4.69]$, where $0.64=(\dinhBracketLo/\dinhBandLo)^{2}$ and $4.69=(\dinhBracketHi/\dinhBandHi)^{2}$: \emph{this is
a tolerance (sensitivity) interval, not an independent plausibility range.} The
honest robustness statement is therefore that the band exits the bracket only if
the over-extraction ratio falls below $0.64\times$ or above $4.69\times$. Note
what does \emph{not} establish that it will not: bounding each device separately,
below by $c\gtrsim1$ and above by Gifford's stated per-gap bound
($\fgap<5\times10^{-4}$, \cite[p.~4498]{b1}), caps the Dinh factor only at
$5\times10^{-4}/1.04\times10^{-4}=4.8$, which is \emph{outside} the window's
$4.69$ edge, and leaves the ratio unbounded below $0.64$ for the same reason.
What makes the ratio plausibly near unity is not those bounds but that the two
devices' over-extraction has the same origin, the local gap-to-channel
resistance ratio, at comparable gap widths ($19\um$ against Gifford's
${\sim}20\um$) and comparable channel widths. That is a \emph{tolerance}
argument resting on a similarity premise, not an independent confirmation. Finally, the concentration-dependent $4.7\um$
leakage (rising $1.8\to6.8\%$ as bead concentration rises $106\to1714\,\mu$L$^{-1}$)
is crowding (bead--bead interaction), not a cutoff shift. Dinh 2024 reads the same
sweep the same way: ``\emph{As the particle concentration increases, the
interactions between particles also increase and, therefore, the recovery of large
particles in the filtrate should also increase}'' \cite[p.~917]{b5}.
$\rightarrow$ \emph{answered: bounded.}

\item[CQ3: Is the band the self-consistent cutoff? (anchor dependence).]
No. Because the calibration curve returns $\kapp$ \emph{as a function of the
cutoff}, the operating law is \emph{implicit}, $\dcrit=\kapp(\dcrit/\wc)\,\delta$,
and the cutoff is the fixed point of $u=\sqrt{\fgap/3}\,\kapp(u)$. Solved
honestly at the bracket device's initial gap fraction, held constant as the band
holds it, the only self-consistent value is the sub-steric foot
($\dcrit\approx0.41$--$0.59\um$), \emph{below} the bracket and inconsistent with
the $1\um$ bead tracking the fluid. On that reading the reported band lands in the
bracket only because it is read off the Gifford plateau anchor: its lower edge takes the
near-knee curve value $\kappa\approx\dinhKappaLocal$ at the outlet station and its upper edge
the plateau anchor $\kappa=\plateauKappa$ at the inlet, the curve value being the same at
both stations because $\kapp$ is evaluated at $\kplat\sqrt{\fgap/3}$, which does
not depend on $\wc$. It is not iterated to
that device's own fixed point. Along the linear taper, where Dinh's schedule raises
the gap fraction, the self-consistent value rises into the bracket instead, when the
rise is read with the Xia scaling Dinh 2024 cites, in its corrected form (C5); on
either reading the band is the plateau-anchored value and not the iterated one,
which is what this question records. The diagnostic, including why the threshold
is specific to that gap fraction and what the schedule changes, is worked out in the
subsection below (Fig.~\ref{fig:c7-fixedpoint}). $\rightarrow$ \emph{disclosed;
conclusion downgraded to anchor-contingent.}

\item[CQ4: How severe is the test? (could it have failed?).]
Barely. The bracket is $4.7\times$ wide and the over-extraction tolerance is
$7.3\times$ wide, so almost any admissible parameterisation lands inside; the
bracket edges are moreover \emph{soft} (bead-population averages, not sharp
cutoffs). The robustness asserted in CQ2 and the weakness asserted here are two
readings of the \emph{same} insensitivity. Passing this test is weak
corroboration of the law. $\rightarrow$ \emph{answered: low severity.}

\end{description}

% ----------------------- DESIGN CHOICE / RESIDUAL / STATUS --
\field{Design choice (promoted from caveat).} The clean anchors for the knee are
\emph{plateau} devices (Gifford's larger beads, and Xia's OR-7 with the RS-8 and
RS-9 reference triples, all at the $100\um$ reference width), where a stable
cutoff fixed point exists on the plateau. The RS devices hold that width only as a
reference: they scale the gap fraction as $(100\um/\wc)^{2}$ (eqn~(1) as
corrected; C5) along a wider taper (RS-7's runs $300\to150\um$
\cite[p.~10 (Sci Rep PDF pagination)]{b2}; RS-8's and RS-9's starting widths are not
printed), so
on RS-7's geometry only the narrower sections are plateau. Any Track-B
finite-particle
study should run on plateau geometries, not on Dinh. (C3b's CQ2 downgrades Xia's
$\kapp\approx\xiaKappaMean$ to a lead, on grounds of circularity and non-diagnosticity.
That does not bear here: what places these anchors on the plateau is their
\emph{published} $\fgap = 5.8$--$8.6\times10^{-4}$ at $\wcref=100\um$
\cite[pp.~6 and~10 (Sci Rep PDF pagination)]{b2}, read straight off the paper,
and not the prefactor computed from them.)

\field{Residual risk.} None \emph{open}: the graph records
\texttt{open\_exceptions: []} for this node, and every critical question above is
answered. The mapping is one-to-one: CQ2 and CQ3 are the node's two Walton
exceptions (\texttt{doubly\_sensitive\_depth\_and\_feff} and
\texttt{anchor\_contingent\_fixed\_point\_below\_bracket}), CQ4 is its assumption
(\texttt{sufficient\_correlation}), and CQ1 belongs to the computation half
(\texttt{c7\_dinh\_computation}), which carries no critical question at all. What remains is not risk but \emph{weakness}, and it is the weakness the
status word already carries: anchor-contingency (CQ3), double confounding with
$\feff$ dominant and depth secondary (CQ2), and low severity against a
$4.7\times$ soft bracket (CQ4). Three answered questions whose answers bound how
much this band can be asked to carry, which is why it is a band and not a test.

\field{Status.} \textbf{Warranted only as a low-severity, anchor-contingent
consistency band}: a corroboration of the cutoff law on a held-out device, not a
test that could have discriminated against it.

% ===========================================================================
\subsection*{Why the band is anchor-contingent: the implicit cutoff (backing for CQ3)}

The operating law fixes the cutoff as $\dcrit=\kapp\,\delta$, and for the
bracket device at its initial gap fraction $\delta=\wc\sqrt{\fgap/3}$ is known
($0.31$--$0.44\um$). The subtlety is that
$\kapp$ is not a free constant: the calibration curve was built by giving each
Gifford device its own apparent prefactor $\kapp=\dcrit/\delta$ and plotting it
against \emph{that device's own} normalised cutoff $\dcrit/\wc$. The curve
therefore returns $\kapp$ \emph{as a function of the cutoff}, and the operating
law becomes \emph{implicit}: the cutoff appears on both sides,
$\dcrit=\kapp(\dcrit/\wc)\,\delta$. Writing $u=\dcrit/\wc$ and dividing by $\wc$
turns this into a one-line fixed-point equation whose root is the predicted
cutoff,
\[
  u \;=\; \sqrt{\tfrac{\fgap}{3}}\;\kapp(u)\;=:\;g(u),
  \qquad \sqrt{\tfrac{\fgap}{3}}=0.00589 .
\]
On the plateau this has a clean stable root; the question is whether the bracket
device does at its initial gap fraction, held constant along the taper as the band
holds it.

It does not. Starting the natural way, the plateau value $\kapp=\plateauKappa$ predicts
$\dwc=0.029$; but $\plateauKappa$ is the mean over the plateau bin $\dwc\ge\plateauUwc$,
and the fitted curve itself does not reach it until $\dwc\approx0.051$. At
$0.029$ the curve returns $\dinhKappaLocal$, which predicts a \emph{smaller} cutoff
($\dwc\approx0.024$), where the curve is smaller still ($\approx3.7$), so the
estimate drifts down at every step ($0.029\to0.024\to0.022\to\cdots$). The only
self-consistent value is the sub-steric foot ($\kapp\to\footKappa$),
$u^{*}\approx0.0079$, i.e.\ $\dcrit\approx0.41$--$0.59\um$ (the bare Hill fit,
having no foot, drives the iteration to $0$ instead). Either way the
self-consistent cutoff lands \emph{below} the measured bracket and would predict
the $1\um$ bead retained, contradicting its observed fluid-tracking
($54\%\approx52\%$).

This is specific to the \emph{bracket} device's initial gap fraction. A
self-consistent cutoff \emph{above the foot} exists only once the rising curve
clears the knee, that is, only once $g$ touches the diagonal:
$\sqrt{\fgap/3}\,\kapp(u)=u$ together with $\sqrt{\fgap/3}\,\kapp'(u)=1$. A
tangency is a derivative condition, so it is solved on the full-precision fit.
It puts the threshold at $\fgap=\dinhSaddleFgap$, with the tangency itself on
the rising curve at $\dwc\approx0.0145$. Read as the plateau prediction at that
gap fraction, $\kplat\sqrt{\fgap/3}$, the threshold is $\dwc=\dinhSaddleU$, which
numerically coincides with the $\dwc\approx0.03$ knee C3 and C3c already name,
though it is a starting value for the iteration and the knee is a place on the
curve. The two independent
re-digitisations of Fig.~3 give $1.042$ and $1.046\times10^{-4}$. Xia's OR-7
($5.8\times10^{-4}$, plateau prediction $\dwc=0.068$) sits far above it.
\textbf{At its initial gap fraction the bracket device sits on it}: its plateau prediction is
$\dwc=0.02893$ against a threshold of $\dinhSaddleU$, so it is below on all three
digitisations, but by less than $0.6\%$ on every one. That is the quantitative
content of ``the bracket device operates below the knee,'' and it is a knife edge rather than
a margin. What does \emph{not} turn on the margin is the conclusion CQ3 records,
that the band is read off the plateau anchor rather than iterated to a
self-consistent cutoff: that holds on either side of the knee. What does turn on
it is the descent to the foot, which holds on all three digitisations and would
not if the device sat a fraction above the threshold. What the knife edge does
forbid is reading the threshold as a comfortable separation, and it puts a near-fixed
point in the descent at $\dwc\approx0.0145$ ($\dcrit\approx0.75$--$1.09\um$),
where $g$ comes within $4\times10^{-6}$ of the diagonal. That near-fixed point is
the tangency point itself, where a stable root would appear were the gap fraction a
hair higher, so the three steps
displayed above understate how slowly the estimate falls to the foot. (The Gifford \emph{calibration}
array deliberately runs down to $\fgap=6.4\times10^{-5}$, returning measured
thresholds $\fgap^{*}$ as low as ${\sim}10^{-4}$ \cite[p.~4500]{b1}, which is
exactly why the curve \emph{has} a foot to fall onto; its plateau-anchoring
enrichment devices sit at $\fgap\ge2.2\times10^{-4}$
\cite[p.~4501, Fig.~5 caption]{b1}.) At that gap fraction the reported band
$[\dinhBandLo,\dinhBandHi]\um$ therefore lands inside the bracket only because it is read off
at the plateau (the first step above), not iterated to the bracket device's own
sub-bracket fixed point. That dependence on the anchor, not the containment itself, is what
CQ3 records, and it is why the clean anchors are the plateau devices, not Dinh.

\emph{Dinh's other two designs are on the far side of the threshold.} The same caption
reports $\fgap^{*}=1.12$ and $1.28\times10^{-4}$ beside the
$1.04\times10^{-4}$ of the bracket device \cite[p.~917, Fig.~2 caption]{b5}, and
both clear the threshold: each acquires a stable fixed point, at
$\dcrit=\sibALo$--$\sibAHi\um$ and $\sibBLo$--$\sibBHi\um$ respectively, inside the measured
bracket, and the iteration started at the plateau converges to it rather than
drifting to the foot. Both roots still lie below the plateau edge $\plateauUwc$, so
neither design is a clean anchor either: a stable fixed point is not a plateau.
The anchor-contingency recorded here is therefore a property of the one device
the bead bracket was measured on, at its initial gap fraction, not of the paper.
Nor can those two stand in for it: the rigid-bead runs were made on the
$1.04\times10^{-4}$ device alone \cite[ESI, Fig.~S1]{b5}.

\emph{Held constant is the band's reading, not Dinh's schedule.} Everything above
holds $\fgap$ at $\fgap^{*}$ along the whole taper, as the band does. Dinh's
schedule does not: the gap fraction is adjusted gap by gap ``\emph{to maintain the
desired size cut-off throughout the device}'' (P-C7.1), and in the linear section
that is done by ``\emph{progressively increasing the width of the side channels and
decreasing}'' that of the middle channel \cite[p.~916]{b5}. Under the Xia
scaling that design sentence cites, in its corrected form (C5),
$\fgap(i)=\fgap^{*}\,(75\um/\wc)^{2}$, the gap
fraction clears the threshold once the middle channel has narrowed to
$\wc\approx74.98\um$, and from there the iteration started at the plateau settles on
a stable root at $\dcrit\approx1.09$--$2.02\um$ along the linear section, rising
from the tangency point to the outlet and lying inside the bracket. The descent to
the foot, and with it the contradiction with the $1\um$ bead, is then confined to
the stretch before the middle channel begins to narrow, where $\fgap$ still equals
$\fgap^{*}$; how much it weighs depends on that stretch's share of the extracted
fluid, which the paper does not print. What the schedule does not change is what
CQ3 records: on either
reading the band is the plateau-anchored value rather than the iterated one, and
the paper prints the rule but not the $\fgap(i)$ that would let the iterated value
be computed from Dinh alone.

% ----------------------- ARGUMENT GRAPH 2 (interval) --------
\begin{figure}[h]
\centering
\begin{tikzpicture}[font=\scriptsize, x=2.5cm, y=1cm]
% --- axis 0..5 um ---
\draw[-{Stealth[length=1.6mm]}] (-0.05,0) -- (5.25,0) node[right]{$\dcrit\ (\mu\mathrm{m})$};
\foreach \x in {0,1,2,3,4,5} \draw (\x,0.06) -- (\x,-0.06) node[below,font=\tiny]{\x};
% --- measured bracket (1.0, 4.7): open span above axis, with end caps ---
\draw[thick] (1.0,0.95) -- (4.7,0.95);
\draw[thick] (1.0,0.78) -- (1.0,1.12);
\draw[thick] (4.7,0.78) -- (4.7,1.12);
\node[above, font=\tiny] at (2.85,1.22)
  {measured rigid-bead bracket $(\dinhBracketLo,\dinhBracketHi)\um$: $4.7\times$ wide, soft edges
   ($1\um$ tracks fluid $\to$ $4.7\um$ retained)};
% --- computed band [1.25, 2.17]: filled, nested, straddling the bracket line ---
\fill[black!18] (\dinhBandLo,0.78) rectangle (\dinhBandHi,1.12);
\draw (\dinhBandLo,0.78) rectangle (\dinhBandHi,1.12);
\node[below, font=\tiny] at (1.71,0.66) {computed band $[\dinhBandLo,\dinhBandHi]\um$};
% --- fixed point ~0.4-0.6 um: marker below axis, outside (below) the bracket ---
\fill[black!55] (0.5,0) circle (1.7pt);
\draw[densely dotted] (0.5,0) -- (0.5,-0.42);
\node[below, font=\tiny, align=center] at (0.5,-0.44)
  {self-consistent\\ fixed pt.\ ${\approx}0.4$--$0.6\um$};
\end{tikzpicture}
\caption{Interval diagnostic for C7. The computed band $[\dinhBandLo,\dinhBandHi]\um$ (shaded)
is nested inside the measured rigid-bead bracket $(\dinhBracketLo,\dinhBracketHi)\um$, so the held-out
check passes, but the self-consistent fixed point at the device's initial gap
fraction ($\approx0.4$--$0.6\um$, the filled marker) sits \emph{below} the bracket
entirely. At that gap fraction the band agrees with the
bracket only because it is anchored to the Gifford plateau, not iterated to the
bracket device's own (sub-bracket) cutoff; and the bracket is $4.7\times$ wide with soft
($1\um$-tracking vs $4.7\um$-retained) edges. All three facts (containment, the
sub-bracket fixed point, and low severity) are visible at once.}
\label{fig:c7-fixedpoint}
\end{figure}


\clearpage
\section*{C4 --- Separation lemma: $F\indep\fgap$ (the outlet identity carries no
cutoff information)}

\noindent\textbf{Summary.} The device-level filtrate fraction $F$, the share of
inlet flow leaving through the two side outlets, is fixed by the outlet
resistance ratio alone, $F=2\Rc/(\Rs+2\Rc)$, with the per-gap schedule, the gap
count $N$, the taper, and the cutoff-law scaling exponent all provably absent.
This is \emph{proven} symbolically, \emph{validated} against 3-D COMSOL on two
Dinh variants (the 1-D residual closing to within $1\sigma$ of the published
values), and \emph{recovered} twice: as the emergent behaviour of the full
resistor network, and by the iterated 1-D recursion, including the strongest
case, a variable-$\fgap$ schedule where $\fgap$
changes at every gap yet $F$ is unmoved to within $2\times10^{-5}$. The same fraction can also be written as
a cumulative product, $F=1-\prod_{i}(1-2\fgap(i))\approx1-e^{-2N\fgap}$, which
\emph{is} $\fgap$- and $N$-dependent; the two forms agree only because a given
schedule fixes the outlet geometry. The lemma's content is therefore precise: at
fixed outlet geometry, $F$ is independent of \emph{how} the device reached it,
and $F$ carries no information about the local per-gap fraction (or the
over-extraction $c$) that sets the cutoff. Recovering $\kapp$ or $c$ from $F$ is a
category error: the inversion this project retracted.

% ----------------------- ARGUMENT GRAPH ---------------------
\begin{figure}[t]
\centering
\begin{tikzpicture}[
  font=\scriptsize,
  ante/.style={draw, rounded corners, align=center, inner sep=4pt,
               text width=10.2cm, minimum height=9mm},
  ded/.style={draw, double, double distance=0.7pt, rounded corners,
              align=center, inner sep=4pt, text width=8.6cm, minimum height=11mm},
  sup/.style={draw, rounded corners, align=center, inner sep=3pt,
              text width=4.4cm, minimum height=14mm},
  sep/.style={draw, rounded corners, align=center, inner sep=4pt,
              text width=10.6cm, minimum height=11mm},
  down/.style={-{Stealth[length=2mm]}, thick},
  up/.style={-{Stealth[length=2mm]}, thick},
  blk/.style={-{Stealth[length=2mm]}, thick, dashed},
  nt/.style={font=\scriptsize\itshape, text=black!60, align=center},
]
% --- antecedents ---
\node[ante] (an) at (0,0)
  {\textbf{Shared CIF architecture (premises).}\ $\Qs(0)=0$ (single inlet)
   $\;\cdot\;$ device mass conservation $\;\cdot\;$ outlet pressure
   equilibration (PE)};
% --- deductive identity ---
\node[ded] (fid) at (0,-1.95)
  {\textbf{Current divider (deductive)}\ $\Rightarrow$\ $F=2\Rc/(\Rs+2\Rc)$:
   fixed by the outlet ratio alone; \ $\fgap$, $N$, taper, exponent provably
   absent};
% --- corroboration ---
\node[sup] (cm) at (-4.6,-4.25)
  {\textbf{Validated (3-D reality)}\\ COMSOL within $1\sigma$:\\
   $F_{\mathrm{closed}}\,0.56/0.72\to F_{\mathrm{COMSOL}}\,0.51/0.63$};
\node[sup] (nw) at (4.6,-4.25)
  {\textbf{Verified (network, recursion)}\\ network-emergent $\netEmergentPct\%$;\\ recursion,
   variable-$\fgap$ $|\Delta F|{=}2{\times}10^{-5}$};
% --- separation payoff ---
\node[sep] (sp) at (0,-6.55)
  {\textbf{Separation (proven): $F\indep$ cutoff.}\ The cumulative form
   $1-\prod(1-2\fgap)$ recovers only the product $N\fgap$ ($N$ unpublished),
   never the local over-extraction $c$. Reading $\kapp$ or $c$ out of $F$ (the
   retracted ``$0.3516$'') is a \emph{category error}.};
% --- edges ---
\draw[down] (an) -- (fid);
\draw[up] (cm) -- (fid);
\draw[up] (nw) -- (fid);
\draw[blk] (fid) -- (sp) node[midway,right,nt]{licenses};
\end{tikzpicture}
\caption{Argument graph for C4. The outlet identity is \emph{deduced} (doubled
border) from three shared-architecture premises, \emph{validated} against 3-D
reality by COMSOL, and \emph{verified} twice: as the emergent behaviour of the
full network, and against the iterated 1-D recursion, most stringently on the
variable-$\fgap$ schedule. Its payoff is the
separation: because $F$ sees only the outlet ratio (or, written as a cumulative
product, only the product $N\fgap$), it carries no information about the local
per-gap over-extraction $c$ that sets the cutoff. Border language follows the
overview: doubled $=$ deductive, single $=$ warranted; the dashed arrow marks the
licensed boundary (what the lemma permits and forbids).}
\label{fig:c4-argmap}
\end{figure}

% ----------------------- CLAIM ------------------------------
\field{Claim.} At fixed outlet geometry $(w_c,w_s)$, $F=2\Rc/(\Rs+2\Rc)$ is
fixed by the outlet resistance ratio alone. It is independent of the per-gap schedule $\{\fgap(i)\}$, the gap count
$N$, the taper, the cutoff-law scaling exponent, and the inlet flow rate. That
holds exactly under the three premises of P-C4.1; a finite tapered network realizes
it to within a residual of the order of the gap fraction, which vanishes as the
ladder is refined (P-C4.3). Up to that residual, $F$ carries no information about
the local per-gap fraction, or the over-extraction $c$, that sets the cutoff. Recovering $\kapp$ (or $c$) from $F$ is a category
error.

% ----------------------- SCHEME -----------------------------
\field{Scheme.} Argument from Established Rule (deductive): the two-resistor
current-divider rule (hydraulic--electrical analogy \cite{b13,b14}) applied to the CIF outlet
under its shared architectural premises (WRM~2008). The lineage applies the rule
at the outlet itself: for devices with $450\um$ final side channels and a
$100\um$-wide, $125\um$-deep centre channel, Gifford puts the ratio of the two
side channels' cumulative output flow to the centre channel's at theoretically
${\sim}$23~:~1, calculated from the duct-resistance series (his eqn~(5)),
``\emph{provided the hydrostatic pressures in the output collection areas are
kept equal}'' \cite[p.~4500]{b1}: a flow split fixed by the outlet resistances
alone, at equal outlet pressures.

% ----------------------- PREMISES ---------------------------
\field{Premises.}
\begin{description}

\item[P-C4.1 (the rule and its antecedents).]
The current-divider form follows deductively from three premises shared across
the CIF lineage: (1) the boundary condition $\Qs(0)=0$ (sample enters the centre
channel only, a single inlet); (2) device mass conservation $\Qc(0)=\Qc(N)+
2\Qs(N)$; (3) outlet pressure equilibration $\Qs(N)\Rs=\Qc(N)\Rc$. With
$F\equiv2\Qs(N)/[\Qc(0)+2\Qs(0)]=2\Qs(N)/\Qc(0)$ (using premise~1), three
algebraic steps give $F=2\Rc/(\Rs+2\Rc)$; equivalently, parallel-combining the
two side branches ($\Rs\parallel\Rs=\Rs/2$) reaches the identical result from
the same premises by an independent algebraic route. The two are the same
physics in two notations, the hydraulic--electrical analogy mapping one onto the
other, so their agreement is a consistency check rather than independent
evidence; its value is that the current-divider form makes $F$'s dependence on
outlet geometry alone immediately visible.
The limits are physical ($F\to1$ as $\Rs\to0$, $F\to0$ as
$\Rs\to\infty$, $\partial F/\partial\Rs<0$), and $\fgap$ is symbolically absent.
\prov{Routt 2026 \S II; code: network.F, network.verify\_symbolic, which asserts
\texttt{f\_gap not in network.F.free\_symbols} (separation lemma C4)}

\item[P-C4.2 (3-D physical validation, with both distances stated).]
The closed form predicts $F_{\mathrm{closed}}=\FclosedVone$ (V1)$/\FclosedVthree$ (V3), above
the published $\FexpVone\pm\FsdVone/\FexpVthree\pm\FsdVthree$ \cite[p.~919, Fig.~4 caption]{b5}
by $\FexcessPctVOne\%$ (V1) and $\FexcessPctVThree\%$ (V3);
3-D outlet-bifurcation physics closes the residual
($F_{\mathrm{COMSOL}}=\FcomsolVone/\FcomsolVthree$), and the \emph{1-D-vs-3-D} residual it
closes is $\residPctVOne\%$ (V1) and $\residPctVThree\%$ (V3), a different comparison from the one
above, and the two should not be quoted for each other. The glossary requires both $\sigma$-scales
wherever one is given, so both are here. On the raw-$\mathrm{SD}$ scale the 3-D
solve is within $1\sigma$ ($\EsdComsolVOne/\EsdComsolVThree$); on the stricter $\mathrm{SEM}$ scale it
is within $2\sigma$ ($\EsemComsolVOne/\EsemComsolVThree$). The \emph{calibration-free} 1-D closed form
sits $\EsdClosedVOne\sigma$ (V1) and $\EsdClosedVThree\sigma$ (V3) above the measurement:
consistent within $2\sigma$, but \textbf{not a match}, and the claim this node
makes is that the yield law needs no device-specific fitting, not that the 1-D
form reproduces the measurement. The residual is itself an outlet-geometry effect
(the V1/V3 asymmetry $\alpha_{\mathrm{3D}}=\alphaThreeDVone$ vs $\alphaThreeDVthree$), so the correction
introduces no $\fgap$ coupling. The same caption bears on the Claim's flow-rate
clause directly: Dinh run the V1 device at two inlet flow rates, and
``\emph{Volume recovery in filtrate was 50.5 $\pm$ 3.3\% for 30 $\mu$L\,min$^{-1}$
and 49.9 $\pm$ 1.8\% for 90 $\mu$L\,min$^{-1}$}''
\cite[p.~919, Fig.~4 caption]{b5}. A threefold change in flow moves the yield by
less than either standard deviation. \prov{Routt 2026 \S IV-B; COMSOL Multiphysics 6.3;
glossary \xref{app:units_validation}{A.6} validation block}

\item[P-C4.3 (network-emergent and schedule-independent).]
The full Kirchhoff solve's emergent $F$ equals the symbolic identity at the
outlet to $\netEmergentPct\%$ across ramp shapes (bridge N3); the closed form matches the
iterated recursion to $|\Delta F|\le\dFrecurConst$ across four
constant-$\fgap$ devices and to $|\Delta F|=\dFrecurVar$ on the
variable-$\fgap$ Xia~RS-7 case, the most stringent test, since $\fgap$ enters
every recursion step yet drops out of the closed form. This is the direct
empirical content of ``$F\indep\fgap$'': varying $\fgap$ at every gap, at a fixed
outlet, leaves $F$ at the closed form to the figures above. The separation
itself is symbolically exact, and the full network of bridge N3 realizes it to
$\netEmergentPct\%$; the recursion figures above are a different comparison, against a 1-D
recursion that assumes equilibration at every gap \cite[p.~4497]{b1}. The
network's departure is the taper's own: an untapered ladder realizes the identity
to rounding, while on the linear tapers probed, Dinh's own geometry among them, the
network holds an outlet pressure imbalance that moves $F$ by between about a quarter
of and about four times the gap fraction, depending on the taper and on the thickness
of the wall each gap crosses, which the papers do not state; it falls as $1/N$ when
the same taper is spread over more rows, so it vanishes in the continuum limit. The
decoupling test departs
further: by at most $1.7\times10^{-3}$ from its
own run-outs' divider across its sweep, the leak of the run-out mismatch it
carries, which falls as its rung-to-outlet conductance ratio rises (C8, P-C8.3).
CQ2's $\peMarginMin\times$ margin is on a different axis, the pressure-equilibration
half-life against the device's length, and is not a bound on that departure: the
test's ladder clears the half-life criterion and carries the departure anyway. The graph
keeps the symbolic separation and the executable test apart, giving the test a
conclusion of its own, \texttt{f\_indep\_fgap\_executable}, distinct from the
symbolic \texttt{f\_indep\_fgap} this node argues.
\prov{Routt 2026 Table~I, \S 8.4; code:
\_emergent\_F\_from\_solve, verify\_against\_notebook N1--N3}

\end{description}

% ----------------------- CRITICAL QUESTIONS -----------------
\field{Critical questions.}
The current-divider identity is deductive \emph{given} its premises and carries
no critical questions of its own; the questions below test those premises
(pressure equilibration; $\Qs(0)=0$) and the empirical accuracy of $F$'s
\emph{magnitude}, never the separation $F\indep\fgap$, which is symbolically
exact.
\begin{description}

\item[CQ1: Exception (does the 1-D residual couple $F$ to $\fgap$?).]
No. That residual ($\residPctVOne\%$ at V1, $\residPctVThree\%$ at V3; the realized-yield factor
$\eta = \etaVal$ is the representative single magnitude, matching V1 closely rather
than holding at both) concerns $F$'s own magnitude accuracy, not the
separation, and it closes via 3-D outlet-bifurcation fields that are themselves
geometry-dependent (P-C4.2); the separation is symbolically exact. A
\emph{sibling} 3-D-junction effect, junction hydrodynamics of the same
\emph{kind} as the physics that sets $F$'s residual, does perturb the per-gap
$\feff$ the 1-D ladder estimates, contributing an \emph{estimated}
${\sim}\siblingSqrtcPct\%$-scale uncertainty on $\sqrt{c}$ (to be bounded by a 3-D junction
solve), distinct from $F$'s residual and not entering $F$. \emph{Sibling} is a
statement of kind, not of identity: whether this yield-side residual and the
cutoff-side residual in fact share a common three-dimensional origin is an
\emph{open question}, deferred to future work in the glossary, and nothing in
either appendix turns on the answer. The two are treated throughout as
corrections on independent axes. The
graph carries that estimate as the one open exception on this node,
\texttt{sibling\_3d\_junction\_uncertainty\_on\_sqrt\_c}, attached to $F$'s
\emph{magnitude} and not to the separation.
$\rightarrow$ \emph{separation holds.}

\item[CQ2: Is the PE premise satisfied?]
Yes; PE is tested, not assumed. It is the assumption the lineage's recursion rests
on: Gifford treats ``\emph{the gaps between the flow channels of the array as nodes
which allow equalization of pressure across the width of the device}''
\cite[p.~4497]{b1}, and reports that ``\emph{This assumption has been validated by
observations of fluid flow in this geometry}'' \cite[p.~4498]{b1}. The companion
preprint's diagnostic fits the
half-life $k_{1/2}$ of the cross-width pressure imbalance and requires
$k_{1/2}/N_{\mathrm{total}}<0.05$, at least twenty half-lives within the device.
Its five reference devices are Gifford 2014's bead device, Dinh 2024's three
variants (V1, V2 and V3) and Xia~RS-7: five devices over three published papers.
In the preprint's computation all five clear the criterion with a safety margin
of at least $\peMarginMin\times$ (Xia~RS-7 tightest). Those half-lives depend on the
thickness of the wall each gap crosses, which the papers do not state: the
preprint's code takes it equal to the gap's width, and its 3-D COMSOL models of
V1 and V3 use the same wall, so the agreement of their outlet half-lives with the
1-D network's, to $\peHalfAgreeVOne\%$ (V1) and $\peHalfAgreeVThree\%$ (V3), tests the network and not the
wall. With the thicker walls the design code uses, every half-life is longer and
every device still clears the criterion. What the
criterion measures is the decay of an imbalance along the ladder; it does not see
one the outlet row itself holds, such as the decoupling test's or a taper's (P-C4.3). The code's
\texttt{ARG\_MATH.REFERENCE\_DEVICES}, read by the Reynolds-number gates, is a
different list of five (two Gifford widths, Lezzar, a Dinh row, and a row
labelled Xia~RS-7 that carries Dinh's outlet geometry); the preprint reports the
criterion for its own five, not for that list.
\prov{Routt 2026 \S III-A7 and Table~II (the \S 8.7 diagnostic); code: reference
implementation \S 8.7, pe\_consistency\_check (gap\_p $=$ gap\_size)}
$\rightarrow$ \emph{premise holds for the imbalance the criterion measures, on all
five devices.}

\item[CQ3: Is $\Qs(0)=0$ a genuine rule?]
Yes; it is a device-level architectural boundary condition, a single
centre-channel inlet with dual side outlets, shared by every device in the
published lineage (Gifford, Xia, Lezzar, Dinh), not a modelling choice.
\prov{Routt 2026 \S II} $\rightarrow$ \emph{rule established.}

\end{description}

% ----------------------- WHY THIS LEMMA EXISTS --------------
\field{Why this lemma exists (the cautionary tale).} The canonical equation set
carries two expressions both labelled $F$: the outlet identity
$F=2\Rc/(\Rs+2\Rc)$ ($\fgap$-free) and the cumulative product
$F=1-\prod_{i}(1-2\fgap(i))\approx1-e^{-2N\fgap}$ ($\fgap$- and $N$-dependent).
They are the same quantity only at the self-consistent outlet but expose
different variables, and conflating them is exactly the error this project
retracted. The retracted ``$\feff\!\cdot\!N=0.3516$ fixed by $F$'' line read
$N\fgap$ out of $F$ via $\ln(1-F)=-2N\fgap$: that relation belongs to the
\emph{cumulative} form, recovers only the product $N\fgap$ ($N$ unpublished), and
never touches the local over-extraction $c$ that sets the cutoff; moreover,
$0.3516$ is that product, a cumulative-yield number, not the outlet $F$ at all. The
inversion fails on three counts.

\emph{Where the figure comes from, and why it carries no macro.} $0.3516$ is
historical: it is the value the retracted construction $\feff=0.3516/N$ asserted,
recorded with that construction in the glossary's revision record. It is
deliberately not given a value macro. The generated value block holds the
project's \emph{live} values, and a retracted figure that tracked them would
rewrite the record of what was retracted; two of its four appearances are inside
verbatim quotations of the retracted line, where substitution is disallowed in
any case. The $0.505$ the check that follows is fed is kept literal for the same
reason. It is Dinh's measured V1 yield, and it is also the \emph{input} the
retracted construction was built on, since $-\ln(1-0.505)/2$ is where $0.3516$
came from; the input to a retracted computation is as historical as its output.
Substituting the value macro would print the same characters today, and would
leave a frozen figure paired with a live antecedent that the next regeneration
could move on its own. It is checkable where it stands, which is what a macro would have
bought: $\ln(1-F)=-2N\fgap$ returns $N\fgap=0.3516$ at $F\approx0.505$, one
residual-width below the band the closed form predicts for the two validated
variants (P-C4.2), while reading the same number as $F/2$ in the small-$f$ limit
puts $F\approx0.703$, near the top of that band. One number yielding two values of $F$
nearly $0.20$ apart is the arithmetic signature of the conflation this lemma
exists to block. The lemma is invoked only to license the step
this appendix states elsewhere as \emph{narrow the per-gap over-extraction $c$ by
a local resistance-network calculation, not via $F$} (C3d~CQ3; C7~CQ2).

% ----------------------- STATUS -----------------------------
\field{Status.} The separation $F\indep\fgap$ is \textbf{proven} (symbolically
exact; network-emergent; schedule-independent). $F$'s absolute magnitude is
independently \textbf{validated} by 3-D COMSOL, within $1\sigma$ on the
raw-$\mathrm{SD}$ scale and within $2\sigma$ on the stricter $\mathrm{SEM}$ scale,
while the calibration-free 1-D form sits $\EsdClosedVOne\sigma/\EsdClosedVThree\sigma$ from the
measurement (P-C4.2). The sibling
${\sim}\siblingSqrtcPct\%$ uncertainty on $\sqrt{c}$ is an estimate, to be bounded by a
3-D junction solve; that it is a junction effect of the same \emph{kind} as
$F$'s is not a claim that the two residuals share an origin, which is open
(CQ1).


% ===========================================================================
\clearpage
\section*{C9 --- Transport lemma: why the steric factor is exactly $2$}

\field{Summary.} The steric prefactor $\ksteric=2$ enters C1 as an
\emph{asserted} constant, carried by convergent testimony. It can be
\emph{derived}. Track a fluid element through a gap: its flux fraction measured
from one wall obeys $\varphi'=(\varphi-f)/(1-2f)$, so the centreline
$\varphi=\tfrac12$ is a \emph{repelling} fixed point, fluid migrates
deterministically toward the wall, and a particle is extracted iff
$\varphi_{0}\le F/2$. That repulsion is what makes the selection a hard
\emph{wall constraint} rather than a statistical position criterion, and with one
line of rigid-sphere geometry (a centre cannot come closer to a wall than its own
radius) it pins $\ksteric=2$ exactly. The lemma is not a restatement of C1: it
\emph{upgrades}
C1's testimonial premise to a derivation, and the same solve pays for itself
three more times. Diffusion is shown to broaden the transition without moving
the cutoff, the Fig.~1 ``about half'' caption is deselected on \emph{physical}
grounds independent of C1's textual adjudication, and, within the model's
premises, the residual C3d leaves open is pinned to three model gaps. The node is
deliberately \emph{not} load-bearing: no support path runs from $\ksteric=2$ to
the thesis, because C9 strengthens a premise rather than adding a branch.

% ----------------------- NOTATION ---------------------------
\field{Notation (one collision to forestall).} This section uses the curly
$\varphi$ for the \emph{flux fraction} of a fluid element measured from one wall,
the quantity the transport map advances, and reserves the straight-stroke $\phi$
for the \emph{wall porosity} named in C9d's residual list below. C3d's own
decomposition never mentions porosity, which is part of what C9d adds. The two
are unrelated; the specification distinguishes them the same way.

% ----------------------- ARGUMENT GRAPH ---------------------
\begin{figure}[t]
\centering
\begin{tikzpicture}[
  font=\scriptsize,
  prem/.style={draw, rounded corners, align=center, inner sep=3pt,
               text width=10.0cm, minimum height=13mm},
  aux/.style={draw, rounded corners, align=center, inner sep=3pt,
              text width=4.6cm, minimum height=12mm},
  ded/.style={draw, double, double distance=0.7pt, rounded corners,
              align=center, inner sep=3pt, text width=3.1cm, minimum height=16mm},
  up/.style={-{Stealth[length=2mm]}, thick},
  nt/.style={font=\scriptsize\itshape, text=black!60, align=center},
]
% --- the premise ---
\node[prem] (lem) at (0,0)
  {\textbf{P-C9.1: transport lemma.}\ $\varphi'=(\varphi-f)/(1-2f)$: the
   centreline $\varphi=\tfrac12$ is a \emph{repelling} fixed point, so fluid
   migrates deterministically toward the wall and a particle is extracted iff
   $\varphi_{0}\le F/2$. Checked in code by bisecting the recursion itself:
   $|2\varphi_{\mathrm{crit}}-F|<10^{-6}$.};
% --- the four deductive consequences ---
\node[ded] (a) at (-5.7,-2.9)
  {\textbf{C9a}: constant\\ hard wall $+$ centre\\ exclusion (P-C9.1$'$)
   $\Rightarrow \ksteric=2$\\ \textsf{[proven]}};
\node[ded] (b) at (-1.9,-2.9)
  {\textbf{C9b}: diffusion\\ broadens, does not\\ move the cutoff\\
   \textsf{[proven]}};
\node[ded] (c) at (1.9,-2.9)
  {\textbf{C9c}: caption\\ a hard wall cannot be\\ ``about half''
   $\Rightarrow$ deselected\\ \textsf{[proven]}};
\node[ded] (d) at (5.7,-2.9)
  {\textbf{C9d}: within the model's\\ premises, the residual\\ C3d leaves open\\ lives in three
   model gaps\\ \textsf{[proven]}};
% --- the two auxiliary premises, each drawn below the one conclusion it feeds,
%     as the deductive strand is drawn below its claim in Fig. 2 ---
\node[aux] (exc) at (-5.7,-5.5)
  {\textbf{P-C9.1$'$.}\ centre exclusion: a sphere of radius $a$, carried with
   its centre, cannot bring that centre within $a$ of the wall
   $\Rightarrow$ extracted iff $a\le\delta$};
\node[aux] (dif) at (-1.9,-8.1)
  {\textbf{P-C9.2.}\ diffusion broadens the transition and cannot move the
   cutoff. Per rung, migration ${\sim}\delta/2=0.22\um$ vs diffusion
   $\sqrt{2D\,\Delta t}=0.033\um$, $6.7\times$ smaller (the specification's
   estimate, illustrative); over a transit, not small};
% --- edges ---
\draw[up] (lem) -- (a);
\draw[up] (lem) -- (b);
\draw[up] (lem) -- (c);
\draw[up] (lem) -- (d);
\draw[up] (exc) -- (a);
\draw[up] (dif) -- (b);
\node[nt, anchor=west, text width=4.6cm] at (1.1,-5.6)
  {P-C9.1$'$ is drawn into C9a and P-C9.2, which uses the same exclusion, into
   C9b; C9c and C9d rest on the transport lemma alone.};
\end{tikzpicture}
\caption{Argument graph for C9. One transport lemma, two auxiliary premises, four
deductive consequences. Border style follows the overview map: every consequence
carries the \emph{doubled} border of a proven identity, and there is no dashed box
anywhere: this is the only section of the appendix in which \emph{every}
conclusion is deductive. C9a is the result that matters for the chain, since it
converts C1's asserted $\ksteric=2$ into a derived one; C9b, C9c and C9d are
dividends of the same solve. The two auxiliary premises are drawn \emph{below}
their conclusions, as the deductive strand is in Figure~\ref{fig:c1-argmap},
because each feeds one consequence rather than all four. Note that the frozen
argument graph records only the transport lemma as a premise of
\texttt{c9\_transport}: P-C9.1$'$ is the geometric step its CAES encoding folds
into the conditional, and it is drawn here because a figure that hid it would
make C9a look stronger than it is.}
\label{fig:c9-argmap}
\end{figure}

% ----------------------- CLAIM ------------------------------
\field{Claim.} The streamline-transport solve derives $\ksteric=2$ from the flow
model, taking from the literature only the geometric exclusion of P-C9.1$'$ and
no measurement, no calibration and no reported result; and in doing so it settles
three further questions that C1 and C3d leave open.

% ----------------------- SCHEME -----------------------------
\field{Scheme.} Defeasible Modus Ponens for C9a (the transport lemma together
with centre exclusion implies the constant), and Argument from an Established
Rule for C9b--C9d (the same solved rule applied to three further questions). All
four are \emph{deductive}: every premise they rest on is a closed-form map or a
geometric fact, never a measurement, so the scheme's exceptions slot is empty and
the nodes are labelled \emph{proven} with no open critical question. The
per-rung length scales of P-C9.2 show how narrow the transition is per rung; C9b
does not rest on them.

% ----------------------- PREMISES ---------------------------
\field{Premises.}
\begin{description}
\item[P-C9.1 (the transport lemma).]
A fluid element entering a gap at flux fraction $\varphi$, measured from one
wall, leaves at $\varphi'=(\varphi-f)/(1-2f)$; equivalently
$(\varphi'-\tfrac12)=(\varphi-\tfrac12)/(1-2f)$. The centreline
$\varphi=\tfrac12$ is a fixed point of this map and it is \emph{repelling}:
$|\mathrm{d}\varphi'/\mathrm{d}\varphi|=1/(1-2f)>1$ for every $f>0$. Fluid
therefore migrates deterministically toward the wall, and solving the threshold
gives extraction iff $\varphi_{0}\le\tfrac12(1-P(N))=F/2$, with the two walls
contributing $2\times F/2 = F$. The code checks this by simulation rather than
algebra: the design code's validator M18 finds the critical starting position by
bisecting the recursion itself, gap by gap, and requires
$|2\varphi_{\mathrm{crit}}-F|<10^{-6}$ against $F$ from the cumulative
extraction. \prov{spec \S\,6.2; code: CIF\_v2 critical\_phi (validator M18);
ARG\_MATH transport\_map, critical\_phi}
\item[P-C9.1$'$ (centre exclusion: the step from fluid to particle).]
The lemma above is about a \emph{fluid element}. The step to a particle is one
line of geometry, stated here rather than assumed because without it the factor
$2$ does not follow. The model carries a particle with its centre on the
streamline that centre occupies, and a sphere of radius $a$ cannot bring its
centre closer than $a$ to the wall, so it is extracted if and only if
$a\le\delta$, whence $\dcrit=2\delta$ and $\ksteric=2$. Dinh state the same
criterion for their devices: ``\emph{we would expect particles with radii
exceeding the width of the extracted flow lamina ($a>w$) to remain in the middle
channel}'' \cite[p.~917]{b5}. This is the same exclusion premise C1 takes from
Yamada \& Seki and from Inglis \& Herman; the difference is that here it does no
evidential work of its own. It is a property of rigid spheres, not a finding
about CIF, and it is the \emph{only} thing C9a needs beyond the solved map and
the model's carrying of the centre, whose error is the finite-particle
displacement C9d lists. \prov{spec \S\,6.2, consequence~2}
\item[P-C9.2 (diffusion broadens the transition and cannot move the cutoff).]
Per rung, deterministic migration moves a particle by ${\sim}\delta/2=0.22\um$
while cross-streamline diffusion moves it by
$\sqrt{2D_{\mathrm{diff}}\Delta t}=0.033\um$, a factor $6.7$ smaller. These are
the specification's estimates; it does not give their inputs, and nothing here
recomputes them. Over a whole transit the spread is not small: the specification
puts it at $1$--$4\,\delta$ for calibration beads, and first-order for sub-micron
particles or long transits. What diffusion cannot do is move the cutoff. A
sphere of radius $a>\delta$ cannot bring its centre inside the lamina however it
wanders, so diffusion broadens the transition below $a=\delta$ and leaves the
cutoff where exclusion puts it. \prov{spec \S\,6.2, consequence~1, and \S\,20
(diffusion)}
\end{description}

% ----------------------- CRITICAL QUESTIONS -----------------
\field{Critical questions.}
\begin{description}
\item[CQ1: Does the lemma merely restate C1?]
No, and the distinction is the point. C1 establishes $\ksteric=2$ by
\emph{convergent testimony} across the Shevkoplyas lineage, Yamada \& Seki, and
Inglis \& Herman, and the rule it leans on is itself stated, not derived: Inglis
writes ``\emph{The critical particle diameter, D$_c$, for determining which path a
particle will follow is then phenomenologically given by}'' $D_c=2b$
\cite[p.~656]{b17}. C9 establishes the same constant by \emph{solving the flow
model}, taking from those sources only P-C9.1$'$, a geometric property of rigid
spheres that no measurement and no laboratory is being trusted for. The label does
not change (C1 was already warranted), but what the premise is worth does: a
reader who declines every cited \emph{finding} still has the constant, because
what remains to be granted is that a sphere's centre cannot sit closer to a wall
than its own radius. $\rightarrow$ \emph{answered: it upgrades the warrant, not
the label.}
\item[CQ2: If C9 is the stronger result, why is the thesis not argued through it?]
Because it strengthens a premise rather than adding a branch. On the graph's own
support edges there is no path from $\ksteric=2$ to the thesis: the constant is
consumed by exactly two claims, the mechanism (C1) and the residual
decomposition (C3d), and neither carries through to the thesis either. This is a
design decision, not an omission: making a granted premise sturdier changes how
much confidence it deserves without changing a single label.
$\rightarrow$ \emph{answered: deliberately not load-bearing.}
\item[CQ3: Does the deselection of the caption double-count C1's adjudication?]
No: the two routes are independent, and each is on its own sufficient. C1
deselects the Fig.~1 caption on \emph{textual} grounds (the caption and the main
text are formally inconsistent; the main text is selected, backed by Inglis's
$D_c=2b$). C9 deselects it on \emph{physical} grounds: a hard wall constraint
cannot be ``about half'' of anything. Withdrawing either route alone leaves the
deselection standing and moves no label; withdrawing \emph{both} flips
\texttt{caption\_deselected} to \emph{out} and takes $\ksteric=2$ with it (four
labels move: \texttt{caption\_deselected}, $\ksteric=2$ and its negation, and
C3d's \texttt{residual\_decomp}). That is the precise sense in which the
redundancy is real rather than decorative. $\rightarrow$ \emph{answered:
independent and individually sufficient.}
\end{description}

% ----------------------- RESIDUAL / STATUS ------------------
\field{Residual risk.} None for the four conclusions: each is a deductive
consequence of the recursion, which the code checks by simulating it, and of
centre exclusion. P-C9.2's per-rung figures are the specification's estimates and
illustrate the transition's width; C9b does not rest on them. The residual that
C9d \emph{pins} is C3d's, not C9's.

\field{Consequence for C3d (worth reading across).} Within the model's premises,
C9d closes the list of places the residual C3d leaves open can live: the flat-lamina idealisation, the wall
porosity $\phi$, and finite-particle displacement: of the premises the lemma uses, these
are the ones that can fail once its declared scope (the thesis's Leg~5) is granted, and
a source outside the model, such as how a cutoff is read from data, is not on the
list. C3d first divides out the
factors it measures ($\kinplane=1$, and the depth factor $\Kdepth=R(\alpha)$, which
corrects the flux fraction of a flat slab in a finite-depth duct and is not one of
the three gaps), then names what remains by its effects: per-gap over-extraction
$\sqrt{c}$ and cumulative finite-particle displacement $C$. C9d names the model
gaps that can produce that product, and the wall porosity is named only here.
C9d therefore answers the \emph{exhaustiveness} question C3d leaves open, within
those premises, while
leaving C3d's \emph{identifiability} question ($\sqrt{c}$ and $C$ degenerate
from the published data, being tied to the unreported gap count $N$) exactly
where it was.

\field{Status.} \textbf{Proven}: four deductive consequences of one solved
transport map, with no open critical question.

% ===========================================================================
\clearpage
\section*{C8 (graph) --- The design map: targets to geometry, and three dormant defeaters}

\field{Summary.} Because cutoff and yield are governed by separate laws (C4), and
the plateau schedule holds the cutoff at every row (P-C8.1), the
map from performance targets to device geometry has a \emph{lower-triangular}
Jacobian: changing the gap-row count $N$ leaves the cutoff untouched
($\partial\dcrit/\partial N=0$) while moving the yield, through the terminal width,
and so the outlet ratio, that $N$ itself sets. The zero comes from the
invariant-lamina schedule on the plateau (P-C8.1): S1 holds the lamina at every
row and the plateau prefactor is constant, so the cutoff $\kapp\,\delta^{*}$ does
not depend on $N$. C4 supplies what makes the two targets answer to separate laws:
the yield law carries no per-gap parameter. Neither supplies the
diagonal; the diagonal is non-zero because every row skims a positive fraction, so
the retained flow $\prod_{i}(1-2\fgap(i))$ falls strictly with $N$. The inverse map
therefore exists: \emph{locally}, and throughout the declared envelope. A
triangular Jacobian gives local invertibility and nothing more; this appendix does
\textbf{not} claim a global one-to-one map over an unrestricted domain, and the
distinction is stated here rather than left for a reader to supply. That
conclusion is then reached a second time, by construction rather than by
argument: the design code's forward march, which never inverts the design map,
computes the width schedule for five cutoff targets and finds $N\!\cdot\!D^{2}$ invariant to
$0.29\%$ while $N$ itself moves by a factor of $4.5$. The row count is simply the
price of the target \emph{pair}. Forward substitution then supplies the practical
inverse for a device built to the resulting schedule. Realizability, however, is
\emph{conditional}, and the argument says so before a reader has to ask. Three
defeaters would break the realization if their antecedents held: the device
running below the knee, a run-out mismatch, and a row count pinned by a length
budget. All three are \emph{withheld}: their antecedents are not granted, so they
lie dormant and the design stands. But withheld means ``not established for this
scoped design,'' not ``shown never to occur,'' and published designs are
worked counterexamples to unconditional application: at their wide ends the
plateau law puts Dinh's three designs at $u=\dwc=0.029$--$0.032$
(the bracket device at $0.029$), and Xia's RS-7, on the taper Xia prints, lower
still at $0.023$ (C5), below the plateau floor
$u_{\mathrm{lo}}$ (set here to the plateau-bin edge $\plateauUwc$), and, on the cubic
pair the appendix argues on, the implementation's guard refuses them. This is the
leg on which the headline claim
inherits its conditionality.

% ----------------------- ARGUMENT GRAPH ---------------------
\begin{figure}[t]
\centering
\begin{tikzpicture}[
  font=\scriptsize,
  prem/.style={draw, rounded corners, align=center, inner sep=3pt,
               text width=3.0cm, minimum height=15mm},
  ded/.style={draw, double, double distance=0.7pt, rounded corners,
              align=center, inner sep=3pt, text width=3.0cm, minimum height=15mm},
  dclaim/.style={draw, double, double distance=0.7pt, rounded corners,
                 align=center, inner sep=4pt, text width=9.0cm, minimum height=10mm},
  part/.style={draw, dashed, rounded corners, align=center, inner sep=4pt,
               text width=9.0cm, minimum height=12mm},
  cav/.style={draw, dashed, rounded corners, align=center, inner sep=3pt,
              text width=3.4cm, minimum height=16mm},
  run/.style={draw, rounded corners, align=center, inner sep=3pt,
              text width=3.0cm, minimum height=13mm},
  up/.style={-{Stealth[length=2mm]}, thick},
  cut/.style={-{Stealth[length=1.8mm]}, dashed},
  nt/.style={font=\scriptsize\itshape, text=black!60, align=center},
]
% --- tier 1: the four premises of the deductive route to the inversion ---
\node[prem] (p1) at (-5.1,0)
  {\textbf{Skim law} (C2)\\ $\delta=\wc\sqrt{\fgap/3}$};
\node[prem] (p2) at (-1.7,0)
  {\textbf{Proportionality}\\ (C3a) $\dcrit=\kapp\,\delta$};
\node[ded]  (p3) at (1.7,0)
  {\textbf{Separation lemma}\\ (C4) $F\indep\fgap$: the yield law\\ carries no per-gap parameter};
\node[prem] (p4) at (5.1,0)
  {\textbf{Schedule S1}\\ $\delta(i)=\delta^{*}$ at every row\\
   \textsf{[assumed; M8]}};
% --- the SECOND, independent route: the executable forward march ---
\node[run] (ex) at (6.6,-2.05)
  {\textbf{Forward march} (P-C8.2)\\ $N\!\cdot\!D^{2}$ constant to $0.29\%$ over
   five cutoff targets\\ \textsf{[warranted]}};
% --- tier 2: the inversion ---
\node[dclaim] (tri) at (0,-2.9)
  {\textbf{C8a.}\ the target Jacobian is lower-triangular
   ($\partial\dcrit/\partial N=0$) with non-zero diagonal $\Rightarrow$ a
   \emph{local} inverse exists throughout the envelope\ \ \textsf{[proven]}};
% --- tier 3: the realization ---
\node[part] (proc) at (0,-5.2)
  {\textbf{C8b.}\ forward substitution realizes \emph{both} targets on a
   device built to the schedule, inside the envelope\\ \textsf{[partially closed]}};
% --- tier 4: the three dormant defeaters ---
\node[cav] (d1) at (-5.6,-7.9)
  {\textbf{CQ7} (withheld)\\ below the knee: $\kapp$ becomes a function of the
   cutoff (the C7 fixed point)};
\node[cav] (d2) at (0,-7.9)
  {\textbf{CQ6} (withheld)\\ run-outs mismatched: a $3\times$ arm-level mismatch
   collapses $F$, $\runoutFmatched\to\runoutFmismatched$};
\node[cav] (d3) at (5.6,-7.9)
  {\textbf{CQ5} (withheld)\\ $N$ pinned by a length budget: two targets, one
   free variable};
% --- edges ---
\draw[up] (p1) -- (tri);
\draw[up] (p2) -- (tri);
\draw[up] (p3) -- (tri);
\draw[up] (p4) -- (tri);
\draw[up] (ex) -- (tri);
\draw[up] (tri) -- (proc);
\draw[cut] (d1) -- (proc);
\draw[cut] (d2) -- (proc);
\draw[cut] (d3) -- (proc);
\node[nt, anchor=west, text width=3.0cm] at (4.9,-6.5)
  {each \emph{undercuts} C8b if its antecedent is granted; none is};
\end{tikzpicture}
\caption{Argument graph for C8. The inversion (C8a, doubled border: proven) is
reached twice over. The deductive route is the top row: the skim law,
proportionality, the separation lemma and the S1 schedule together make the
target Jacobian lower-triangular. The second route, entering from the right,
is an executable forward march that never inverts the design map: run at five
cutoff targets, it computes the width schedule that realizes each, which reaches the
same conclusion by construction, as Evidence to a Hypothesis rather than as an
entailment. The two routes share the physics and not the inference, so an error
in the triangularity reasoning would leave the march untouched. Only then does forward substitution realize both targets on a
device built to the resulting schedule (C8b). The three dashed boxes are the envelope boundaries;
their dashed arrows are \emph{undercutters}, not corroboration: each would
withdraw C8b's warrant if its antecedent were granted, and none is granted here.
C8b itself is dashed because its status is \emph{partially closed}, exactly as
C3d is dashed in Figure~\ref{fig:argmap}. Note that the undercutters attach to
C8b and not to C8a: C8a is claimed only inside the envelope, so what the
boundaries qualify is the realization, never the inversion.}
\label{fig:c8-argmap}
\end{figure}

% ----------------------- CLAIM ------------------------------
\field{Claim.} The target-to-geometry map is \emph{locally} invertible throughout
the declared envelope, and forward substitution realizes both the cutoff and the
yield target on a device built to the resulting schedule inside that envelope.

% ----------------------- SCHEME -----------------------------
\field{Scheme.} Two schemes reach C8a, which is why it is drawn with two
incoming routes. Defeasible Modus Ponens for the deductive one: a
lower-triangular Jacobian with non-zero diagonal implies a local inverse, which
is an entailment, not an inference to the best explanation. Evidence to a
Hypothesis for the executable one (P-C8.2): the design code's forward march is
run, it computes the width schedule for every target without inverting anything, and
the hypothesis is confirmed in the defeasible sense. That scheme's own critical
question, whether the invariance is an artifact of the march's model (its $F$ is the
cumulative product $1-\prod_{i}(1-2\fgap(i))$, with no outlet to equilibrate), is
discharged by measurement in P-C8.3, on a network that has one, not assumed away.
That network is a uniform ladder and not the marched taper: the one attempt to put
the taper through the network inverse did not converge, and was reported rather than
tuned. C8b then takes Argument from an
Established Rule: the forward-substitution rule is applied to a case, and the
scheme's exceptions slot is where the three defeaters live. The rule is the
lineage's own design practice, which is what makes it established rather than
proposed here. In Gifford the width of each side channel at gap $i$ is
``\emph{ultimately calculated merely by the desired value of $\fgap$, and their
width at the previous gap}'' \cite[p.~4497]{b1}: a forward recursion, not an
inversion of the design map, whose only solve at each row recovers a width from
the duct-resistance series \cite[pp.~4498--4499]{b1}, as the march's does
(P-C8.2). Xia
state the order that C8a's triangular Jacobian licenses, the cutoff's parameter
first and the yield after it: ``\emph{Once the desired value of $\fgap$ (and
dimensions of the centre channel, posts, and gaps) are selected, a CIF-based
device can be rapidly patterned numerically, and structured into a paperclip-like
format (Fig. 1a) to achieve the preferred amount of particle
concentration/filtration in a compact footprint}''. Along the device the side
channels widen to a final width that ``\emph{is determined by the degree of
particle concentration (or, equivalently, the amount of filtration) preferred for
a given application}'' \cite[p.~2 (Sci Rep PDF pagination)]{b2}.

% ----------------------- KERNEL DECLARATION -----------------
\field{Kernel declaration (required before any of this means anything).} The
glossary's pairing rule is directive: $\kapp$ is meaningless without the kernel it
was fitted with. Two kernels are in use: the leading-order skim branch
$g:\ \delta = \wc\sqrt{\fgap/3}$ with prefactor $\kappa_{g} = \kappVal$, and the
lamina-flux kernel $F_{\mathrm{lam}}$, exact for a flat lamina in an unperturbed duct
(which the skimmed fluid is not), whose prefactor $\kappa_{F}=\kFval$ is not exact
even in that sense: it is transferred from $\kappa_{g}$ at an assumed $80\um$
calibration depth. A pair must be declared and used at all three call sites: the
inversion, the schedule, and the scale $\dcrit = \kappa_{\mathcal K}\,\delta$.
\textbf{This appendix argues on the cubic branch with $\kappa_{g}$}, as
the glossary directs while the calibration depth remains unmeasured: every prefactor
fitted and every cutoff computed in C1--C7 is a cubic-branch number, and C8 uses
the other pair only where it says so, in the design code's march (P-C8.2) and the
kernel comparison below.

That declaration matters here more than anywhere else in the appendix, because the
design code carries both pairs: its \texttt{design} function defaults to the cubic
pair, while its physics layer, the one C4's design-map test marches (P-C8.2), runs
on $F_{\mathrm{lam}}$. The two are not interchangeable and the
mismatch is \emph{not} detected automatically: at Dinh's geometry and the bracket
device's gap fraction, comparing the exact inversions of the two kernels as
the glossary does, the pair choice alone moves the cutoff by $+\choiceIn\%$ at the
inlet row,
rising to $+\choiceOut\%$ at the outlet, and crossing a prefactor onto the other
kernel's lamina shifts it a further $+\crossUp\%$ or $\crossDown\%$. Switching costs
$+26.9\%$ on the row count at the assumed $80\um$ calibration depth, falling to
$+3.9\%$ at $150\um$: figures the glossary flags as \emph{ungated}, coming from a
calibration-depth sweep the derivation record does not reproduce, so they are
carried here as an order of magnitude rather than as verified constants. So the
claim C8b makes is that forward substitution realizes both targets \emph{with one
declared pair used throughout}; a design that inverted on one kernel and scaled on
the other would be a silent modelling error of that size, and it is the
specification's pairing rule, not this argument, that prevents it. The glossary
records the transfer bias as an uncertainty band rather than applying a correction
whose calibration cannot be verified, and this appendix does the same.

% ----------------------- PREMISES ---------------------------
\field{Premises.}
\begin{description}
\item[P-C8.1 (the invariant-lamina schedule).]
The canonical schedule S1 holds the diverted lamina constant along the array,
$\delta(i)=\delta^{*}$ at every row. This project's own design specification
states it as ``\emph{S1 --- invariant lamina (canonical). Exact at every aspect
ratio}'' (\S\,7); the quotation is from an internal document, not the literature,
and is quoted only so that the schedule the argument relies on is fixed in
writing. It is this invariance, on the plateau where the prefactor is a constant,
that makes the cutoff a property of the \emph{schedule} rather than of the row
index, and so makes $\partial\dcrit/\partial N=0$ exact rather than approximate.
Below the plateau it does not: $\kapp$ then depends on $\dwc$, so a lamina held
constant no longer holds the cutoff. C7's computation for Dinh's \emph{bracket
device} is the worked case: under the schedule Dinh describes, read in the
corrected Xia form, the leading-order lamina stays at $0.442\um$ while the
self-consistent cutoff rises from $1.09$ to $2.02\um$ along the linear section. The
plateau guard
(CQ7) keeps a design in the plateau bin, where the prefactor is taken as constant;
C7's smooth fit of $\kapp(\dwc)$ reaches the plateau value only at
$\dwc\approx0.051$, so just above the floor the premise is exact on the bin's model
rather than on that fit. The schedule is an assumption of the design, and the
validator that tests it exactly is~M8.
\prov{spec \S\,7 (S1); validator M8}
\item[P-C8.2 (the forward march: a second route to C8a).]
This premise does not invert the design map, and that is the point: the march
\emph{computes} $f(i)$ from $\delta^{*}$ rather than solving for it, so it reaches
the same conclusion as the Jacobian argument by a different \emph{inference}. The
two are not physically independent: the march is the design code's physics
layer, on the pair its own lineage spec declares (the lamina-flux kernel
$F_{\mathrm{lam}}$ with $\kappa_{F}=\kFval$, the Kernel declaration above), which
is the lamina physics of the skim law at finite depth. But they fail differently, which is what makes the
redundancy worth having: an error in the triangularity reasoning would not move
the march, and a convergence failure in the network inverse cannot reach a
computation that never calls it (the march's only solve, for the channel width at
each row, is a bisection on a monotone function). Sweeping the cutoff across five
targets while holding the yield target fixed leaves $N\!\cdot\!D^{2}$ constant to $0.29\%$
($[14764,\,14776,\,14781,\,14791,\,14807]$, tolerance $2\%$) while $N$ itself
spans $1446$--$6562$, a factor of $4.5$, as the specification's closed form for the
canonical taper requires ($N\propto D^{-2}$, \S\,13); the converse sweep meets every
yield target while the row count moves. Most of this holds by construction and is
not evidence: the march fixes $\delta^{*}$, and so the cutoff, at every row (S1), so
$\partial\dcrit/\partial N=0$ is its input rather than its finding; it stops at the
first row that meets the yield target, so meeting the target is its stopping rule,
and the converse sweep's $<10^{-3}$ with it; and its taper step is the closed form's
own, so $N\propto D^{-2}$ follows to leading order. What it does test is that the
computation is feasible across the targets, with the taper not exhausted before the
yield target is met, and that on the exact kernel the row count stays
$0.26$--$0.55\%$ above the closed form, which is written for the leading-order
schedule (S1$'$). That is the sense in which it reaches C8a a second time: the inverse is
computed rather than argued, and the computation agrees with the specification.
The row count is therefore the price of the target \emph{pair}, and the two
targets are independently specifiable. The row counts are $F_{\mathrm{lam}}$-pair numbers, not cubic-branch
ones: the pairing rule applied, one declared pair used throughout. The march
also holds its prefactor constant and applies no plateau guard, and it runs on the
$75\um$ lineage inlet, where only its largest target, $3.2\um$, reaches the
plateau floor $u=\dwc\ge\plateauUwc$; the other four, the converse sweep's fixed
cutoff among them, are targets the plateau guard would refuse. What it
corroborates is therefore the triangular structure of the constant-prefactor law,
which is what C8a asserts inside the envelope, and not a realization below the
knee.
\prov{code: tests/test\_C4\_design\_map.py}
\item[P-C8.3 (decoupling, executable, and the discharge of P-C8.2's critical question).]
The separation lemma is not assumed here but measured: the filtrate fraction $F$
stays at the value the outlet resistance ratio fixes while the \emph{measured}
$\fgap$ spans $3.63\times10^{-5}$ to $1.93\times10^{-3}$, a range of $1.73$
decades. Across that sweep the worst residual is $1.7\times10^{-3}$, at the
weakest outlet equilibration reached ($\mathrm{PE}=59\times$), inside the
$O(1/\mathrm{PE})$ bound the test sets there, $2.6\times10^{-3}$, and it falls to
$2.3\times10^{-5}$ at $\mathrm{PE}=5093\times$. The sweep moves PE and $\fgap$
with one knob, the rung width, so that fall alone cannot say which of the two the
residual follows. Holding the rung, and so $\fgap$, fixed and scaling only the
outlet resistances, from the test's own up to sixteenfold, settles it: the residual
times PE stays constant to under $1\%$ at every rung of the sweep, while across the
sweep, where $\fgap$ moves $1.73$
decades, it moves only from $0.10$ to $0.12$. The residual is a
finite-$\mathrm{PE}$ effect, not a coupling, and a graded negative control,
coupling the outlet to $\fgap$ at the $2\%$ level, is caught. It exists because
the test carries a run-out mismatch: its run-outs impose $F=0.5$ on
a ladder that would, left to itself, send $2.9\%$ of the flow down its rails, so
the outlet row holds a pressure imbalance, and $F(1-F)$ times that imbalance is
the residual to first order; with run-outs matched to the ladder's own split it
falls to the solver's floor. The PE here is the test's own axis, a rung-to-outlet
conductance ratio; the $48\times$ margin of C4's CQ2 is the preprint's, on a
different axis (the pressure-equilibration half-life against the device's
length), so the two numbers are not comparable, and the half-life does not bound
this residual. The ladder's half-life depends on its channel and rung
conductances alone: $k_{1/2}=0.366$ rows at the weakest rung, $109$ in the
test's $40$ rows, which clears the preprint's criterion. Yet the outlet row's
imbalance does not fall as rows are added, although the preprint's exponential
model would have it fall: that model covers an imbalance carried in from
upstream, and this one the outlet holds.

One caveat travels with this premise and the glossary states it explicitly: that
$1.73$-decade span is a \emph{deliberately widened numerical sweep} built to
demonstrate decoupling over a decade, not a range any built device occupies.
Gifford's calibration array caps at $\fgap = 5.76\times10^{-4}$ and no published
lineage device exceeds $1.0\times10^{-3}$. At the sweep's top,
$\fgap = 1.93\times10^{-3}$, the flux-rescaled gap Reynolds number would be
$\approx0.010$, at the $\mathrm{Re}_{\mathrm{stokes}}$ gate, which is
where the creeping-flow premise stops having room. The decoupling result is
therefore comfortable across every band a device occupies and marginal only at the
top of the test sweep, which is the honest way to report it.
\prov{code: tests/test\_C4\_decoupling\_v2.py; glossary \xref{app:units_equations}{A.9}}
\end{description}

% ----------------------- CRITICAL QUESTIONS -----------------
\field{Critical questions.} Three, all \emph{open by design}, and numbered CQ5--CQ7
rather than CQ1--CQ3 because these are the numbers the argument graph gives them
(\texttt{c8\_def\_pinned}, \texttt{c8\_def\_runout}, \texttt{c8\_def\_knee}), and
renumbering them here would break the correspondence this appendix exists to
maintain. Their ground atoms sit in the graph's withheld block, which records
them as possible without granting them: each computes \emph{out} until its
antecedent is asserted, so each defeater lies dormant. Closing them by fiat would
assert that a CIF device is
always above the knee, always run-out-matched and never length-budgeted, and the
first of those is false of a published device.

\begin{description}
\item[CQ5: Is the row count free? \textup{(\texttt{pinned\_row\_count})}]
Not always. The inversion consumes $N$ as the free variable that meets the yield
target. If a device-length budget pins $N$ externally, two targets meet one free
variable and the system is over-constrained. Nothing in the physics prevents
this; it is a packaging constraint, and the lineage has met it. In Xia's words the
recursion ``\emph{enables patterning devices of any total length that is
needed}'', and the next sentence reads ``\emph{In practice, however, one must
partition a lengthy CIF device into `legs' in order to fit its overall footprint
onto a standard substrate mould}''
\cite[p.~2 (Sci Rep PDF pagination)]{b2}. It is the cheapest of the three to check
before a design is committed: the specification checks it as G9, each leg's
length and the array's total against the leg budget. The design code does not
enforce it. Its G9 compares the array's total with the budget and reports the
result without failing, and only for the validator battery's own short test
device; \texttt{design} itself refuses nothing on length, and a design's report
prints the length beside the budget.
$\rightarrow$ \emph{open, withheld: an envelope boundary, not a defect.}
\item[CQ6: Are the run-outs matched? \textup{(\texttt{runout\_mismatch})}]
They must be built to match the schedule $\rho(N)$, and the tolerance is tight: a
$3\times$ mismatch collapses the filtrate fraction from $\runoutFmatched$ to
$\runoutFmismatched$, a factor of two in delivered yield (reading the mismatch at
arm level, as the specification defines it; trebling only one of the two side
run-outs gives $0.418$). The specification calls
this the most sensitive input its previous version left un-gated, and now gates it
at design time: G13 is a pre-flight assertion, before the march, that the run-out
resistances match the outlet ratio. The design code has no such pre-flight check,
since its G13 runs on the validator battery's own short test device. It meets the
condition by construction: it builds both run-outs at the last row's widths over
one length, so their ratio is that row's resistance ratio, and the march's closure
check refuses any row whose ratio misses its target, the outlet ratio at the last
row. What no gate sees is the device as built. $\rightarrow$ \emph{open,
withheld: gated by the specification at design time (G13) and met by construction
in the design code; a fabrication tolerance after it.}
\item[CQ7: Is every row on the plateau? \textup{(\texttt{below\_knee})}]
This is the one with a worked counterexample, and C7 has already done the work.
If a row runs below the knee, $\kapp$ becomes a function of the cutoff and the
operating law turns \emph{implicit}, $\dcrit=\kapp(\dcrit/\wc)\,\delta$; the
first triangular row is then the fixed point of
Figure~\ref{fig:c7-fixedpoint} rather than a direct solve. The published lineage
taper does exactly this at its wide end: Dinh's bracket device, a $75\um$ inlet
read at $\dcrit=\dinhBandHi\um$, gives $u=\dwc=0.029$ against the plateau floor
$u_{\mathrm{lo}}$, set here to the plateau-bin edge $\plateauUwc$, and the plateau
guard refuses it. Dinh's other two designs give $0.030$ and $0.032$ on the same
reading; those fall in the gap in Gifford's calibration data between $u=0.0294$
and $0.0353$, where the bin edge could sit anywhere, but their own stable roots,
which C7 reports as cutoffs of $1.02$--$1.47\um$ and $1.26$--$1.81\um$ across the
taper, are $u\approx0.020$ and $0.024$, below it wherever it sits. That
$\dinhBandHi\um$ is the \emph{upper} edge of C7's band, and the device's own
self-consistent cutoff at that end is lower still: at its initial gap fraction only
the sub-steric foot is self-consistent ($0.59\um$ at $75\um$), and under the
schedule Dinh describes, read in the corrected Xia form, the first root above the
foot, at $1.09\um$, appears only once the channel has narrowed to $74.98\um$ (C7).
The margin below the floor is therefore if anything understated here, whether the
gap fraction is held at its initial value or follows the schedule. C7 is thus not
merely a
corroborating band but a \emph{counterexample to unconditional application} of C8.
Nor is the $75\um$ lineage alone in this. On the taper Xia prints, RS-7 runs
$\dwc=0.023$ at its $300\um$ wide end (C5~CQ4), further below the floor than any of
Dinh's three designs, and the guard refuses it on the same test.
One caveat rides along, the one the Section~2 row states: the \emph{need} for the
guard is derived, from C7, but its \emph{value} is not. $u_{\mathrm{lo}}$ is the
plateau-bin edge, an empirical boundary nobody derived; what it is \emph{for} is
checked (above it the calibration points collapse onto one $\kapp$ about $22$ times
more tightly than below it), but the number itself is an external input, so the
guard is only as sharp as that number. $\rightarrow$ \emph{open, withheld: bounded by
a code-enforced guard, and shown on published designs.}
\end{description}

% ----------------------- RESIDUAL / STATUS ------------------
\field{Residual risk.} The three envelope boundaries above, none excluded and each
named: \texttt{below\_knee}, \texttt{runout\_mismatch}, \texttt{pinned\_row\_count}.
The \emph{inversion} carries no residual inside the envelope it is stated for: C8a
is proven, and it survives every
rival in the graph including the one still open, the
size-dependent $\feff$ of C3c.

\field{Status.} \textbf{Partially closed.} The map is proven \emph{locally}
invertible; its
realization is warranted \emph{inside the declared envelope}, with three named
boundaries dormant rather than excluded. The headline claim inherits this
conditionality, and states it.

% ===========================================================================
\clearpage
\section*{The thesis, and the fence around it}

\field{Summary.} Everything above is argued in service of one conditional
sentence: \emph{within a declared physical regime (rigid spheres in Stokes flow),
and provided a design stays inside its admissible operating envelope, a CIF device
can be designed directly from its two performance targets, a size cutoff and a
yield, because the laws governing those two targets are \emph{provably
decoupled}.} One
word in that sentence has to be pinned before anything else: throughout this
appendix \emph{yield} means the filtrate fraction $F$, the fraction of \emph{fluid}
sent to the filtrate channels, and never particle recovery. The two are different
quantities and the separation lemma is about the first of them. The thesis is an
\emph{accrual}: it rests on five premises at once, each established in its own
section. The conditions are not fine print. They are load-bearing (the
realizability leg is formally partially closed on exactly these envelope
conditions), and the scope is carried as a claim in the argument rather than as a
sentence only in the abstract, because an examiner reads an unstated exclusion as
an oversight and a stated one as a decision.

% ----------------------- ARGUMENT GRAPH ---------------------
\begin{figure}[t]
\centering
\begin{tikzpicture}[
  font=\scriptsize,
  leg/.style={draw, rounded corners, align=center, inner sep=3pt,
              text width=2.7cm, minimum height=17mm},
  ded/.style={draw, double, double distance=0.7pt, rounded corners,
              align=center, inner sep=3pt, text width=2.7cm, minimum height=17mm},
  part/.style={draw, dashed, rounded corners, align=center, inner sep=3pt,
               text width=2.7cm, minimum height=17mm},
  th/.style={draw, rounded corners, align=center, inner sep=4pt,
             text width=11.4cm, minimum height=13mm},
  up/.style={-{Stealth[length=2mm]}, thick},
  nt/.style={font=\scriptsize\itshape, text=black!60, align=center},
]
\node[leg]  (l1) at (-6.3,0)
  {\textbf{Operating law}\\ $\dcrit=\kapp\,\delta$\\ (C2, C3a--C3c)\\
   \textsf{[warranted]}};
\node[ded]  (l2) at (-3.15,0)
  {\textbf{Separation lemma}\\ $F\indep\fgap$\\ (C4)\\ \textsf{[proven]}};
\node[ded]  (l3) at (0,0)
  {\textbf{Locally invertible map}\\ triangular Jacobian\\ (C8a)\\ \textsf{[proven]}};
\node[part] (l4) at (3.15,0)
  {\textbf{Realizable design}\\ forward substitution\\ (C8b)\\
   \textsf{[partially closed]}};
\node[ded]  (l5) at (6.3,0)
  {\textbf{Declared scope}\\ rigid spheres,\\ Stokes flow\\ \textsf{[proven]}};
\node[th] (th) at (0,-3.2)
  {\textbf{THESIS.}\ CIF is engineerable from its performance targets, not from
   a tuned per-gap schedule, because cutoff and yield obey provably decoupled
   laws, \emph{inside a declared scope and a checked envelope}.\ \
   \textsf{[warranted]}};
\draw[up] (l1) -- (th);
\draw[up] (l2) -- (th);
\draw[up] (l3) -- (th);
\draw[up] (l4) -- (th);
\draw[up] (l5) -- (th);
\node[nt, anchor=north, text width=12.0cm] at (0,-4.35)
  {An accrual adds no inference beyond the conjunction, so the thesis carries no
   residual of its own; what it inherits from Leg~4 it inherits as a
   \emph{condition}, which is why the node is warranted and its statement is
   conditional. Border style follows the overview map: doubled $=$ proven,
   single $=$ warranted, dashed $=$ partially closed.};
\end{tikzpicture}
\caption{The thesis and its five legs. Every arrow is a premise of a single
accrual node: remove one leg and the thesis loses a premise, which is what
distinguishes these five from the corroborating branches (C5, C7, and C6's Davis
cross-check) that have no support path to the thesis at all; nor, as C9's CQ2
records, does $\ksteric=2$. The fifth leg is the scope declaration, drawn beside
the other four rather than beneath them because it is a premise in the same
sense: an argument that did not declare its regime would not be a weaker argument
for this thesis, it would be an argument for a different and false one. Note that
the thesis node is drawn with the single border of a \emph{warranted} claim even
though Leg~4 is dashed: an accrual inherits its premises' conditions as the
conditions of its own statement, and this statement says so in its own words.}
\label{fig:thesis-argmap}
\end{figure}

% ----------------------- CLAIM ------------------------------
\field{Claim.} CIF can be engineered from its performance targets instead of from
a tuned per-gap schedule: its yield follows a closed-form law with no fitted
constant and its selectivity a derived law with a single calibrated prefactor,
and the two are provably decoupled, so the size threshold and the filtrate
fraction are independent design \emph{inputs} rather than coupled outcomes. All
of it holds inside a validity envelope that the design tool checks: it refuses a
design that leaves the plateau or the Reynolds gates, and it reports, without
refusing, a design over its length budget (C8, CQ5).

% ----------------------- SCHEME -----------------------------
\field{Scheme.} Accrual. This is a local tag rather than a WRM~2008 scheme: the
node's premises are other \emph{nodes}, not atoms, and it carries no inference of
its own beyond the conjunction. Its weight is experimental, the weakest
evidential tier among its legs, and it is evaluated under preponderance of the
evidence.

% ----------------------- PREMISES ---------------------------
\field{Premises.} Five, each a claim argued in its own section.
\begin{description}
\item[Leg 1: the operating law (C2, C3a--C3c).]
$\dcrit=\kapp(\dwc,\,\text{depth})\,\wc\sqrt{\fgap/3}$, with $\kapp$ calibrated
rather than universal. Itself an accrual of exactly four premises: the skim law
(C2), proportionality (C3a), the plateau value (C3b), and the \emph{existence} of
the knee (C3c). C1 and C9 are deliberately \emph{not} listed: they establish and
then derive $\ksteric=2$, which underpins the mechanism and C3d's decomposition
but appears in none of those four premises and reaches this node by no path. C3d
is likewise absent: it interprets the prefactor's magnitude rather than
supplying it.
\item[Leg 2: the separation lemma (C4).]
$F=2\Rc/(\Rs+2\Rc)$ depends only on outlet geometry, so yield carries no cutoff
information and the two targets are genuinely two knobs.
\item[Leg 3: the locally invertible design map (C8a).]
The target Jacobian is lower-triangular with non-zero diagonal, so the inverse
exists locally throughout the envelope.
\item[Leg 4: realizability (C8b).]
Forward substitution realizes both targets on a device built to the resulting
schedule, inside the envelope, with three named boundaries dormant. This appendix
reports no device fabricated to such a schedule: the leg rests on the inversion and
three granted assumptions (C8).
\item[Leg 5: declared scope.]
Rigid spheres in Stokes flow at nearly constant viscosity: the regime in which
the skim cubic, the current-divider identity and the steric criterion are all
derived. Deformability, cell crowding and whole blood are \emph{declared} out of
scope, not overlooked.
\end{description}

% ----------------------- SCOPE ------------------------------
\field{Scope, stated as a claim rather than a disclaimer.} Three limitations are
carried inside the argument, each sourced.

\begin{description}
\item[Platelets are discoid, not spherical.]
The motivating application is also the limitation. Dinh describe platelets as
``\emph{$1\um$-thick oblate spheroids (discoids) ranging from 2 to
$7\um$ in diameter}'' \cite[p.~913]{b5} and name ``\emph{the characteristic
tumbling motion oblate spheroids undergo when flowing near the channel walls}''
as a separation challenge in its own right \cite[p.~921]{b5}. Gifford concede
the same (``\emph{approximately discoid}'') and, tellingly, determined the
platelet $\fgap$ by \emph{observation} rather than from the cutoff formula
\cite[p.~4502]{b1}. The diameter entering the cutoff law is therefore
\emph{orientation-dependent}, the rigid-sphere idealisation is an approximation
the founding paper already worked around empirically, and platelet-specific
validity is \textbf{not} established by this argument. That is precisely why the
thesis is scoped to rigid spheres.
\item[The Stokes regime is gated in code, not merely assumed, and the channel gate can bind.]
The numbers, not just the thresholds. Taking the whole flow through the narrow
outlet section, a worst case that bounds what the code computes from the local
flow, the \emph{channel-scale} Reynolds number of a water-like fluid is
${\approx}\ReWorstChannel$ at $30\,\mu\mathrm{L\,min^{-1}}$, rising to
${\approx}\ReHighFlow$ at $90\,\mu\mathrm{L\,min^{-1}}$; platelet-rich plasma, more
viscous than water, runs lower in proportion. The
Reynolds number \emph{at the gap} is obtained by rescaling the \emph{flux}, not
the length: a gap carries only the diverted fraction $\fgap$ of the channel flow,
so $\mathrm{Re}_{\mathrm{gap}}\approx\fgap\,\mathrm{Re}_{\mathrm{channel}}
\approx10^{-3}\times\ReWorstChannel\approx5.2\times10^{-3}\ll1$, evaluated at the
\emph{top} of the published lineage range rather than at a comfortable margin
below it. The gates are $\mathrm{Re}_{\mathrm{stokes}}=10^{-2}$ on the gap and
lamina and $\mathrm{Re}_{\max}=10$ on the channel. The design code computes both
from the geometry and the local flow and refuses a design that violates either
(its M14); the physics auditor separately checks the drawn device's inlet
against its own, stricter channel limit.

Two things are worth saying plainly rather than leaving to be found. First,
\textbf{the channel gate can bind}: for a water-like fluid,
${\approx}\ReHighFlow$ at $90\,\mu\mathrm{L\,min^{-1}}$ exceeds
$\mathrm{Re}_{\max}=10$, and the design code refuses that operating point even
when it uses the local flow, so the envelope is a real constraint on throughput
rather than a formality. For platelet-rich plasma the answer turns on a viscosity
no paper here states: the same point passes above about $1.4\,\mathrm{mPa\,s}$ and
is refused below it, as at the $1.2\,\mathrm{mPa\,s}$ the preprint's code takes for
plasma. Dinh ran V1 with platelet-rich plasma at exactly this rate (P-C4.2), so
whether the design code would refuse the lineage's own operating point turns on
that same unstated number; if it does, it refuses on the channel gate alone, since
the gap and lamina gates, on which the selection physics rests, stay at about
$3\times10^{-3}$ or below there, under their $10^{-2}$.
Second, the gap figure depends on which estimate is used: the flux argument above
is the one the scope gate evaluates, and a \emph{length}-rescaling estimate at the
same velocity gives different numbers (${\approx}1.3$ using the $19\um$ gap
width, ${\approx}0.02$ using the outlet lamina width), which must not be
conflated with it. The Stokes premise on which the whole skim derivation rests is
therefore checked, and checked by a stated method.
\prov{code: CIF\_v2.py \texttt{reynolds}, \texttt{design} (M14); physics auditor
(P8); reference implementation, VISCOSITY\_PLASMA; glossary \xref{app:units_equations}{A.9}}
\item[Two software contributions are out of \emph{this argument's} scope, deliberately.]
The targets-to-photomask tool (Prospectus~C7) and the machine-enforced validity
envelope (Prospectus~C8) carry no argument nodes. They are delivered as verified
software, warranted by their acceptance suite and the physics auditor rather than
by a defeasible inference. The silence is a decision, recorded here so that it
reads as one. \emph{Naming caution:} those are the \emph{prospectus} numbers,
not this appendix's; the collision is declared under
``C7 and C8 are the \emph{graph's} numbers'' (\S\ref{sec:b-notation}).
\end{description}

% ----------------------- CRITICAL QUESTIONS -----------------
\field{Critical questions.}
\begin{description}
\item[CQ1: Is the conditionality doing real work, or is it a hedge?]
Real work. The realizability leg is \emph{partially closed} on three named
envelope boundaries, and one of them has a published counterexample: the lineage
taper, which runs below the plateau floor at its wide end (C8~CQ7). An
unconditional version of this thesis would be false of a device that already
exists. $\rightarrow$ \emph{answered: the conditions are load-bearing, and one is
demonstrably binding.}
\item[CQ2: Does the one live open question threaten the thesis?]
No, and this is checked rather than argued. The single undefeated mechanistic
rival, a size-dependent $\feff(d)$ as the cause of the knee, moves exactly four
labels when granted: the rival and its mechanism come \emph{in}, the
finite-particle reading and the specific residual decomposition go \emph{out}. The
thesis, the operating law, the exponents and the \emph{existence} of the knee are
all untouched, because a size-dependent $\feff$ is itself a form of
size-dependence and so cannot defeat the variation it would merely reattribute.
The four-label blast radius is asserted and re-checked on every build.
$\rightarrow$ \emph{answered: the engineering thesis survives whichever way the
mechanism resolves.}
\item[CQ3: Does the thesis depend on the appendix's chosen weights?]
Barely, and the dependence is measured rather than assumed. Collapsing all six
evidential tiers to any common value above the proof-standard threshold moves
\emph{no} label; reversing the ranking entirely (expert testimony above
deductive proof) moves \emph{one} of $151$. The labelling therefore depends on
the tier weights through a single binary distinction, whether an argument
supporting a beyond-reasonable-doubt or clear-and-convincing conclusion clears
the threshold, and not through the ordering at all. $\rightarrow$
\emph{answered: the ordinal ranking is inert; only the threshold does work.}
\end{description}

% ----------------------- RESIDUAL / STATUS ------------------
\field{Residual risk.} Inherited, not new: the three envelope boundaries of
Leg~4, and the open evidentiary questions of Legs~1 and~2. The thesis adds no
residual of its own: an accrual carries no inference beyond the conjunction of
its premises.

\field{Status.} \textbf{Warranted, conditionally}: every leg is established
within its own section; the weakest of them, realizability, is partially closed
on three named and dormant envelope boundaries; and the headline claim inherits
that conditionality in the same breath as it is stated.


\clearpage
% >>> BEGIN GENERATED BIBLIOGRAPHY -- gen_B_bibliography.py
% Carried from glossary_official.tex. DO NOT EDIT: regenerate with
% gen_B_bibliography.py; validate.sh gate 15j diffs this region against a
% fresh generation and fails on any hand edit.
%
% THE KEY IS THE INVARIANT, NOT THE NUMBER. b5 is Dinh here exactly as it is
% in A.1-A.12. The numbers printed in a STANDALONE build are local to this
% file, because B cites a subset of the glossary's works; on integration this
% whole block is deleted and every \cite resolves against the manuscript's
% bibliography. See the INTEGRATION NOTE at the head of this file.
\begin{thebibliography}{99}

\bibitem{b1}
S.~C. Gifford, A.~M. Spillane, S.~M. Vignes, and S.~S. Shevkoplyas,
``Controlled incremental filtration: A simplified approach to design
and fabrication of high-throughput microfluidic devices for selective
enrichment of particles,'' \emph{Lab Chip}, vol.~14, no.~23,
pp.~4496--4505, 2014, doi: 10.1039/C4LC00785A.

\bibitem{b2}
H.~Xia, B.~C. Strachan, S.~C. Gifford, and S.~S. Shevkoplyas,
``A high-throughput microfluidic approach for 1000-fold leukocyte
reduction of platelet-rich plasma,'' \emph{Sci. Rep.}, vol.~6,
p.~35943, 2016, doi: 10.1038/srep35943.

\bibitem{b3}
H.~Xia, B.~C. Strachan, S.~C. Gifford, and S.~S. Shevkoplyas,
``Author Correction: A high-throughput microfluidic approach for
1000-fold leukocyte reduction of platelet-rich plasma,''
\emph{Sci. Rep.}, vol.~12, p.~14689, 2022, doi:
10.1038/s41598-022-18449-5.

\bibitem{b4}
D.~L. Lezzar \emph{et~al.}, ``A high-throughput microfluidic device
based on controlled incremental filtration to enable
centrifugation-free, low extracorporeal volume leukapheresis,''
\emph{Sci. Rep.}, vol.~12, p.~13798, 2022, doi:
10.1038/s41598-022-16748-5.

\bibitem{b5}
M.~T.~P. Dinh, A.~Mukhamedshin, K.~Abhishek, F.~W. Lam, S.~C. Gifford,
and S.~S. Shevkoplyas, ``Separation of platelets by size in a
microfluidic device based on controlled incremental filtration,''
\emph{Lab Chip}, vol.~24, no.~4, pp.~913--923, 2024, doi:
10.1039/D3LC00842H.

\bibitem{b6}
S.~Yang, A.~\"{U}ndar, and J.~D. Zahn, ``A microfluidic device for
continuous, real time blood plasma separation,'' \emph{Lab Chip},
vol.~6, no.~7, pp.~871--880, 2006, doi: 10.1039/B516401J.

\bibitem{b7}
COMSOL Multiphysics\textregistered\ v.~6.3, COMSOL AB, Stockholm,
Sweden, \url{www.comsol.com}.

\bibitem{b12a}
A.~H. Routt, ``CIF verified derivation record'' [software], release tag
and content hash recorded in the dissertation reproducibility archive,
2026. \emph{Numerical authority: every constant is recomputed on each run
and pinned by content hash, and the whole record is regenerated and
re-verified on every commit.}

\bibitem{b13}
K.~W. Oh, K.~Lee, B.~Ahn, and E.~P. Furlani, ``Design of pressure-driven
microfluidic networks using electric circuit analogy,'' \emph{Lab Chip},
vol.~12, no.~3, pp.~515--545, 2012, doi: 10.1039/C2LC20799K.

\bibitem{b14}
Z.~Li, C.~Liu, and J.~Sun, ``Hydraulic--electric analogy for design and
operation of microfluidic systems,'' \emph{Lab Chip}, vol.~23,
pp.~3311--3327, 2023, doi: 10.1039/D3LC00265A.

\bibitem{b15}
M.~Yamada and M.~Seki, ``Hydrodynamic filtration for on-chip
particle concentration and classification utilizing microfluidics,''
\emph{Lab Chip}, vol.~5, no.~11, pp.~1233--1239, 2005, doi:
10.1039/B509386D.

\bibitem{b17}
D.~W. Inglis, J.~A. Davis, R.~H. Austin, and J.~C. Sturm, ``Critical
particle size for fractionation by deterministic lateral displacement,''
\emph{Lab Chip}, vol.~6, no.~5, pp.~655--658, 2006, doi:
10.1039/B515371A.

\bibitem{b18}
D.~W. Inglis and N.~Herman, ``A scalable approach for high throughput
branch flow filtration,'' \emph{Lab Chip}, vol.~13, no.~9,
pp.~1724--1731, 2013, doi: 10.1039/C3LC50192B.

\bibitem{b21}
V.~Doyeux, T.~Podgorski, S.~Peponas, M.~Ismail, and G.~Coupier,
``Spheres in the vicinity of a bifurcation: elucidating the
Zweifach--Fung effect,'' \emph{J. Fluid Mech.}, vol.~674,
pp.~359--388, 2011, doi: 10.1017/S0022112010006567.

% REWRITTEN from the glossary's entry: The glossary's entry points at the
%   glossary's own three-family table. Appendix B has no such table, so the
%   reference would dangle; the recomputation is named by the appendix that
%   holds it instead.
\bibitem{b22}
J.~A. Davis, \emph{Microfluidic Separation of Blood Components through
Deterministic Lateral Displacement}, Ph.D. dissertation, Princeton
University, 2008. (Eqn~2-22; recomputed in the glossary for the cross-family
comparison.)

% --- Appendix B only: not in the glossary, with the reason for each ---
% Appendix B's scheme authority. Every node in this appendix is an instance
%   of one of its schemes; the glossary is a units-and-equations record and
%   cites no argumentation theory at all.
\bibitem{bwrm}
D.~Walton, C.~Reed, and F.~Macagno, \emph{Argumentation Schemes}. Cambridge, U.K.: Cambridge Univ.\ Press, 2008.

% Named in C2 for the depth-integrated rectangular-duct velocity field (his
%   eqn~(15)), and identified there as Inglis \& Herman's ref.~17, which is
%   what makes the duct-field provenance shared rather than merely parallel.
%   The glossary reaches the duct law through Yang (b6), a different object
%   -- the flow-resistance solution, not the velocity field -- and so
%   carries no entry for Brody. No issue number or DOI: the paper prints
%   neither, and an unverified identifier is worse than none.
\bibitem{bbrody}
J.~P. Brody, P.~Yager, R.~E. Goldstein, and R.~H. Austin, ``Biotechnology at low Reynolds numbers,'' \emph{Biophys.\ J.}, vol.~71, pp.~3430--3441, 1996.

% Cited in C1 (CQ4) for the centre-position picture in pinched-flow
%   fractionation, which it writes as y_P = w_P - D/2 and tests against the
%   paper it comments on. The glossary does not cite it. Byline, volume,
%   pages and DOI read off the in-project PDF (one author, not "et al." as
%   ledger DG-6 had it).
\bibitem{bmort}
N.~A. Mortensen, ``Comment on `Pinched flow fractionation: Continuous size separation of particles utilizing a laminar flow profile in a pinched microchannel','' \emph{Anal.\ Chem.}, vol.~79, no.~23, pp.~9240--9241, 2007, doi: 10.1021/ac7018816.

% The companion preprint. Appendix B cites it for the COMSOL validation of
%   $F$'s magnitude; the glossary reaches the same numbers through the
%   derivation record (b12a) and so has no entry for the preprint.
\bibitem{bpre}
A.~H. Routt, ``A closed-form outlet identity and open-source reference implementation for controlled incremental filtration,'' engrXiv, 2026, doi: 10.31224/7152.

\end{thebibliography}
% <<< END GENERATED BIBLIOGRAPHY

\end{document}